\documentclass[11pt]{article}
\usepackage[T1]{fontenc}
\usepackage[utf8]{inputenc}
\usepackage{lmodern}
\usepackage[margin=1in]{geometry}
\usepackage{amsmath,amssymb,amsthm,mathtools,bm}
\usepackage{booktabs,array,longtable}
\usepackage{enumitem}
\usepackage{hyperref}
\usepackage{xurl}
\usepackage{seqsplit}
\newcommand{\OPHPaperReleaseID}{r1520}
\newcommand{\OPHPaperReleaseDate}{July 9, 2026}
\newcommand{\OPHPaperReleaseBanner}{%
\begin{center}
\small\textbf{Paper release:} \texttt{\OPHPaperReleaseID}\quad\textbf{Released:} \OPHPaperReleaseDate
\end{center}
}


\hypersetup{colorlinks=true,linkcolor=blue,citecolor=blue,urlcolor=blue}
\setlength{\emergencystretch}{3em}
\setlength{\LTleft}{0pt}
\setlength{\LTright}{0pt}
\setlength{\tabcolsep}{5pt}
\renewcommand{\arraystretch}{1.08}

\newtheorem{theorem}{Theorem}[section]
\newtheorem{corollary}[theorem]{Corollary}
\newtheorem{proposition}[theorem]{Proposition}
\newtheorem{lemma}[theorem]{Lemma}
\newtheorem{definition}[theorem]{Definition}
\newtheorem{axiom}[theorem]{Axiom}
\newtheorem{assumption}[theorem]{Assumption}
\newtheorem{remark}[theorem]{Remark}

\newcommand{\SU}{\mathrm{SU}}
\newcommand{\U}{\mathrm{U}}
\newcommand{\ellbar}{\bar{\ell}}
\newcommand{\Conf}{\mathrm{Conf}}
\newcommand{\ophid}[1]{\path{#1}}
\newcommand{\decnum}[1]{\vcenter{\hbox{\parbox{0.68\linewidth}{\small\ttfamily\seqsplit{#1}}}}}

\title{Deriving the Particle Zoo from\\
Observer Consistency}
\author{B.\ M\"uller \and Alexander Osika \and Mario Poneder \and Kai Xue \and Maarten Antonie Visser \and David Matscheko}
\date{\OPHPaperReleaseDate}

\begin{document}
\maketitle
\OPHPaperReleaseBanner

\begin{abstract}
Observer Patch Holography (OPH) proposes to read one screen cell as an outside area and as
the electromagnetic observation scale seen from within. Matching the two readings defines a
declared numerical pixel map; its physical endpoint relation is unconstructed. This paper carries the Standard
Model structure of the companion compact paper into particle calculations:
charge quantization, conditional gauge and gravitational carrier modes, the
electroweak bridge, the hierarchy relation, and the Higgs/top quantitative
branch. The source-audit weak-boson law gives the running/chart coordinates
\((80.330,91.119)\)~GeV (closure row CL-5). They are not commensurate with the
PDG mass-dependent-width Breit--Wigner parameters or with complex-pole masses, and therefore
provide no physical near-hit evidence. The stated criticality branch gives
\((m_H,m_t)=(125.77,172.63)\)~GeV. Gauge structure is imported from the
companion compact paper as an input, and no nonzero particle mass is
emitted as a prediction; the per-sector ledger carries that status row by row.

The same framework provides sharp failure tests. A reciprocal quark-shape
candidate misses its held-out Yukawa coordinates by \(21.56\%\), and the
weighted-cycle neutrino candidate fails the stated oscillation-data test. For
charged leptons, a connected six-dimensional response register and its
tracial GNS representation derive the Koide balance \(Q=2/3\). Its attachment
to a physical chiral mass response is work in progress; the \(2/9\) phase and
individual mass ratios are comparison inputs. Absolute charged-lepton masses
and hadron masses are reported only on their declared empirical-closure
surfaces.
\end{abstract}

\section*{What This Paper Contributes}

The structural particle content is inherited from the compact SM/GR paper: receipt-conditional compact gauge
reconstruction, the realized Standard Model quotient, exact hypercharge, three colors, three
generations, and the connection/metric carrier roles. The compact paper gives classical massless
carrier modes only after explicit action/background/phase receipts and keeps quantum-particle
poles behind a stronger gate. This paper asks what that structure does after the
local declared-map coordinate is supplied.

The OPH addition is the forward particle readout. A single screen-cell coordinate \(P_\star\)
feeds the electromagnetic branch, the electroweak transport surface, and the mass rows that can be
certified on the declared branch. The hierarchy result is the cleanest example: the product
adjoint fixes \(m_{\rm rep}=2(8+3+1)=24\), and the selected exact branch emits the dimensionless hierarchy
\(v/E_\star\), with Higgs naturality defect \(\epsilon_H=0\).  A physical weak
scale in GeV requires an independently source-closed \(E_\star\). The whole chain runs at zero
continuous dials: each declared pixel map carries a machine-certified interval proof of a unique
fixed point. The falsification program covers only mature structural
non-cosmological rows. The same
paper keeps weaker entries visible without upgrading them. Charged
lepton absolute masses and first-principles hadron masses are outside the emitted theorem,
benchmark checks are named as checks, and measured endpoint data are isolated when a transport
bridge uses them.

\section{Introduction}

Observer-Patch Holography asks a concrete follow-up question after the gauge structure is fixed.
Can one local screen constant organize the particle spectrum, or are the masses and mixings a
collection of unrelated inputs?

The structural companion papers start from observer consistency on a finite screen. Physical data are organized in local patch algebras, nearby observers must agree on overlaps, realized states obey a local maximum-entropy and refinement principle, generalized entropy is recoverable, and the realized low-energy branch is fixed by Minimal Admissible Realization (MAR). On the controlled BW/geometric branch, that basis yields the Lorentzian and Einstein branches; on the realized compact-gauge branch it yields the Standard Model gauge structure
\[
\frac{\SU(3)\times\SU(2)\times\U(1)}{\mathbb Z_6},
\qquad
N_c=3,
\qquad
N_g=3.
\]
The Lorentzian branch, Einstein recovery, and the realized Standard Model gauge structure written
above are imported from \cite{ophcompact} with their stated branch conditions. That same paper also derives the hypercharge lattice and
the realized color triplet \(N_c=3\) and generation count \(N_g=3\). This paper uses those results as input and summarizes only the
downstream particle-spectrum continuation. If that derivation is correct, it should also constrain the particle spectrum.
The substrate-neutral normal-form framework of Ref.~\cite{observableNormalForms} isolates the
same-source confluence, cross-source boundary-identification, liveness, and refinement criteria
used upstream.  It supplies no particle-sector selector or numerical particle output; every such
claim below is conditional on the compact-paper branch and the particle-specific receipts.

The particle branch studied here is a one-fixed-point forward reconstruction problem with an
explicit closure matrix. Instead of treating masses, mixings, and Yukawa parameters as independent
inputs, it tries to propagate one dimensionless pixel ratio \(P\) through the spectrum. The
pipeline does not emit photon, gluon, or graviton particle masses from symmetry labels. It records
their conditional classical carrier modes and open quantum-particle gates, while the Higgs
candidate is on its declared quantitative surface. The hierarchy/naturality bridge is closed on its
selected branch. The restricted quark source-spread non-identifiability theorem, the common-scale
rejection of its reciprocal-ray candidate, and the weighted-cycle neutrino lane are
visible as obstruction or comparison surfaces. Target-anchored quark mass textures, the mixed-convention
formula diagnostics, and compare-only absolute neutrino attachments are withheld from public
prediction tables. The \(P\)-closure root,
the electroweak \(W/Z\) rows, charged leptons, and hadrons have the sector-specific boundary
statuses recorded below. This paper develops the
derivation chain from the axioms and candidate \(P\) trunk to the reported particle outputs, with
the closure boundary stated by sector.

\subsection{Why this particle content is selected}

On the declared branch, particles are not inserted by hand as a list of elementary
ingredients. They are the stable observer-visible excitations and carrier modes that
are left once three layers of structure have been fixed:
\begin{enumerate}[leftmargin=2em]
\item overlap consistency on the observer-patch network;
\item compact gauge reconstruction from the theorem-produced transportability criterion and fixed-cutoff bosonic sector category generated under one common stagewise strict representative, together with the compact paper's explicit refinement receipt for the refinement/fiber ladder and realized cofinal MAR-admissible witness data;
\item minimal admissible realization of the low-energy branch.
\end{enumerate}
The first layer says what kinds of local data can be compared consistently across neighboring
patches. The second assembles persistent zero-obstruction sector data into a compact gauge
structure through the compact paper's transportability, category, and refinement/fiber theorems.
The third selects which of the admissible low-energy branches is physically realized.
At that point one is not choosing an arbitrary particle catalog. One is reading off the
carrier and matter content of the realized branch.

\emph{Recovering Relativity and the Standard Model from Observer Overlap Consistency}~\cite{ophcompact} is therefore central to this draft. That paper does
the structural work that fixes the realized low-energy package
\[
\frac{\SU(3)\times\SU(2)\times\U(1)}{\mathbb Z_6},
\qquad
N_c=3,
\qquad
N_g=3,
\]
together with one admissible Higgs doublet and the realized chiral matter content. The
construction does not begin by guessing photons, gluons, quarks, leptons, and weak bosons and
then asking how to fit them. It first derives the realized gauge-and-matter branch, and that
branch determines which elementary carriers and matter families are available at low energy.

On that selected one-Higgs branch the scalar carrier has a canonical local screen-chart model.
Let \(C_{\rm EW}\cong\mathbb{CP}^1\) be the support-visible electroweak chart and fix the positive
Hopf line-bundle convention by the neutral component condition \(Q(\phi^0)=0\). Borel--Weil gives
\[
H_{\rm OPH}=H^0(C_{\rm EW},\mathcal O(1))\cong\mathbb C^2,
\]
the first nontrivial holomorphic section space. With OPH's hypercharge and \(\mathbb Z_6\)
normalization this is exactly the \((1,2)_{1/2}\) Higgs carrier. Projectivization classifies a
nonzero section direction as a point of \(\mathbb P(H_{\rm OPH})\cong\mathbb{CP}^1\), but it
forgets the scalar hypercharge phase. For the lower-component vacuum vector
\[
\phi_0=\frac{v}{\sqrt2}\binom{0}{1},\qquad v\ne0,
\]
one has
\[
e^{i\alpha T_3}e^{i\beta Y}\phi_0
=e^{i(\beta-\alpha)/2}\phi_0.
\]
Consequently \([\phi_0]\) has the projective two-torus stabilizer
\(\U(1)_{T_3}\times\U(1)_Y\), modulo the inherited finite center, whereas the vector
\(\phi_0\) is fixed only when \(\beta=\alpha\), locally. Its connected stabilizer is therefore the
electromagnetic diagonal \(\U(1)_Q\), generated by \(Q=T_3+Y\). Thus the projective geometry
explains the carrier-ray classification, while the chosen nonzero vector supplies the
symmetry-breaking statement. This explains the representation, charge, and symmetry-breaking
geometry of the one-Higgs slot. It does not derive \(m_H\), the quartic, \(v\), or
Coleman--Weinberg dynamics; the
weak scale, Higgs/top quantitative surface, and \(\epsilon_H=0\) hierarchy/naturality closure are
the OPH quantitative branch.

This also explains the phrase ``and no others.'' The structural OPH claim concerns the realized
low-energy branch selected by the stated admissibility conditions. Once one imposes those
conditions, the realized low-energy branch is the
Standard-Model branch. Additional connected gauge generators,
additional light Higgs multiplets, extra light chiral families, or low-energy supersymmetric
partners do not appear on that realized branch because they would enlarge the admissible package
that MAR selects against. In that sense the observed light particle zoo is not an arbitrary menu;
the representation/content package is the minimal realized one. Propagating gauge-particle
existence requires the action, phase, and quantum-pole gates below.

The individual families then emerge for different reasons. The realized electromagnetic, color,
and dynamical-metric branches identify connection and metric carrier roles. Explicit Maxwell,
perturbative pure-Yang--Mills, and pure-Einstein quadratic actions then yield classical transverse
or TT massless modes on their stated backgrounds and phases. Quantum photon, gluon, or graviton
states require the separate physical-Hilbert-space, pole-residue, and asymptotic/phase receipt. The weak bosons arise when
the electroweak gauge sector is propagated through its quantitative closure branch. Quarks and leptons arise from the
realized chiral matter package together with the three-generation structure \(N_g=3\). Their
family splittings are then read from the deeper overlap-transport and excitation machinery, not
from a second arbitrary postulate that says ``copy the family three times and assign masses by
hand.'' Once color is realized and confined, stable hadrons are composite readout channels of the
quark/gluon sector.

The role of the pixel constant \(P\) is also important to state clearly. Here, \(P\) does not
decide whether an electromagnetic particle pole exists, whether there are three colors, or whether
quarks come in three generations. The electromagnetic representation role and the color/generation
facts belong to the gauge-and-matter branch fixed before the quantitative closure step. What \(P\)
does is set the shared quantitative scale on which receipt-certified content is read out
numerically. So the logic of the paper is:
\[
\begin{aligned}
\text{observer consistency}
&\to
\text{realized gauge/matter branch} \\
&\to
\text{representation roles}
\to
\text{conditional action-level carrier modes} \\
&\to
\text{common scale }P
\to
\text{receipt-gated quantitative spectrum}.
\end{aligned}
\]
That is how the derivation explains why a universe with the realized branch and the
shared pixel scale exhibits this particle zoo.

\subsection{Logical basis and reading rule}

The particle derivation uses the same OPH basis as the SM/GR derivation in Ref.~\cite{ophcompact}. Its five
axioms are the screen net, overlap consistency, local MaxEnt with refinement stability,
recoverable generalized entropy, and Minimal Admissible Realization (MAR). On that basis, the
SM/GR derivation in Ref.~\cite{ophcompact} supplies the structural chain used here: receipt-conditional compact gauge reconstruction as a classification step, MAR selection of the realized
Standard Model quotient, the hypercharge lattice, the realized color triplet \(N_c=3\), and the generation count \(N_g=3\).
Refs.~\cite{ophconsensus,ophmicrophysics} supply the patch-net, repair, measurement, and
regulated screen language for flavor transport and observer-facing readout.

The quantitative particle side adds one local closure variable. The common pixel ratio \(P\),
reported by the synthesis paper's incomplete outer/inner declared map~\cite{ophsynthesis}, feeds the forward
electroweak map and its descendants. The public empirical comparison branch displays
\[
\alpha^{-1}(0)=137.035999177(21),
\qquad
\alpha(0)\simeq0.00729735256433,
\qquad
P_C\simeq1.6309682094.
\]
The comparison pixel \(P_C\) is defined from the measured endpoint; it is not a source-root
coordinate.
The computation has a fixed source order: golden-ratio entropy balance gives
\(\varphi\), boundary Gaussian normalization supplies the \(\sqrt{\pi}\) width,
a trial \(P\) feeds the source map through unification, running, and electroweak
anchoring, and Ward-projected electromagnetic transport gives the Thomson
endpoint used by the outer/inner pixel fixed point.
The declared numerical map fixes a unique root on its certified interval. The map is
incomplete because no derivation identifies its numerical endpoint with the physical
electromagnetic readout. The
first-principles diagnostic trunk records the certified coordinates
\[
\begin{aligned}
P_{\mathrm{fwd}}&=1.630972095858897\ldots,\\
\alpha_{\mathrm{root}}^{-1}
&=136.994835177413\ldots.
\end{aligned}
\]
The root is the interval-certified unique fixed point of the declared numerical map
(CL-6, closed).
The approximately \(0.041\) inverse-alpha difference to the public Thomson endpoint is a residual
of that incomplete map. It neither derives a QCD/hadronic contribution nor identifies a physical
endpoint relation.
The QCD-free hierarchy witness is the cleaner first-principles stress test. Combining the source/root
value with the finite-screen unified gauge-width contribution evaluated at the CODATA-derived
comparison pixel, \(\alpha_U(P_C)\), gives the mixed-provenance no-hadron diagnostic
\(A_{\alpha_U}^{\mathrm{fp}}=137.0359595136\ldots\), below the Thomson endpoint. It is
mixed-provenance comparison bookkeeping rather than a source-only fine-structure prediction. The certified
self-consistent gauge-width fixed point is \(\alpha^{-1}=137.035660136946577\ldots\)
(CL-2); the mixed diagnostic combines the inner value at \(P_{\mathrm{fwd}}\) with
\(\alpha_U\) at the SL-3 pixel and is excluded from the physical output ledger.
The detailed endpoint table appears in Section~\ref{sec:particle-input-ledger}.
The separate proposed
cosmic record-closure target \(N_{\mathrm{CRC}}=F(N_{\mathrm{CRC}})\), alongside the separate
declared count-density selector
\(N_\star=\arg\max_N[\log|\Omega^{\mathrm{sc}}_N|-N]\), belongs to the cosmological-capacity branch and
supplies no particle theorem. The Newton coupling uses the separate
selected no-\(G\) scale certificate \(\gamma_\star=\ell_\star\nu_{\mathrm{Cs}}/c\), equivalently
\(B_\star=3\pi/\ell_\star^2\), so the particle branch does not derive \(G\) by back-solving the
pixel ratio. The physical map \(F\) has not been constructed. Its identification with the
electroweak bridge is conditional on CP-1 to CP-3; the bridge value
\(N_{\mathrm{EW}}\simeq3.532\times10^{122}\) is about \(6.6\) percent above the
Planck-\(\Lambda\) central capacity \(N_\Lambda\simeq3.313\times10^{122}\). If the missing
physical identifications are supplied, the pair \((P_\star,N_{\mathrm{CRC}})\) determines
dimensionless curvature products such as
\(\Lambda_\star a_{\mathrm{cell}}=3\pi P_\star/N_{\mathrm{CRC}}\), not the SI scale product
\(\Lambda_\star N_{\mathrm{CRC}}\) by itself. The rounded
\(N_\Lambda\simeq3.313\times10^{122}\) display is a cosmological central-value label, not the
high-precision input for \(G\). Structural carriers, quantitative outputs, exact sidecars, and
continuation lanes are therefore downstream branches of the declared theorem checklist, with
particle-facing quantitative burden carried by the local declared-map root \(P_\star\).
The public fine-structure display row is a comparison endpoint conditioned on the measured
Thomson value. The first-principles diagnostic row and empirical hadron closure
rows are separate support records.
Section~\ref{sec:particle-input-ledger}
records that boundary in full.

\section{P-Closure and the Reverse-Engineering Strategy}

Before the claim-tier table, the reverse-engineering claim should be stated plainly. In
ordinary phenomenological use, the Standard Model does not derive the observed particle masses and
mixings from one common microscopic constant. It is usually presented with \(19\) free parameters
in the minimal massless-neutrino theory, and with at least \(26\) once neutrino masses and mixing
are included, depending on neutrino-sector conventions. The OPH particle program aims to organize
that situation around one universal pixel ratio \(P\) from the outer/inner closure relation
described in the synthesis paper. Whether any parameter compression is realized follows
from counting the declared settings against the independent landed basins. The same \(P\) must drive all
downstream bosonic, quark, lepton, neutrino, and hadronic branches.

One clarification matters. The broader self-closure formulation proposes a global screen capacity
\[
N_{\mathrm{CRC}}=F(N_{\mathrm{CRC}}),
\]
where \(F(N)\) would be the active cosmic record capacity read back by observers inside the universe
supplied with capacity \(N\). This physical map has not been constructed. G2-GAP-1 identifies it
with the electroweak bridge only conditionally on CP-1 to CP-3. A separate declared count-density selector is
\[
N_\star=\arg\max_N\left[\log|\Omega^{\mathrm{sc}}_N|-N\right],
\]
whose interior-maximum and smoothness premises are declared inputs that the lexicographic
minimality axiom does not license. Its differentiable mathematical map is
\(T_\eta(N)=N+\eta\,d(\log|\Omega^{\mathrm{sc}}_N|-N)/dN\) under the stated strict-concavity
certificate. No proof identifies this mathematical map with \(F\). The conditional bridge value
\(N_{\mathrm{EW}}\simeq3.532\times10^{122}\) is about \(6.6\) percent above the
Planck-\(\Lambda\) central capacity \(N_\Lambda\simeq3.313\times10^{122}\), a
propagated \(2.4\) to \(2.5\) one-dimensional sigma; read through
\(\Lambda_\star\ell_\star^2=3\pi/N\), the same bridge is the zero-dial relation
\(\Lambda\ell_P^2=3e^{-6\pi/(P\alpha_U(P))}\) up to a menu-priced one-term
correction, with the same exponent that runs the weak hierarchy of this paper. The
particle-spectrum derivation studied here uses the local declared-map root \(P_\star\) in place of
a long particle-by-particle parameter list.

Observation is allowed to supply a branch hint. That is a consequence of the OPH picture, where
the universe is a closed fixed structure and internal observers can read approximate coordinates
from experiments. The proof obligation is stricter: the exact \(P_\star\) used by the particle
branch is computed by the self-referential numerical pixel equation, and the root is unique on the
declared interval. A further derivation is required before that root is a physical endpoint.
The project consistency stack places this reconstruction inside the forcing
chain C1--C10: the pixel equation is the C9 loop-closure row, and the uniqueness statement
for the closure map (a self-map and derivative bound certified by interval arithmetic on the
declared interval) is carried by the shipped interval contraction certificate. The
domain-global statement is discharged as well: the companion domain-global certificate bounds
\(\sup|g'|<1\) on every piece of a 256-piece subdivision of the declared numerical domain
(\(\alpha^{-1}\in[100,200]\), both readout maps, empty exceptional set), so each declared readout
map has exactly one fixed point on that numerical domain. This certificate does not provide the
missing physical endpoint relation.

\subsection{Why a single \(P\) matters}

The structural theorems determine \emph{what kind} of particle world is
realized, but not \emph{what absolute scale in GeV} that world occupies. The overlap,
modular, gauge, anomaly, and admissibility arguments recover the realized Standard Model branch,
the exact hypercharge pattern, three generations, and three colors. The carrier-mode theorem is an
additional action-level result and does not emit zero-GeV quantum particles. These structural
results do not by themselves tell us the numerical values
of the \(W\) boson mass, the Higgs boson mass, or the up-quark mass.

The role of the pixel closure is to supply one common quantitative scale
variable for the entire downstream spectrum. In the synthesis paper, the
fine-structure lane asks for the nonzero detuning of a holographic screen cell
such that the cell's outer geometric displacement from perfect self-similar
equilibrium equals the electromagnetic observation scale emitted by the
universe living on that same screen.
The outer side of the closure is
\[
P=\varphi+\alpha_{\mathrm{in}}(P)\sqrt{\pi}.
\]
The first-principles computation is a five-step source chain. First, the
golden-ratio entropy balance of the local screen cell supplies
\(\varphi=(1+\sqrt5)/2\). Second, boundary maximum-entropy normalization fixes
the width of the leading electromagnetic detuning to \(\sqrt{\pi}\). Third, a
trial \(P\) is sent through the source map
\[
M_U(P)=E_P e^{-2\pi}P^{1/6},\qquad E_{\mathrm{cell}}(P)=\frac{E_P}{\sqrt P},
\]
followed by the heat-kernel closure
\[
\bar\ell_{\mathrm{SU}(2)}(t_2(P))+\bar\ell_{\mathrm{SU}(3)}(t_3(P))=\frac{P}{4},
\]
which selects \(\alpha_U(P)\) and the running family
\(\alpha_i(m_Z;P)\). Fourth, electroweak mixing gives the source anchor
\[
a_0(P)=\alpha_2^{-1}(m_Z;P)+\frac53\alpha_1^{-1}(m_Z;P).
\]
Fifth, Ward-projected \(\U(1)_Q\) transport gives the Thomson endpoint
\[
A_T(P)=T_Q(a_0(P),F_{\mathrm{src}}(P))=\alpha_{\mathrm{em}}^{-1}(0;P),
\]
and the cell closes when
\[
P=\varphi+\frac{\sqrt{\pi}}{A_T(P)}.
\]
The numerical solve is therefore a root problem for
\[
H(P):=P-\varphi-\frac{\sqrt{\pi}}{A_T(P)}.
\]
The branch interval is localized from the observer-facing data, the solver evaluates \(A_T(P)\)
from the declared source map at each candidate point, and the accepted value is the unique zero of
\(H\) on that interval. The measured endpoint can locate the interval. It cannot replace the root
condition.
This proves uniqueness only for the declared numerical map. The derivation connecting its root to
the observer-supporting physical electromagnetic endpoint is missing.
The CODATA-conditioned comparison readout
\(\alpha^{-1}(0)=137.035999177(21)\) gives
\[
P\simeq1.6309682094.
\]
The same equation is the local ruler for the downstream particle rows. The five-layer
OPH diagnostic trunk emits the certified source-side diagnostic point
\[
\begin{aligned}
P_{\mathrm{fwd}}&=1.630972095858897\ldots,\\
\alpha_{\mathrm{root}}^{-1}
&=136.994835177413\ldots.
\end{aligned}
\]
The root is an incomplete declared-map output (CL-6, closed), not a physical fine-structure
prediction.
The detailed table records how the diagnostic trunk, endpoint residual, and
source-spectral payload fit together.
A separate hardware note reports an optical-cavity check of the same declared geometry. It is a
non-discriminating engineering check and carries no physical endpoint promotion.
Once the diagnostic map coordinate is set, the particle question becomes whether all downstream readouts
can be written as
\[
X_j = G_j(P)
\]
with no new sector-specific constants inserted by hand. This paper studies exactly that downstream
map.

\subsection{How the electroweak branch works}

In the implementation used here, the quantitative electroweak branch is the first major readout from \(P\). Its algebraic
output family is
\[
\{M_W^{(10)},\ M_Z^{(10)},\ \alpha_{\mathrm{em}}^{-1}(q^2),\ \sin^2\theta_W(q^2),\ v\},
\]
evaluated from one declared input \(P\) on the printed running/matching/threshold/scheme package.  The superscript
\((10)\) denotes a D10 mass-chart coordinate, not a certified complex pole. The construction
is straightforward in spirit: start from \(P\), build the electroweak running family, select the physical carrier point on that family, and read off the \(W\) boson and \(Z\) boson pair from that selected point. The Higgs/top critical stage
inherits that same electroweak core and does not introduce a new
free-input sector. On the declared Ward-projected transport branch, the low-energy electromagnetic row is represented by the Thomson endpoint
\[
\alpha_{\mathrm{Th}}^{-1}(P)=\lim_{q^2\to 0}\alpha_{\mathrm{em}}^{-1}(q^2;P)
\]
of that same electromagnetic transport family.
The forward logical order on this branch is:
\[
P
\longmapsto
\bigl(M_U(P),E_{\mathrm{cell}}(P)\bigr)
\longmapsto
\alpha_U(P)
\longmapsto
\bigl(t_U(P),t_{\mathrm{tr}}(P)\bigr)
\longmapsto
\bigl(t_2(P),t_3(P),v(P)\bigr)
\longmapsto
\alpha_i(\mu_\ast;P).
\]
The forward transmutation certificate records that the same runtime basis reconstructs the unified diffusion parameter \(t_U(P)=4\pi^2\alpha_U(P)\) and transmutation exponent \(t_{\mathrm{tr}}(P)=2\pi/((N_c+1)\alpha_U(P))\) as the pixel-closure solve itself. The runtime subgraph reads measured electroweak data only on its validation surface, but no frozen provenance receipt excludes target dependence or selects one repair law.

The electroweak branch is the place where the named pixel coordinate organizes the weak-sector rows. The
reference-fitted inverse adapter displays
\[
M_W=80.3625~\mathrm{GeV},
\qquad
M_Z=91.1879~\mathrm{GeV}.
\]
The selected-carrier chart, the runtime-target-separated formula candidate, and
the reference-fitted coherent repair diagnostic are recorded separately in the support
table below.

\subsection{What the electroweak branch fixes}

On the declared electroweak running/matching/threshold/scheme surface, the pixel variable \(P\) supplies the
source basis
\((\alpha_U,\alpha_{2,m_Z},\alpha_{Y,m_Z},\eta_{\mathrm{source}},v)\) used by the
candidate chart. The displayed value law is an exact implication of the
five-part quotient-path certificate stated below, but the finite carrier does
not emit that certificate. The candidate law evaluates the mass-side pair
\((W,Z)\) together with the running-family anchor \((a_0,s_0,v)\). The
electromagnetic row is physically read only after Ward projection to the unbroken \(\U(1)_Q\)
channel; its low-energy value is the \(q^2\to0\) endpoint of the same transport family. The
derivation also includes two explicit benchmark surfaces beneath the public theorem output: the
exact selected-carrier chart and the reference-fitted coherent repair diagnostic. The same support
logic applies to flavor continuations, whose support levels differ, whereas the
electroweak quantitative lane has an explicit declared-surface contract. Its source
payload, same-scheme remainder, and interval-certificate records are public bookkeeping artifacts
for the declared bridge from \(P\) to the fine-structure endpoint.

The empirical closure surface for that bridge is populated. A documented piecewise
\(e^+e^-\to\mathrm{hadrons}\) compilation, built from resonance parameters and perturbative QCD,
evaluates the subtracted dispersion integral to
\(\Delta\alpha_{\mathrm{had}}^{(5)}(m_Z)=0.0267591200153\pm0.0007524955365\).
Inserting it into the endpoint map with the frozen source anchor and lepton transport packet gives
\(\alpha^{-1}=136.2636589431\) on
\([136.1583832834,136.3690975240]\), at \(P=1.6310415204\), on the
empirical-closure row class.
The comparison to the measured endpoint sits in an explicitly compare-only block: the measured
value lies outside the interval, and the certified same-scheme anchor gap is
\([0.6485541111,0.8547920666]\) inverse-alpha units. The anchor \(a_0(P)=\alpha_{\mathrm{em}}^{-1}(m_Z^2;P)\) is the
one-loop renormalization-group value run from the OPH unification scale to \(m_Z\); the physical
five-flavor on-shell value is \(128.939\), and the difference from the OPH anchor \(128.308\) is
\(0.631\), at the lower edge of the certified gap. The gap records transport missing from the
one-loop anchor; it is not evidence by itself that a source chain closes. A source-only bridge
that closes this gap reduces to the OPH hadronic spectral measure,
which needs a working hadron construction; the empirical-closure endpoint is the working surface in
the interim.
On the paper surface, the transmutation factor is \(\beta_{\mathrm{EW}}=N_c+1\); overloaded \(\beta\)-ratios appear only on benchmark readouts and are not part of the theorem contract.

\section{Results at a Glance}

The particle derivation carries the local pixel scale into the directly comparable rows below.
Detailed scope statements are collected in the support sections that follow.

\begin{table}[h]
\centering
\scriptsize
\setlength{\tabcolsep}{3pt}
\begin{tabular}{@{}>{\raggedright\arraybackslash}p{0.30\textwidth}>{\raggedright\arraybackslash}p{0.37\textwidth}>{\raggedright\arraybackslash}p{0.25\textwidth}@{}}
\toprule
Chain & OPH output & Role \\
\midrule
\(P\)-closure and fine structure & certified incomplete-map root \(136.994835177413\ldots\) (CL-6); mixed diagnostic \(137.0359595136\ldots\); certified gauge-width fixed point \(137.035660136946577\ldots\) (CL-2); empirical endpoint \(136.2636589431\in[136.1583832834,136.3690975240]\); measured \(137.035999177(21)\) & distinct support classes: incomplete declared-map output, mixed diagnostic excluded from physical output, empirical hadron closure with anchor gap \([0.6485541111,0.8547920666]\), and measurement \\
Classical carrier modes & two transverse Maxwell modes; \(2\dim G\) perturbative pure-Yang--Mills modes; two pure-Einstein TT modes & conditional quadratic-action results; quantum particle gates open \\
Electroweak \(W/Z\) & source-audit running/chart coordinates \((80.330,\,91.119)\)~GeV (closure row CL-5); candidate value-law coordinates \((80.3770000154,\,91.1879780779)\)~GeV; no source-only physical mass emitted & these coordinates are not commensurate with the PDG mass-dependent-width Breit--Wigner masses or with complex-pole masses; the stated complex-pole conversion diagnostic places the \(W\) chart \(0.5\) propagated experimental standard deviations from the converted PDG 2026 pole \(80.3340\)~GeV and the \(Z\) chart about \(17\) from its converted pole, with no theory covariance. QT1--QT5, the finite-carrier certificate, source-law selection, matching prescription, and physical pole readout are open; no numerical proximity is counted as mass evidence \\
Electroweak hierarchy/naturality & \(\begin{gathered}v/E_\star\ \text{on each named pixel branch}\\N_{\mathrm{EW}}\simeq3.53235\times10^{122}\\ \mathcal B_{\mathrm{EW}}=0,\ \epsilon_H=0\end{gathered}\) & selected dimensionless bridge; identifying \(N_{\mathrm{EW}}\) with cosmic \(F\) is conditional on CP-1 to CP-3. The Planck-\(\Lambda\) central capacity is \(N_\Lambda\simeq3.313\times10^{122}\), a \(6.6\)-percent mismatch; physical \(F\) and \(E_\star\) are separate gates \\
Higgs/top relation & double-criticality branch \((m_H,m_t)=(125.77,\,172.63)\)~GeV at the frozen boundary-scale candidate \(E_\star e^{-\pi}P^{-1/6}\) (two loops), with \(m_H=125.72\)~GeV on the fit-free curve at the measured top; target-anchored declared-surface fit \((125.1995304097,\,172.3523553288)\)~GeV kept separate; no source-only physical mass emitted & the criticality family has zero continuous parameters over the boundary scale; the boundary-scale selection is a theorem modulo two finite carrier facts CF1/CF2 (the variational midpoint principle is proved and its quadratic-cost and placement premises reduce to the axioms; CF1/CF2 are the D11 carrier census the \(W/Z\) law also needs); the declared-surface fit is target-anchored and never predicts; the source-root, scale, rigidity, provenance, uncertainty, and complex-pole gates are open \\
Charged leptons & target-anchored witness values withheld; finite eight-path digital carrier and empirical transport intervals retained as diagnostics & the model proves schema satisfiability and a central record dilation, not physical source selection; the intervals use target-anchored ratios, measured \(\alpha\), and empirical transport \\
Selected-frame quarks & public numeric rows withheld; reciprocal-ray product \(0.835323\) and held-out error \(21.56\%\) at \(M_Z\) & common-scale data reject the reciprocal-ray candidate; the generic interface has six scalars and requires a source-derived flavor-orbit selector \\
Neutrino absolute attachment & compare-only absolute masses withheld; scale-free and mixing comparison rows retained & weighted-cycle comparison branch \\
Hadrons & first-principles strong-binding descent plus empirical \(e^+e^-\to\mathrm{hadrons}\) closure surface & OPH hadron construction for first-principles masses; measured hadron data for empirical closure rows \\
H3 record localization and worldlines & certified cap-response localization balls followed by cross-boundary record-token continuation, or \(\mathrm{AMBIGUOUS}\) & conditional observer-record theorem; not a species, mass, or coupling input \\
\bottomrule
\end{tabular}
\end{table}

\begin{table}[h]
\centering
\tiny
\setlength{\tabcolsep}{2pt}
\begin{tabular}{@{}>{\raggedright\arraybackslash}p{0.14\textwidth}>{\raggedright\arraybackslash}p{0.11\textwidth}>{\raggedright\arraybackslash}p{0.19\textwidth}>{\raggedright\arraybackslash}p{0.20\textwidth}>{\raggedright\arraybackslash}p{0.27\textwidth}@{}}
\toprule
Sector & Support status & Derivation stage & Public output(s) / exact sidecar(s) & Caveat / theorem boundary \\
\midrule
Carrier roles and modes & structural group/content theorem plus conditional action theorem & explicit transverse/TT quadratic kernels with no promoted particle mass & action, background, and phase are hypotheses; quantization, positive-residue pole, and asymptotic/deconfinement gates are separate \\
Electroweak bosons & running/chart audit plus conditional quotient-transport theorem and comparison adapter & source-audit zero-selector law; selected carrier \(\rightarrow\) two-coordinate chart; QT1--QT5 \(\Rightarrow\) conditional value law & no nonzero mass in the prediction ledger; source-audit coordinates \((80.330,\,91.119)\)~GeV and candidate value-law coordinates \((80.3770000154,\,91.1879780779)\)~GeV are chart outputs & the chart coordinates are not commensurate with PDG mass-dependent-width Breit--Wigner masses or complex-pole masses. QT1--QT5, source-law selection, matching, scale, threshold, tadpole, and pole receipts are open \\
Electroweak hierarchy/naturality & selected-branch dimensionless theorem & local \(P\to\alpha_U\to v/E_\star\) hierarchy lane\newline \(\rightarrow\) product-adjoint \(m_{\rm rep}=24\)\newline \(\rightarrow\) mathematical capacity bridge and RG/Higgs square; conditional screen branch: 12 ports and an oriented 24-slot register & \(N_{\mathrm{EW}}=3.5323546226929906511\ldots\times10^{122}\), \(\mathcal B_{\mathrm{EW}}=0\), \(\epsilon_H=0\) & Supports the dimensionless hierarchy and naturality identities on the named branch. The two counts are aligned, not physically identified without a receipt.\newline Equality with cosmic capacity is conditional on CP-1 to CP-3 and an unconstructed \(F\); the Planck-\(\Lambda\) central value is lower by about \(6.6\) percent. An independent physical \(E_\star\), the public Thomson endpoint, theorem-level \(W/Z\), and the other listed mass rows are separate surfaces. \\
Higgs/top stage & double-criticality family from the gauge sector + declared-surface fit sidecar & criticality law \(\lambda=0,\ \beta_\lambda=0\) at one source scale \(\rightarrow\) fully constrained \((m_H,m_t)\) family \(\rightarrow\) frozen boundary-scale candidate & no nonzero mass in the prediction ledger; double-criticality branch \((m_H,m_t)=(125.77,\,172.63)\)~GeV at the frozen boundary-scale candidate, with \(m_H=125.72\)~GeV on the fit-free curve at the measured top; the target-anchored declared-surface fit \((125.1995304097,\,172.3523553288)\)~GeV stays in technical audit prose & Strict promotion inherits the two boundary-scale carrier facts CF1/CF2 (the midpoint selection theorem is proved and its premises reduce to the axioms modulo these) plus the source-root, physical-scale, QT1--QT5, RG/scheme, rigidity, provenance, and complex-pole gates.\newline The companion top coordinate is not a separate public top-mass prediction row. \\
Quark family & common-scale rejection plus restricted source non-identifiability theorem & one-scheme dimensionless Yukawas \(\rightarrow\) reciprocal-ray test\newline \(\rightarrow\) six-scalar generic interface and flavor-selector boundary & no public numeric quark row; at \(M_Z\), \(\rho_u\rho_d=0.835323\) and the endpoint-granted held-out error is \(21.56\%\) & The reciprocal-ray candidate is physically rejected across the tested running scales.\newline The \((\mathbb R_{>0})^2\) theorem is a restricted lower-bound obstruction.\newline Mixed-convention template and RSCC residuals are target-anchored diagnostics.\newline A flavor-orbit selector, quark--Higgs carrier, and common-scale source transport are absent \\
Charged leptons & continuation + exact sidecar witness & shared excitation dictionary \(\rightarrow\) ordered charged carrier \(\rightarrow\) exact centered readback \(\rightarrow\) determinant-line lift on theorem-level physical charged data & exact centered readback; same-family target-anchored charged triple withheld from public prediction tables & public charged masses are not emitted from \(P\).\newline The theorem lane does not emit a theorem-level sector-isolated charged determinant exponent vector.\newline It does not attach a source-side determinant character to the physical charged determinant line.\newline The determinant-line lift and algebraic mass readout apply only on theorem-level physical charged data \\
Neutrinos & target-informed continuation candidate + auxiliary checks & template family transport \(\rightarrow\) same-label scalar certificate \(\rightarrow\) weighted-cycle candidate \(\rightarrow\) compare-only bridge and absolute attachment \(\rightarrow\) shared-basis identity & frozen scale-free candidate and PMNS/Majorana comparison rows; absolute mass attachments withheld from public prediction tables & exact linear algebra conditional on the declared template and selectors.\newline The candidate fails the NuFIT 6.1 correlated \((\sin^2\theta_{23},\delta_{\mathrm{CP}})\) profile, and neither the scale-free point nor the bridge invariant has prediction or theorem status~\cite{nufit61} \\
Hadrons & strong-binding construction absent; empirical closure policy emitted & OPH-QCD backend contract, including quotient ensemble, source QCD parameter map, Ward current ledger, spectral exports, and empirical \(e^+e^-\to\mathrm{hadrons}\) payload schema & no first-principles hadron masses; empirical hadron closure rows are separate from first-principles OPH rows & First-principles hadron prediction requires a working OPH hadron construction, such as GLORB/Echosahedron, with the Ward-projected two-current spectral measure, higher-point/transition spectral sectors, same-scheme remainder, and systematics. The empirical closure surface uses separate \(e^+e^-\to\mathrm{hadrons}\) input and cannot upgrade the first-principles theorem \\
H3 record-worldline stitching & conditional certificate theorem & declared \(H^3_R\) chart atlas and observer clock atlas\newline \(\rightarrow\) cap-response localization certificate for each descended token and clock slice\newline \(\rightarrow\) localization balls with positive \(\Delta_{\mathrm{loc}}\) for point outputs\newline \(\rightarrow\) overlap descent\newline \(\rightarrow\) real transverse interface-crossing germs\newline \(\rightarrow\) common-chart sector/gauge transport\newline \(\rightarrow\) ID-independent one-to-one assignment gap\newline \(\rightarrow\) coarse/fine contraction check & nonbranching localized record-token paths, rays, or isolated records when both localization and stitch certificates pass; otherwise \(\mathrm{AMBIGUOUS}\), \(\mathrm{REJECTED}\), or \(\mathrm{INTERACTION\ REQUIRED}\) & This is a continuation theorem for observer-visible records in \(H^3\), not a derivation of particle species, masses, gauge charges, or scattering amplitudes. Repeated implementation IDs, nearest-neighbor fits, pre-existing coordinates without response inversion, and file-boundary coincidences are not admissible evidence. \\
\bottomrule
\end{tabular}
\end{table}

Two points are worth stating explicitly. First, the theorem rows and the exact-hit rows are not
the same object. The theorem table carries a two-modulus quark non-identifiability result and a
corpus-limited charged no-go boundary, while the exact-hit surface above it contains target-anchored
same-family witnesses and diagnostic sidecars.
Second, the electromagnetic, color, and metric carrier
roles have conditional classical-mode receipts but no promoted quantum mass rows; the \(W/Z\)
numbers are convention-dependent chart and adapter coordinates; and the
Higgs boson sits on the declared Higgs/top critical surface.

This support boundary governs the rest of the paper. Whenever a sector chapter gives more
detail about a continuation family, it should be read through this table. The table records the
constructive chain and keeps its support status explicit.

\paragraph{Non-hadron output surface.}
The paper-facing non-hadron bundle splits by lane as follows:

\begin{table}[h]
\centering
\scriptsize
\setlength{\tabcolsep}{3pt}
\begin{tabular}{@{}>{\raggedright\arraybackslash}p{0.14\textwidth}>{\raggedright\arraybackslash}p{0.17\textwidth}>{\raggedright\arraybackslash}p{0.31\textwidth}>{\raggedright\arraybackslash}p{0.30\textwidth}@{}}
\toprule
Lane & Exact output(s) & Exact chain on the paper surface & Caveat \\
\midrule
Carrier roles and classical modes & no photon/gluon/graviton particle-mass output & axioms \(\rightarrow\) realized electromagnetic/color/dynamical-metric roles; additional action/background/phase receipt \(\rightarrow\) transverse/TT quadratic modes & quantum-particle receipt not supplied; confined color has no free asymptotic-gluon claim \\
Electroweak chart and sidecar & no nonzero source-only mass output & source tuple \(\rightarrow\) exact two-coordinate chart; QT1--QT5 \(\Rightarrow\) candidate value law; inverse adapter separate & all numerical coordinates are confined to the audit discussion; QT1--QT5 are a finite-carrier certificate obligation and no pole pair is promoted \\
Hierarchy/naturality bridge & exact dimensionless bridge residual and zero Higgs naturality defect & axioms \(\rightarrow P\)-fixed source branch \(\rightarrow\) product-adjoint \(m_{\rm rep}=24\) \(\rightarrow\) mathematical \(N_{\mathrm{EW}}\) bridge map \(\rightarrow\) source-to-Higgs settled-form square; conditional screen branch: 12 ports and an oriented 24-slot register & selected-branch theorem for \(v/E_\star\) and the naturality identities; the equal counts are not physically identified without a receipt; equality with the cosmic capacity is conditional on CP-1 to CP-3 and an unconstructed \(F\); it does not supply a physical \(E_\star\) or promote the mass rows \\
Higgs/top exact sidecar & no nonzero source-only mass output & gauge core \(\rightarrow\) declared Jacobian \(\rightarrow\) exact inverse slice & benchmark inverse slice only; the conditional coordinates are outside the prediction ledger and inherit every open full-chain gate \\
Charged exact witness & exact same-family charged triple withheld & axioms \(\rightarrow\) shared excitation dictionary \(\rightarrow\) ordered charged carrier \(\rightarrow\) exact centered readback \(\rightarrow\) closed quadratic readout theorem \(\rightarrow\) same-family exact witness & same-family target-anchored audit witness only; theorem lane carries exact centered readback plus the closed common-shift no-go. The available derivation has a no-go boundary: a theorem-level \(\widehat C_e\) lift emitting the physical scalar \(\mu_{\mathrm{phys}}(Y_e)\) is not part of the emitted theorem surface \\
Quark physical boundary & numeric prediction rows withheld; common-scale reciprocal-ray failure leads the status; mixed-convention formula residuals are audit-only & six dimensionless Yukawas in one scheme and scale \(\rightarrow\) reciprocal-ray rejection \(\rightarrow\) six-scalar interface and selector no-go & at \(M_Z\), \(\rho_u=1.110889\), \(\rho_d=0.751941\), product \(0.835323\); even with four endpoints granted, the held-out miss is \(21.56\%\). The two-spread theorem is a restricted non-identifiability result \\
Neutrino candidate & frozen scale-free comparison point; absolute mass attachments withheld & template family transport \(\rightarrow\) same-label scalar certificate \(\rightarrow\) target-informed weighted-cycle law \(\rightarrow\) compare-only bridge and absolute attachment \(\rightarrow\) shared-basis identity & rejected by the NuFIT 6.1 correlated profile; target dependence, source closure, basis orientation, and absolute normalization fail promotion gates~\cite{nufit61} \\
\bottomrule
\end{tabular}
\end{table}

Concretely, the public values and benchmark checks shown on this surface are
\[
\alpha^{-1}(0)=137.035999177(21),
\qquad
P\simeq1.6309682094,
\]
\[
N_{\mathrm{EW}}
=
3.5323546226929906511\ldots\times10^{122},
\qquad
\mathcal B_{\mathrm{EW}}=0,
\qquad
\epsilon_H=0,
\]
This bridge value is about \(6.6\) percent above the Planck-\(\Lambda\) central capacity
\(N_\Lambda\simeq3.313\times10^{122}\), and the physical equality requires CP-1 to CP-3 plus a
constructed readback map \(F\). No nonzero particle mass appears in this prediction ledger. The selected-carrier, value-law,
inverse-adapter, and conditional Higgs coordinates appear only in the technical audit
sections below. Target-anchored charged-lepton
and quark witness values, including the mixed-convention quark target packet, appear only in
the exact-fit audit artifacts. Compare-only absolute neutrino attachments are withheld
from public prediction tables. The neutrino comparison surface keeps representative splitting
checks such as
\[
\begin{aligned}
\Delta m_{21}^2&=7.488059465106851\times10^{-5}\,\mathrm{eV}^2,\\
\Delta m_{31}^2&=2.5123118727618473\times10^{-3}\,\mathrm{eV}^2,\\
\Delta m_{32}^2&=2.4374312781107786\times10^{-3}\,\mathrm{eV}^2.
\end{aligned}
\]
The table is lane-based. It states the derivation chain that emits each exact-hit surface and the caveat that prevents a stronger theorem claim.

\paragraph{How to read mismatches and exact hits.}
Exact numerical agreement is not by itself a proof of a blind prediction. The provenance
record classifies the \(W/Z\) row as target-used frozen-reference reproduction, the charged-lepton
triple as an empirically anchored current-family witness, and the mixed-convention quark packet as a
target-anchored selected-frame audit rather than a public source-only prediction. The
same record classifies the hierarchy/naturality formula as a zero-residual
identity on its declared selected surface, with \(\epsilon_H=0\); the strict
source-root and physical-scale certificates are separate open gates. Where a
row does not match, or where it matches only on a constrained surface, the explanations are
as follows. The source/root certificate gives
\(\alpha_{\mathrm{root}}^{-1}=136.994835177413\ldots\) and
\(P_{\mathrm{fwd}}=1.630972095858897\ldots\) (CL-6, closed), but the declared map is
incomplete. The certified gauge-width map gives
\(\alpha^{-1}=137.035660136946577\ldots\), with CL-2 residual
\(2.5\times10^{-6}\) relative to the measured \(137.035999177(21)\). The no-hadron packet
that combines different pixel coordinates is no fixed point and is excluded from physical
output. Source-only closure requires Ward-projected hadronic spectral transport. The separate
endpoint-accounting diagnostic is
\[
\Delta^{\mathrm{fp}}_{\mathrm{H,cal}}
=\alpha_U(P_C)\,C_{24,Q},
\qquad
C_{24,Q}
=1.0009647859732326253849511140702475\ldots.
\]
The \(C_{24,Q}\) factor is endpoint accounting rather than a source-emitted theorem. At the public endpoint pixel value
\[
P\simeq1.6309682094
\]
the endpoint residual package requires
\[
\Delta_{\mathrm{source}}(P)=0.041465861005223389053448715357314044\ldots.
\]
That scalar belongs to the source-side hadronic spectral transport and
scheme-remainder map. The declared empirical route would use a separately labeled
\(e^+e^-\to\mathrm{hadrons}\) spectral input. No same-scheme integrated endpoint payload is emitted
here. The first-principles transport row excludes an inserted comparison endpoint. The
electroweak support record names the residual map and the
RG/matching/threshold/scheme packet.
An observer inside the branch measures the dressed Thomson coupling, not the
undressed source diagnostic \(1/136.994835\ldots\), because a zero-momentum
electromagnetic measurement sees the Ward-projected \(\U(1)_Q\) current after
charged-lepton vacuum polarization, confined-quark/hadron spectral transport,
and same-scheme finite endpoint matching have been included.
The electroweak \(W/Z\) row is therefore a chart and prescription audit, not a physical mass
benchmark. Charged leptons carry a corpus-limited
no-go because the determinant trace-lift attachment from the D10 descendants of \(P\) to physical charged data is absent. The auxiliary
direct-top PDG row differs from the theorem coordinate by \(0.20673301656674425~\mathrm{GeV}\),
or \(0.28458848947515303\) combined standard deviations, because Q007TP and Q007TP4 are distinct
extraction codomains; the direct-top response kernel has a corpus-limited no-go boundary. Neutrino absolute masses are
not directly measured, and PMNS-angle residuals are visible comparison tension outside the
theorem branch. First-principles hadron masses have no emitted prediction because the
required Ward-projected hadronic spectral measure must come from a working OPH
hadron construction. Empirical hadron closure rows carry a separate support class.

\paragraph{Local unification surface.}
The bosonic rows carry one cross-lane statement. The same pixel input \(P\), fixed on the
synthesis-paper outer/inner closure relation, fixes the
D10/D11 bosonic trunk
\[
P
\longmapsto
\alpha_U(P)
\longmapsto
\bigl(t_U(P),t_{\mathrm{tr}}(P)\bigr)
\longmapsto
v(P)
\longmapsto
\left(M_W,M_Z\right).
\]
\[
\sigma_{D11,\mathrm{HT}}
=
\alpha_U(P)\cos(2\theta_{W0})/\sqrt{\pi}
\longmapsto
\left(m_H,m_t\right).
\]
Here \(\alpha_U(P)\) is the branch value selected by the same forward D10 pixel-closure solve, so
the bosonic trunk contains no inverse electroweak readback of the internal transmutation data.
On the companion gravity side, the same pixel law packages
\[
\ellbar_{\mathrm{SU(2)}}(t_{2,\mathrm{run}})
\;+\;
\ellbar_{\mathrm{SU(3)}}(t_{3,\mathrm{run}})
=
P/4.
\]
The gravity normalization itself is emitted by the selected scale certificate
\(\gamma_\star=\ell_\star\nu_{\mathrm{Cs}}/c\), equivalently
\(B_\star=3\pi/\ell_\star^2\). With
\(a_{\mathrm{cell}}=P\ell_\star^2\), the Newton area law gives
\(G_{\mathrm{geom}}=\ell_\star^2\), so the pixel cancels.
The local unification surface separates one explicit familiar-unit readout package from three
support surfaces: the \(W/Z\) benchmark sidecar, the conditional D11
declared-surface Higgs/top coordinate,
and the gravity-side readout with its strict classical-regime clause.
On the gravity side, the stated local extension surface uses the lifted product presentation of the
realized quotient branch and identifies
\[
\ellbar_{\mathrm{shared}}
=
\ellbar_{\mathrm{SU(2)}}(t_{2,\mathrm{run}})
\;+\;
\ellbar_{\mathrm{SU(3)}}(t_{3,\mathrm{run}}).
\]
On that same surface the D10 pixel law fixes \(\ellbar_{\mathrm{shared}}=P/4\),
and the local SI readout is
\[
G_{\mathrm{SI}}=\frac{c^3\ell_\star^2}{\hbar}
\]
from the selected scale certificate. The \(\chi_\nu\) protected-reserve branch
uses the same public endpoint convention as
\[
P_\chi=P_{\mathrm C}=1.630968209403959\ldots ,
\qquad
\epsilon_{\rm res}=\frac{P_\chi}{24},
\]
after the same-collar shared-edge budget and one-class \(\mathbb Z_6\) trace
receipts pass. This does not make \(P/4\) a primitive Hilbert-space dimension:
\(P/4\) is the local screen-cell entropy budget in the particle and hierarchy
lanes, while \(P_\chi/24\) is the shared scalar-reserve density used by the
\(\chi_\nu\) collar branch. On that declared extension surface the same
familiar-unit package reads
\[
L_{\mathrm{loc}}=\sqrt{a_{\mathrm{cell}}}\,\widehat L(P),
\qquad
t_{\mathrm{loc}}=\frac{\sqrt{a_{\mathrm{cell}}}}{c}\,\widehat T(P),
\]
\[
E_{\mathrm{loc}}=\frac{\hbar c}{\sqrt{a_{\mathrm{cell}}}}\,\widehat E(P),
\qquad
\Theta_{\mathrm{loc}}=\frac{\hbar c}{k_B\sqrt{a_{\mathrm{cell}}}}\,\widehat\Theta(P),
\]
with dimensionless \(\widehat L,\widehat T,\widehat E,\widehat\Theta\). Thus, at fixed \(P\), the
local ruler is \(\sqrt{a_{\mathrm{cell}}}\); seconds are that ruler divided by the structural
Lorentz output \(c\); and GeV and Kelvin are downstream familiar-unit displays of the inverse
local ruler through \(\hbar\) and \(k_B\). On that declared extension
surface the local scale-readout bundle is
\[
\begin{aligned}
c&=299792458\,\mathrm{m/s},\\
G&=6.674299995910528\times10^{-11}\,\mathrm{m^3\,kg^{-1}\,s^{-2}},
\end{aligned}
\]
No nonzero particle-mass coordinate belongs to this prediction ledger. The \(W/Z\) inverse adapter
and the conditional Higgs/top coordinates are audit-only quantities discussed in their technical
sections. The structural Lorentz output is the common invariant null cone and speed
\(c_\star\). The decimal \(299792458\,\mathrm{m/s}\) is exact by the SI definition of the metre,
not a predicted magnitude. The paper keeps the
\(G\) row as an exact scale-readback value, with no literal zero-difference identity against the
rounded benchmark \(6.6743\times10^{-11}\).

The detailed inventory below records the public theorem/continuation rows sector by sector. Whenever an exact sidecar or same-family witness is stronger than the public theorem row, the caveat is the one stated in the exact-hit table above.

\subsection{Detailed particle inventory}

\tiny
\setlength{\tabcolsep}{2pt}
\begin{longtable}{@{}>{\raggedright\arraybackslash}p{0.135\textwidth}>{\raggedright\arraybackslash}p{0.145\textwidth}>{\raggedright\arraybackslash}p{0.135\textwidth}>{\raggedright\arraybackslash}p{0.215\textwidth}>{\raggedright\arraybackslash}p{0.265\textwidth}@{}}
\toprule
Family & Particle & Support status & OPH value or strongest derived benchmark & Open condition \\
\midrule
\endfirsthead
\toprule
Family & Particle & Support status & OPH value or strongest derived benchmark & Open condition \\
\midrule
\endhead
\bottomrule
\endfoot
Conditional carrier modes & electromagnetic & action/phase receipt & two transverse classical \(k^2=0\) modes; \(\mu_{Q,\mathrm{quad}}^2=0\) only in the displayed Maxwell action & physical Hilbert space, positive-residue pole, and stable asymptotic photon receipt absent \\
Conditional carrier modes & color & perturbative/pre-confinement action receipt & \(2\dim G\) transverse classical \(k^2=0\) modes; \(\mu_{\mathrm{YM},\mathrm{quad}}^2=0\) only in the displayed pure-Yang--Mills expansion & positive-residue pole and deconfined asymptotic colored sector absent; no free gluon is claimed on the confining branch \\
Conditional carrier modes & tensor & action/background receipt & two Einstein transverse-traceless classical \(k^2=0\) modes; \(\mu_{\mathrm{EH},\mathrm{TT}}^2=0\) only on the pure-Einstein branch & metric quantization, positive physical Hilbert space, positive-residue pole, and asymptotic/EFT particle interpretation absent \\
Electroweak branch & \(W\) boson & no source-only physical mass & withheld from the prediction ledger & selected-carrier, conditional value-law, and inverse-calibration coordinates are audit-only; no pole receipt \\
Electroweak branch & \(Z\) boson & no source-only physical mass & withheld from the prediction ledger & same boundary as the \(W\) row \\
Higgs/top critical stage & Higgs boson & no source-only physical mass & withheld from the prediction ledger & independent source root and physical scale, QT1--QT5, RG/scheme and threshold transport, target-independent D11 rigidity, complex-pole, uncertainty, and no-target dependency-DAG gates open \\
Quark family & top quark & separate target-audit coordinate & withheld from public prediction table & cross-section pole-mass extraction coordinate; it is not a sixth entry in one common running-mass chart and is not selected by the source spread laws \\
Quark family & up quark & reciprocal-ray candidate rejected & withheld from public prediction table & common-scale dimensionless Yukawa test fails; a source-derived flavor-orbit selector is absent \\
Quark family & down quark & reciprocal-ray candidate rejected & withheld from public prediction table & same common-scale rejection and six-scalar interface boundary \\
Quark family & strange quark & held-out reciprocal-ray failure & withheld from public prediction table & part of the \(19.9\%\) to \(21.6\%\) held-out failure across tested scales \\
Quark family & charm quark & held-out reciprocal-ray failure & withheld from public prediction table & stored GeV-valued matrices are mixed-convention mass textures, not physical Yukawas \\
Quark family & bottom quark & reciprocal-ray candidate rejected & withheld from public prediction table & common-scale source transport, Higgs normalization, and flavor-orbit selection absent \\
Charged leptons & electron & continuation gap & \(n/a\); exact same-carrier centered readback exists once a charged source pair is emitted, but the absolute scale is blocked by the closed common-shift no-go; benchmark target \(g_e^\star=0.0457789\), equivalently \(\Delta_e^{\mathrm{abs},\star}=3.00398633\) & close branch-generator splitting, then emit the lift whose descended scalar is \(\mu_{\mathrm{phys}}(Y_e)\); within that lift the exact smaller forcing object is the physical identity-mode equalizer, after which \(\widetilde C_e(Y_e)=\widehat C_e(Y_e)+\mu_{\mathrm{phys}}(Y_e)\,\mathbf 1\), \(A_{\mathrm{ch}}(Y_e)=\mu_{\mathrm{phys}}(Y_e)\), and the readouts \(g_e\), \(\Delta_e^{\mathrm{abs}}\) follow \\
Charged leptons & muon & continuation gap & \(n/a\); same exact centered-readback / common-shift-no-go frontier as the electron row & same \(\widehat C_e^{\mathrm{cand}} \rightarrow\) branch-generator splitting \(\rightarrow\) lift \(\rightarrow \mu_{\mathrm{phys}}(Y_e) \rightarrow A_{\mathrm{ch}} \rightarrow g_e\) closure chain, with the physical identity-mode equalizer beneath the descended scalar \\
Charged leptons & tau lepton & continuation gap & \(n/a\); same exact centered-readback / common-shift-no-go frontier as the electron row & same \(\widehat C_e^{\mathrm{cand}} \rightarrow\) branch-generator splitting \(\rightarrow \mu_{\mathrm{phys}}(Y_e) \rightarrow A_{\mathrm{ch}} \rightarrow g_e\) closure chain, with the physical identity-mode equalizer beneath the descended scalar \\
Neutrino comparison & declared \(f\)-basis weighted-cycle matrix & rejected target-informed candidate & no flavor-particle mass row; the frozen declared-basis coordinate gives \(\theta_{12}=34.2259^\circ\), \(\theta_{23}=49.7228^\circ\), \(\theta_{13}=8.68636^\circ\), \(\delta=305.581^\circ\), \(J=-0.02753\), and, under its declared normal-ordering labels, \(\Delta m_{21}^2/\Delta m_{32}^2=0.03072111\) & the upstream flavor kernel is a hand-written template; the physical charged basis and mass-label rule are open; the target-informed selector precludes a blind-prediction claim; the NuFIT 6.1 correlated profile rejects the candidate; Majorana and absolute-attachment values are compare-only \\
Hadrons & proton & strong-binding construction absent; empirical closure policy emitted & no first-principles emitted prediction & first-principles prediction requires a working OPH-QCD hadron backend, such as GLORB/Echosahedron, plus Ward-projected two-current, higher-point, and transition spectral exports, same-scheme remainder, and production systematics; empirical closure rows use separate \(e^+e^-\to\mathrm{hadrons}\) input \\
Hadrons & neutron & strong-binding construction absent; empirical closure policy emitted & no first-principles emitted prediction & same OPH construction gate, with isospin-resolved hadron systematics in the construction support class \\
Hadrons & neutral pion proxy & strong-binding construction absent; empirical closure policy emitted & no first-principles emitted prediction & same OPH construction gate; local stable-channel surrogate output is not a paper prediction \\
Hadrons & \(\rho(770)^0\) proxy & strong-binding construction absent; empirical closure policy emitted & no first-principles emitted prediction & same OPH construction gate, plus finite-volume resonance extraction in the construction support class \\
\end{longtable}
\normalsize

\section{Detailed Closure Values, Screen Architecture, and Theorem Packages}\label{sec:particle-input-ledger}

This section records the detailed closure-value and theorem checklist used by the particle derivation: the five axioms,
the screen-capacity closure, the local pixel closure, the regulated screen realization, and the imported
OPH theorem packages, ordered as axioms, screen language, closure values, and theorem packages.

\subsection{Canonical OPH basis}

The particle-spectrum derivation uses the same five axioms as the compact reconstruction
paper~\cite{ophcompact}. They are restated here for local readability because every
particle-family chapter depends on them either directly or through the structural OPH chain
imported from that paper.

\begin{axiom}[Screen Net]\label{ax:screen}
Physical data is organized on a horizon screen \(S^2\) carrying a net of local algebras
\[
P \longmapsto \mathcal A(P)
\]
for connected patches \(P\subset S^2\), with isotony
\[
P\subset Q \implies \mathcal A(P)\subset \mathcal A(Q).
\]
\end{axiom}

\begin{axiom}[Overlap Consistency]\label{ax:overlap}
For overlapping patches \(P_1\cap P_2\neq\varnothing\), the local states induced on the shared
algebra agree:
\[
\omega_{P_1}|_{\mathcal A(P_1\cap P_2)}
=
\omega_{P_2}|_{\mathcal A(P_1\cap P_2)}.
\]
\end{axiom}

\begin{axiom}[Local MaxEnt and Refinement Stability]\label{ax:maxent}
At the regulator scale \(\ell_{\mathrm{UV}}\), the realized branch is selected by maximizing
entropy subject to the finitely many homogeneous global-sum constraints
\[
\Bigl\langle\sum_x O_a(x)\Bigr\rangle=C_a,
\qquad a=1,\dots,N_{\mathrm{con}},
\]
built from gauge-invariant local densities of support radius \(O(\ell_{\mathrm{UV}})\); there is
one constraint and one Lagrange multiplier per density label \(a\), not per regulator cell, so the
multiplier count is cutoff-independent by construction. The axiom additionally asserts a
refinement-closure clause: under refinement, the coarse-grained realized state again lies in the
exponential family generated by the same finite density list at the coarser scale. Closure is a
substantive renormalization condition, since coarse-graining a finite-range Gibbs family
generically generates interactions outside any fixed finite list. The closure defect, the
moment-matching I-projection realizing the induced refinement map on the multiplier space, and
the Pinsker trace-norm residual bound controlling a nonzero defect are constructed in the
companion SM/GR derivation and synthesis papers~\cite{ophcompact,ophsynthesis}; the clause is
exactly the statement that this defect vanishes along the realized branch, and that vanishing is
assumed there rather than proved. Granting it, the realized
states belong to one common finite-dimensional MaxEnt family instead of unrelated maximizers.
\end{axiom}

\begin{axiom}[Recoverable Generalized Entropy]\label{ax:entropy}
A generalized entropy functional exists on caps,
\[
S_{\mathrm{gen}}(C)=S_{\mathrm{bulk}}(C)+\langle L_C\rangle,
\]
where \(L_C\) is a positive edge-center entropy functional and the semiclassical branch
identifies its leading coarse-grained contribution with \(A(\partial C)/(4G)\). The functional
obeys the recoverability and focusing structure required by the collar and null-modular
arguments.
\end{axiom}

\begin{axiom}[Minimal Admissible Realization]\label{ax:mar}
Among admissible realized low-energy sector packages \(\mathfrak S\) consisting of the connected
Lie gauge-sector image relevant in the EFT branch, its admissible light chiral matter content,
and one Higgs doublet, the realized package is lexicographically minimal under
\[
C(\mathfrak S)=\bigl(\chi_{\mathrm{cpl}},\,N_{\mathrm{nonab}},\,N_c,\,N_g\bigr),
\]
subject to loop coherence, anomaly freedom, refinement-stable light chiral matter,
single-Higgs Yukawa completable structure with one connected abelian charge factor, intrinsic
CP capability, and weak-sector UV completeness.
\end{axiom}

For particle physics, the first four axioms provide the observer-centric kinematic and
entropic/modular background. The fifth axiom, MAR, is the selector that turns ``some compact
gauge group reconstructed from the transportable edge-sector category on a cofinal tail carrying
the compact paper's explicit refinement receipt, coherent pullback ladder, symmetry, and
forgetful-fiber conditions'' into the realized Standard Model branch. This
paper therefore uses MAR constantly, but only after the compact-gauge reconstruction step has
been separated from realized-branch selection.

\subsection{Regulated screen architecture and patch language}

The SM/GR derivation in Ref.~\cite{ophcompact}
states the axioms in algebraic language. Ref.~\cite{ophmicrophysics} supplies one concrete
regulated realization of that language: a federation of finite observer patches with echosahedral
multi-port interfaces, recurrent toroidal subchannels, exposed overlap packets, record algebras,
and local repair instruments. A spherical screen is an observer-visible regulator chart, not a
literal microscopic computer. This regulated architecture is important for the particle derivation
for three reasons.

First, it makes the screen picture operational. A patch is a finite local
algebra in a bounded carrier, and an overlap is a boundary-visible algebra with declared shared
observables. Second, it makes the edge-sector
language concrete. The particle derivation uses edge sectors, transport, refinement,
fusion, and overlap data to build the gauge and flavor branches; Ref.~\cite{ophmicrophysics} shows
how those objects can be realized at fixed cutoff by finite gauge-aware patch carriers.
Third, it clarifies the measurement interface. The observer-facing measurement package is a
theorem-bearing fixed-cutoff statement about a central record algebra, Born probabilities for its
event projectors, and L\"uders conditioning on that same commuting algebra.

This regulated architecture should be read as a reference implementation, not as the unique OPH
UV completion. Ref.~\cite{ophmicrophysics} is explicit on that boundary: it selects a federated
patch-carrier architecture because it is local, finite-dimensional, gauge-aware, naturally
compatible with collars, overlap observables, and repair maps, and does not ask the reader to
identify the spherical regulator chart with the literal substrate. Hardware evidence is imported
only when it appears in a public OPH evidence bundle with stable hashes and verifier receipts.
This paper inherits that architecture as a concrete screen-side realization of the OPH language,
while keeping microscopic uniqueness separate from the support-visible continuum theorem surfaces
supplied by the compact paper.

\subsection{Quantitative inputs and where they enter the particle derivation}

The quantitative derivation developed here uses one incomplete declared-map coordinate and one
conditional capacity coordinate together with the
selected no-\(G\) scale certificate:
\begin{align}
P &\equiv a_{\mathrm{cell}}/\ell_\star^2,\\
N_{\mathrm{CRC}} &\equiv \log \dim \mathcal H_{\mathrm{tot}},\\
\gamma_\star&\equiv\frac{\ell_\star\nu_{\mathrm{Cs}}}{c},
\qquad
B_\star\equiv\frac{3\pi}{\ell_\star^2}.
\end{align}

These closure values have distinct downstream dependencies.

\paragraph{Pixel area \(P\).}
The pixel area is the declared local UV-area ratio. In this derivation it feeds the heat-kernel /
edge-law input stage and, through that route, the forward D10 electroweak surface. It is the
common upstream numerical variable for the electroweak, Higgs/top, flavor, quark, and
hadron-facing branches. The particle paper therefore treats \(P=P_\star\) as the inherited
root of the incomplete outer/inner declared map
\[
P_\star=\varphi+\frac{\sqrt{\pi}}{A_T(P_\star)}
\]
from the synthesis paper~\cite{ophsynthesis}, not as a quantity derived again inside the
particle-spectrum argument itself.

\paragraph{Conditional screen capacity \(N_{\mathrm{CRC}}\).}
The screen capacity is not part of the local recovered-core gravity/gauge derivation. It enters
the separate screen-capacity branch through the cosmic record-closure target
\[
N_{\mathrm{CRC}}=F(N_{\mathrm{CRC}}),
\qquad
\Lambda_{\mathrm{CRC}}\ell_\star^2=\frac{3\pi}{N_{\mathrm{CRC}}},
\qquad
\Lambda_{\mathrm{CRC}}=\frac{3\pi}{G N_{\mathrm{CRC}}}.
\]
The physical readback map \(F\) is not constructed. Identifying the electroweak bridge with this
cosmic fixed point requires CP-1 to CP-3. The bridge value near
\(3.532\times10^{122}\) is about \(6.6\) percent above the Planck-\(\Lambda\) central capacity
near \(3.313\times10^{122}\). The last equality is therefore a conditional scale-certified
display; it does not determine an SI curvature scale from \(P_\star\) and
\(N_{\mathrm{CRC}}\) alone.
The separate declared count-density selector is
\[
N_\star=\arg\max_N\left[\log|\Omega^{\mathrm{sc}}_N|-N\right],
\]
with its interior-maximum and smoothness premises declared inputs.
If \(\Omega^{\mathrm{sc}}_N\) is independently constructed, then with
\(\ell(N)=\log|\Omega^{\mathrm{sc}}_N|-N\), the map
\(T_\eta(N)=N+\eta\ell'(N)\) has a unique stable mathematical fixed point when the normalized
density is strictly concave. No proof identifies this candidate with the physical \(F\).
In the particle context, this matters only for the cosmological-capacity discussion that
frames why local null data do not determine the
cosmological constant. The implemented branch uses the static-patch normalization
\[
N_{\mathrm{patch}}=\left(\frac{r_{\mathrm{dS}}}{\ell_P}\right)^2\simeq1.05\times10^{122},
\qquad
N_{\mathrm{scr}}=\pi N_{\mathrm{patch}}\simeq3.313\times10^{122}.
\]
This capacity bookkeeping supplies no neutrino result; the weighted-cycle neutrino derivation is
a separate lane.

\paragraph{Selected scale certificate.}
The gravity scale is a separate observation-located certificate. On the declared branch,
\[
\gamma_\star=\frac{\ell_\star\nu_{\mathrm{Cs}}}{c},
\qquad
B_\star=\frac{3\pi}{\ell_\star^2},
\qquad
G_{\mathrm{SI}}=\frac{c^3\ell_\star^2}{\hbar}.
\]
If the declared-map root acquires its missing physical identification, the local pixel
\(P_\star\) identifies \(a_{\mathrm{cell}}=P_\star\ell_\star^2\) and
\(\ellbar_{\mathrm{shared}}=P_\star/4\). In the Newton area law \(P_\star\) cancels, so this paper
does not treat \(G\) as a particle-sector consequence of \(P_\star\).
\(\ellbar_{\mathrm{shared}}\) is a shared-cut density, not the logarithm of an autonomous cell
Hilbert factor. The particle branch uses it for cell/edge consistency; it does not select a fixed
primitive observer capacity.

\paragraph{Why these are the declared continuous closure/readback values.}
One of the discipline constraints of the OPH derivation is that particle-sector freedom should not
be hidden in a large uncontrolled parameter set. This paper exposes the incomplete declared-map
root \(P_\star\), the proposed \(N_{\mathrm{CRC}}\) readback fixed point with \(N_\star\) as a
separate count candidate, and the
selected scale certificate, then asks
the overlap, modular, gauge, and admissibility machinery to do the rest. The continuation branches
state their compliance with that discipline explicitly. For the particle-spectrum branch, the nontrivial
common quantitative burden sits on \(P_\star\): \(N_{\mathrm{scr}}\) enters the separate
capacity/cosmology side and \(\gamma_\star\), equivalently \(B_\star\), enters the separate
gravity-scale side, whereas the
reverse-engineering claim for masses and couplings is that one shared declared-map coordinate should
drive the spectrum instead of a long sector-by-sector input list.

\subsection{Technical theorem data carried into the particle-spectrum derivation}

This paper uses the same theorem checklist as
\emph{Recovering Relativity and the Standard Model from Observer Overlap Consistency}~\cite{ophcompact}, together with
the regulator-language clarification integrated into the consensus and microphysics
appendices. The ledger matters only
insofar as it separates branch-internal statements from declared inputs.

The third axiom internalizes several pieces of infrastructure that often appear as separate
assumptions: the local Gibbs form of the regulator-scale state, the
quasi-local propagation bound, the endpoint-control estimate for bounded intervals, and the
meaning of refinement stability itself. In regulator language, one works with finite local
Hilbert spaces and a boundary-fixed overlap action on cut data, often summarized as the R0/R1
presentation. Beyond that regulator package, the compact paper supplies the support-visible BW
scaling theorem for the Lorentz/null-modular/Einstein branch. Transportability and the fixed-cutoff bosonic sector category are theorem-produced in the compact paper; the refinement/fiber
ladder and cofinal MAR-admissible compact-gauge witness additionally require its explicit compact-gauge refinement receipt. On the central branch the matching gluing
obstruction is the combined zero-obstruction transport criterion: vanishing triangle defect plus
trivial represented holonomy of the strictified edge \(1\)-cocycle. It introduces no second theorem-side input. Any
local-Gibbs or collar-mixing language on that branch is fixed-cutoff recoverability/support
control; it is not the source of the compact-gauge witness.
For the BW/geometric side, the target is the support-visible extracted prime geometric subnet.
The compact paper proves the needed scaling theorem by combining regularized support-visible
modular transport, weak-\(*\)/GNS extraction of the cap pair, support-readable modular covariance,
BW framing, held-out oriented cross-ratio rigidity, and independently normalized geometric
\(2\pi\)-KMS convergence with wrong-scale controls. The unregularized full-algebra common-floor
route is deliberately not claimed, because off-support directions may collapse without affecting
observer-facing matrix elements. Bare finite consensus does not by itself produce the finite
cap-normal BW certificate used by that theorem.

For the particle-spectrum derivation, this has a practical consequence. Structural statements such
as compact-gauge reconstruction, the realized Standard Model quotient, the exact hypercharge
lattice, the color count \(N_c=3\), and the generation count \(N_g=3\) are not floating
independently of the declared input and theorem checklist. The additional classical carrier modes
live on explicit quadratic action/background/phase receipts, and particle poles require the
stronger quantum receipt. These statements live on a
controlled chain whose scope conditions need to be visible in this paper,
especially for quantitative closure, continuation, and
nonperturbative sectors.

\subsection{Imported core theorem packages}

This paper uses the following OPH theorem packages without re-proving them in full.

\begin{enumerate}[leftmargin=2em]
\item \textbf{Patch-net and overlap language.} Ref.~\cite{ophconsensus} organizes overlap repair,
fixed-point language, cycle holonomy, and gauge-as-gluing in a form that is especially useful
for flavor transport and observer-centric interpretation.
\item \textbf{Screen-side regulated architecture.} Ref.~\cite{ophmicrophysics} gives the
federated echosahedral patch-carrier model, the regulated patch-net embedding theorem, and the
fixed-cutoff edge heat-kernel / Casimir theorem.
\item \textbf{Measurement interface.} Ref.~\cite{ophmicrophysics}, together with the integrated
measurement appendices and the synthesis paper \emph{Observers Are All You Need}~\cite{ophsynthesis}, supplies the fixed-cutoff central-record, Born-rule, and
L\"uders-conditioning package together with the fixed-cutoff Bell/CHSH theorem stack used when
the measurement sections explain how observations pick out definite observer-accessible outcomes and how
two-wing comparison laws stay on the quantum side of the classical Bell bound.
\item \textbf{Compact gauge reconstruction and Standard Model structure.}
\emph{Recovering Relativity and the Standard Model from Observer Overlap Consistency}~\cite{ophcompact}
and the dedicated gauge-group fragment together provide the path from a receipt-certified cofinal
tail of refinement-stable edge sectors to a compact group, then from MAR to
\[
\frac{\SU(3)\times\SU(2)\times\U(1)}{\mathbb Z_6},
\qquad
N_g=3,
\qquad
N_c=3.
\]
\item \textbf{Conditional topic notes.} The integrated lane appendices carry useful material
on supersymmetry, the finite-quotient baryogenesis source theorem and its open record-generator branch, proton stability, generation structure, and Yukawa hierarchy,
but those sections have to be read with their stated claim boundaries intact. They are sources
for conditional analyses, not recovered-core theorems.
\end{enumerate}

This reading rule governs the whole paper. Structural gauge and carrier results can be used
as structural results. The electroweak branch is a declared quantitative calibration branch. The
Higgs/top critical stage is a conditional downstream calculation on its declared running, matching, and threshold surface. The quark,
charged-lepton, neutrino, and hadron chapters have conditional or comparison status because
their derivation chains are incomplete.

\section{Gauge Architecture, Standard Model Structure, and Structural Carriers}

The particle-spectrum derivation does not start from a list of particles. It starts from the
compact-gauge branch. On a cofinal tail carrying the explicit refinement receipt, OPH reconstructs
a compact gauge group from the tensor-generated zero-obstruction edge-sector category and
forgetful fiber produced by the compact paper's transportability, fixed-cutoff category, and
refinement/fiber theorems. The overlap/transport obstruction calculus is the classification stage,
not the selection stage. The receipt-conditional compact-gauge witness theorem supplies nonempty
realized MAR-admissible branch data, and MAR then selects the
realized low-energy package. That logical order matters for this paper. The
gauge branch fixes the structural carrier roles before any mass readout. It does not fix their
quantum-particle masses; the action-level and quantum spectral gates below decide whether a
classical mode or particle row exists.

On the realized branch, the gauge/content result used throughout this paper is
\[
\frac{\SU(3)\times\SU(2)\times\U(1)}{\mathbb Z_6},
\qquad
N_g=3,
\qquad
N_c=3.
\]
\emph{Recovering Relativity and the Standard Model from Observer Overlap Consistency}~\cite{ophcompact}, the longer main derivation surface, and the dedicated gauge-group proof
fragment all agree on the logical split: compact-gauge reconstruction first, realized Standard
Model selection second. The first stage gives a compact internal symmetry group under the stated
zero-obstruction bosonic sector conditions and compact-gauge refinement receipt. The second stage uses MAR and anomaly cancellation to fix the realized
quotient and hypercharge lattice; the minimal coupled carrier fixes \(N_c=3\); intrinsic CP
capability together with weak-sector UV completeness bounds \(N_g\); MAR then fixes \(N_g=3\);
and Witten parity is a consistency check on the realized triplet-doublet package.

For the particle zoo this structural branch fixes three carrier roles, not three particle masses.
The compact paper's carrier-mode theorem then supplies the following conditional action-level
statements. The Maxwell action on the ordinary unbroken electromagnetic vacuum has two transverse
classical modes. The pure Yang--Mills action about a flat connection has \(2\dim G\) perturbative
transverse modes, but the confined color phase has no asserted gauge-invariant asymptotic gluon.
The pure two-derivative Einstein--Hilbert action about Minkowski space has two classical
transverse-traceless modes. In every case a quantum particle additionally requires a positive-
energy physical Hilbert space, a positive-residue physical two-point pole, and the appropriate
asymptotic or deconfinement receipt.

\begin{table}[h]
\centering
\small
\begin{tabular}{@{}>{\raggedright\arraybackslash}p{0.13\textwidth}>{\raggedright\arraybackslash}p{0.12\textwidth}>{\raggedright\arraybackslash}p{0.12\textwidth}>{\raggedright\arraybackslash}p{0.23\textwidth}>{\raggedright\arraybackslash}p{0.26\textwidth}@{}}
\toprule
Carrier role & Claim tier & Action-level output & Physical degrees of freedom & Quantum-particle boundary \\
\midrule
Electromagnetic connection & conditional Maxwell mode & transverse \(k^2=0\) kernel & two transverse classical modes & photon requires quantization and a positive-residue physical pole \\
Color connection & conditional perturbative pure-Yang--Mills mode & one transverse \(k^2=0\) kernel per generator & \(2\dim G\) perturbative modes & deconfined BRST/LSZ receipt required; no free-gluon claim in confined QCD \\
Metric perturbation & conditional pure-Einstein mode & \(D^{\mathrm{TT}}\propto i\Pi^{\mathrm{TT}}/(k^2+i0)\) & two classical TT modes & graviton requires metric quantization, physical Hilbert space, and pole receipt \\
\bottomrule
\end{tabular}
\end{table}

These conditional classical modes should not be confused with quantum-particle mass rows. The
group and gravity branches identify the carrier roles; the displayed actions and phases supply
their classical propagation; and only a separate quantum spectral receipt can promote them to
photon, gluon, or graviton particles. The boson sections address the \(W\), \(Z\), and Higgs
rows on their own electroweak and Higgs/top quantitative stages.

\section{Electroweak D10 Branch and the Higgs/Top Critical Stage}

The reported bosonic sector distinguishes sharply between emitted numerical rows and scoped
non-emitted rows. The active bosonic stages are the electroweak
D10 prediction stage and the Higgs/top critical stage, while charged-lepton rows
carry their no-go boundary, the neutrino lane contains a rejected weighted-cycle comparison candidate, and hadron rows are backend-gated. The bosonic numerical
sector is therefore concentrated in two places: the D10 electroweak branch for the \(W\) boson and \(Z\) boson, and
the Higgs/top critical stage for the Higgs boson together with its conditional D11 companion top coordinate.

\subsection{Electroweak chart on the D10 quantitative branch}

The electroweak derivation consists of a single-\(P\) running family followed by a
reduced two-scalar carrier, a selector on that carrier, an exact carrier mass
chart, and a conditional quotient-transport theorem beyond the selected
carrier. In practical terms, the construction starts from
the declared pixel input \(P\), builds the running electroweak family, reduces it to the selected
carrier, reads the \(W\) boson/\(Z\) boson pair from the selected point on that carrier,
and then evaluates the mass pair together with the Ward-projected electromagnetic transport family.
On the D10 branch, a source-only prediction must respect that ordering: first certify
the shared pixel input \(P\) on the declared D10 running/matching/threshold/scheme surface, thereby fixing the source basis
\[
(\alpha_U,\alpha_{2,m_Z},\alpha_{Y,m_Z},\eta_{\mathrm{source}},v),
\]
with \(\alpha_U(P)\), equivalently \(t_U(P)\) and \(t_{\mathrm{tr}}(P)\), fixed by the
same forward pixel-closure solve,
and only then emit the electroweak transport family from that source basis. The theorem tier does
not contain this full chain. Its \(W/Z\) rows are data-comparison adapter coordinates,
not a D10 prediction theorem.

The same quantitative lane has one exact golden-ratio benchmark. If one writes
\[
x(C):=\frac{S_{\mathrm{gen}}(C)}{S_{\mathrm{bulk}}(C)}
=
1+\frac{\langle L_C\rangle}{S_{\mathrm{bulk}}(C)}
\]
for the total/bulk/edge hierarchy and imposes exact self-similar balance
\[
\frac{S_{\mathrm{gen}}(C)}{S_{\mathrm{bulk}}(C)}
=
\frac{S_{\mathrm{bulk}}(C)}{\langle L_C\rangle},
\]
then \(x\) obeys \(x^2-x-1=0\), so the unique positive equilibrium point is
\[
x=\varphi:=\frac{1+\sqrt5}{2}.
\]
Equivalently the order parameter \(A_\varphi(x)=x-1-\frac1x\) vanishes there. The synthesis paper
records the unique numerical root of the incomplete outer/inner declared map. The point of
the equilibrium theorem here is that the proximity to \(\varphi\) has a structural reason. Exact
balance is too symmetric to support durable records, structure, and dynamics, so the declared map
assigns a small
equilibrium-breaking detuning away from that balance point.

The compare-only \(W/Z\) adapter rows in the final bundle are
\[
m_W = 80.3625~\mathrm{GeV},
\qquad
m_Z = 91.1879~\mathrm{GeV}.
\]
These numbers sit on the frozen validation surface, not on a target-free prediction theorem.
The exact selected-carrier chart is explicit on disk and emits
the local pair
\[
m_W^{\mathrm{carrier}} = 80.38629169244275~\mathrm{GeV},
\qquad
m_Z^{\mathrm{carrier}} = 91.18290444674243~\mathrm{GeV},
\]
so the mathematics distinguishes the selected-carrier chart, the candidate
value-law candidate \((80.3770000154,91.1879780779)\,\mathrm{GeV}\), the
conditional QT1--QT5 implication, and the measured-reference inverse repair surface.
The rounded pair \((80.377,91.1879780919)\,\mathrm{GeV}\) is a boundary alias
of the candidate, not the inverse adapter. None is a physical pole theorem.

\begin{theorem}[Measured-pair reparametrization (calibration): coherent electroweak inverse fit]
Let the emitted running/core electroweak basis be
\[
Q_{\mathrm{run}}(P)=
\bigl(\alpha_{Y,m_Z}(P),\alpha_{2,m_Z}(P),v(P),\eta_{\mathrm{source}}(P)\bigr),
\]
and write
\[
\tau_Y(\tau_2)
=
-\frac{\tau_2+2\eta_{\mathrm{source}}}{1+4\tau_2^2},
\qquad
n_{\mathrm{fiber}}(\tau_2)
=
1+\frac{\alpha_{Y,m_Z}\tau_Y(\tau_2)+\alpha_{2,m_Z}\tau_2}
{\alpha_{Y,m_Z}+\alpha_{2,m_Z}}.
\]
Freeze one authoritative target pair
\[
\mathcal T^\dagger=(M_W^\dagger,M_Z^\dagger),
\qquad 0<M_W^\dagger<M_Z^\dagger,
\]
and define the charged anchor
\[
\tau_{2,W}^\dagger
=
\frac{M_W^{\dagger\,2}}{\pi v^2\alpha_{2,m_Z}}-1,
\qquad
\delta\alpha_2^\dagger
=
\frac{M_W^{\dagger\,2}}{\pi v^2}-\alpha_{2,m_Z}.
\]
Then the fiber-parallel hypercharge leg is
\[
\tau_Y^\dagger=\tau_Y(\tau_{2,W}^\dagger),
\qquad
n_{\mathrm{fiber}}^\dagger=n_{\mathrm{fiber}}(\tau_{2,W}^\dagger),
\]
\[
M_{Z,\mathrm{fiber}}^\dagger
=
v\sqrt{\pi(\alpha_{Y,m_Z}+\alpha_{2,m_Z})\,n_{\mathrm{fiber}}^\dagger},
\]
and the orthogonal neutral-shear scalar is
\[
\delta M_Z^\perp=M_Z^\dagger-M_{Z,\mathrm{fiber}}^\dagger,
\]
equivalently
\[
\delta n^\dagger
=
\frac{(M_Z^\dagger+M_{Z,\mathrm{fiber}}^\dagger)\,\delta M_Z^\perp}
{\pi v^2(\alpha_{Y,m_Z}+\alpha_{2,m_Z})},
\qquad
\delta\alpha_Y^\perp
=
\frac{(M_Z^\dagger+M_{Z,\mathrm{fiber}}^\dagger)\,\delta M_Z^\perp}{\pi v^2}.
\]
The fiber-parallel hypercharge motion is
\[
\delta\alpha_Y^\parallel
=
\alpha_{Y,m_Z}\,
\frac{8\eta_{\mathrm{source}}(\tau_{2,W}^\dagger)^2-\tau_{2,W}^\dagger}
{1+4(\tau_{2,W}^\dagger)^2}.
\]
Therefore the unique coherent repair package is
\[
\Sigma_{EW}^\dagger
=
\bigl(\delta\alpha_2^\dagger,\delta\alpha_Y^\parallel,\delta\alpha_Y^\perp\bigr),
\]
equivalently \((\tau_{2,W}^\dagger,\delta n^\dagger)\). If the selected carrier anchor is the
\(\tau_2=0\) fiber point, so that
\[
\alpha_{2,*}=\alpha_{2,m_Z},
\qquad
\alpha_{Y,*}=\alpha_{Y,m_Z}(1-2\eta_{\mathrm{source}}),
\]
then the coherent validation couplings are
\[
\alpha_{2,\dagger}=\alpha_{2,*}+\delta\alpha_2^\dagger,
\qquad
\alpha_{Y,\dagger}=\alpha_{Y,*}+\delta\alpha_Y^\parallel+\delta\alpha_Y^\perp,
\]
and they emit one coherent quintet
\[
M_W^\dagger=v\sqrt{\pi\alpha_{2,\dagger}},
\qquad
M_Z^\dagger=v\sqrt{\pi(\alpha_{Y,\dagger}+\alpha_{2,\dagger})},
\]
\[
\alpha_{\mathrm{em},\dagger}^{-1}
=
\frac{\alpha_{Y,\dagger}+\alpha_{2,\dagger}}
{\alpha_{Y,\dagger}\alpha_{2,\dagger}},
\qquad
\sin^2\theta_{W,\dagger}
=
\frac{\alpha_{Y,\dagger}}{\alpha_{Y,\dagger}+\alpha_{2,\dagger}}.
\]
The frozen-target electroweak validation law is one unique coherent value-emission law on top of
the closed fiber map.
\end{theorem}

For the reference-fitted inverse-adapter audit surface,
\[
\begin{aligned}
\alpha_{2,\star}&=0.03377843630219015,\\
\alpha_{Y,\star}&=0.009682831911900495,\\
\eta_{\mathrm{source}}&=0.022147000871961295,
v&=246.76711732749683~\mathrm{GeV},\\
M_W^\dagger&=80.3625~\mathrm{GeV},\\
M_Z^\dagger&=91.1879~\mathrm{GeV},
\end{aligned}
\]
and the law emits
\[
\begin{aligned}
\tau_{2,W}^\dagger&=-0.0005918464744071317,\\
\delta\alpha_2^\dagger&=-1.9991648436440412\times10^{-5},\\
M_{Z,\mathrm{fiber}}^\dagger&=91.16822266259587~\mathrm{GeV},\\
\delta M_Z^\perp&=19.6773374041328~\mathrm{MeV},\\
\delta\alpha_Y^\parallel&=5.9969727526906094\times10^{-6},\\
\delta\alpha_Y^\perp&=1.8756950557734103\times10^{-5},\\
\delta n^\dagger&=4.2716772021678714\times10^{-4}.
\end{aligned}
\]
So the coherent validation couplings are
\[
\begin{aligned}
\alpha_{2,\dagger}&=0.03375844465375371,\\
\alpha_{Y,\dagger}&=0.00970758583521092,
\end{aligned}
\]
and the same coherent branch emits
\[
M_W^\dagger=80.3625~\mathrm{GeV},
\qquad
M_Z^\dagger=91.1879~\mathrm{GeV},
\]
\[
\begin{aligned}
\alpha_{\mathrm{em},\dagger}^{-1}&=132.6344411210187,\\
\sin^2\theta_{W,\dagger}&=0.22333729871365943.
\end{aligned}
\]
These two carrier-side scalars are bookkeeping data on the frozen validation surface. They are not
the public electromagnetic readout on the Ward-projected D10 lane.
That inverse-fit law provides an exact coherent calibration surface: it is a
measured-pair reparametrization. On that
target-frozen surface the reference \(W/Z\) pair is hit exactly because the basis is solved from
that pair; the hit is the defining condition, and the lane is calibration.
It is excluded from every prediction ledger.
The displayed electroweak lane is the forward running/chart solve,
\((80.330,\,91.119)\)~GeV (closure ledger row CL-5). The PDG comparison values are
mass-dependent-width Breit--Wigner parameters, while complex-pole masses define another
convention. As a stated convention diagnostic under the complex-pole conversion of the
PDG 2026 reference values, the chart \(W\) coordinate sits \(0.5\) propagated experimental
standard deviations from the converted pole \(80.3340\)~GeV, and the chart \(Z\)
coordinate about \(17\) from its converted pole; both figures carry no OPH theory
covariance, and the physical readout contract is open. No physical pull or
near-hit is assigned to these coordinates.
The candidate forward value law has a stronger conditional theorem: it follows
uniquely from a finite quotient-path certificate.  The certificate is not
emitted by the five axioms or by the existing carrier artifacts.

\begin{definition}[D10 quotient-path certificate]\label{def:particle-d10-qt}
After quotienting gauge representatives, port labels, scheduler data, and
other hidden presentation coordinates, require:
\begin{enumerate}[leftmargin=2em,label=\textup{QT\arabic*.}]
\item the real, CP-even, color-singlet, charge-preserving response through two
returns is exactly \(\mathbb Rq_2\oplus\mathbb Rq_n\), with no third scalar or
charged-neutral off-block term;
\item the primitive response is bilinear in
\((\eta_{\mathrm{source}},\alpha_U)\), one primitive event has fixed unit
response normalization, and the quotient probability measure gives each of
the four transmutation slots weight \(1/4\), yielding
\(\lambda_{\mathrm{EW}}=\eta_{\mathrm{source}}\alpha_U/4\);
\item explicit carrier path lists, carrying one factor of
\(\eta_{\mathrm{source}}\) per return and uniform color-orbit measure \(1/3\),
have primitive, one-return, and two-return incidences
\((1,2/3,1)\) for \(q_2\) and \((1,4/3,2)\) for \(q_n\), while the diagonal
\(\mathbb Z_6\) projector in the declared response representation has
normalized trace \(1/6\), source amplitude \(\rho_{\mathrm{EW}}\), and subtracts
\((\rho_{\mathrm{EW}}/6)\eta_{\mathrm{source}}^2\) from both degree-two
characters;
\item mismatch descent fixes the charged sign as contracting and the neutral
sign as uplifting, and the parallel hypercharge fibre is the least-norm
minimizer with Gram factor \(1+4\tau_2^2\) and residual pairing
\(\tau_2+2\eta_{\mathrm{source}}\);
\item the listed paths exhaust the admissible class: deeper paths,
alternative central weights, and mixed terms are absent, quotient-trivial, or
separated by a positive target-independent MAR gap.
\end{enumerate}
\end{definition}

\begin{theorem}[Conditional D10 quotient-transport value law]\label{thm:particle-conditional-d10-qt}
Let the D10 source-only emitted basis be
\[
Q_{\mathrm{src}}(P)=
\bigl(\alpha_U(P),\alpha_{2,m_Z}(P),\alpha_{Y,m_Z}(P),
\eta_{\mathrm{source}}(P),v(P)\bigr),
\]
and define
\[
\rho_{\mathrm{EW}}
:=
\frac{\alpha_{2,m_Z}-\alpha_{Y,m_Z}}{\alpha_{2,m_Z}+\alpha_{Y,m_Z}},
\lambda_{\mathrm{EW}}
:=
\frac{\eta_{\mathrm{source}}\alpha_U}{4}.
\]
Assume that every entry is emitted on the same source branch, that
\(\eta_{\mathrm{source}}=\rho_{\mathrm{EW}}\alpha_U\), and that no measured
\(W/Z\), measured \(v\), fitted electroweak parameter, or calibrated proxy is
an ancestor of the tuple.  Hence
\(\lambda_{\mathrm{EW}}=\eta_{\mathrm{source}}^2/(4\rho_{\mathrm{EW}})\).
Assume also the finite certificate of Definition~\ref{def:particle-d10-qt} and the
physical-domain inequalities
\(\alpha_2'>0\), \(\alpha_Y'>0\).
Then the beyond-selected-carrier transport map is uniquely
\[
\tau_2^{\mathrm{exact}}
=
-\lambda_{\mathrm{EW}}
\left(
1+\frac23\eta_{\mathrm{source}}
+
\left(1-\frac{\rho_{\mathrm{EW}}}{6}\right)\eta_{\mathrm{source}}^2
\right),
\]
\[
\delta n^{\mathrm{exact}}
=
\lambda_{\mathrm{EW}}
\left(
1+\frac43\eta_{\mathrm{source}}
+
\left(2-\frac{\rho_{\mathrm{EW}}}{6}\right)\eta_{\mathrm{source}}^2
\right),
\]
with fiber hypercharge transport
\[
\tau_Y^{\mathrm{fiber}}
=
-\frac{\tau_2^{\mathrm{exact}}+2\eta_{\mathrm{source}}}
{1+4(\tau_2^{\mathrm{exact}})^2}.
\]
The coherent transport shifts are
\[
\delta\alpha_2=\alpha_{2,m_Z}\tau_2^{\mathrm{exact}},
\qquad
\delta\alpha_Y^{\parallel}
=
\alpha_{Y,m_Z}
\frac{8\eta_{\mathrm{source}}(\tau_2^{\mathrm{exact}})^2-\tau_2^{\mathrm{exact}}}
{1+4(\tau_2^{\mathrm{exact}})^2},
\]
\[
\delta\alpha_Y^{\perp}
=
(\alpha_{2,m_Z}+\alpha_{Y,m_Z})\,\delta n^{\mathrm{exact}}.
\]
With
\[
\alpha_{2,*}=\alpha_{2,m_Z},
\qquad
\alpha_{Y,*}=\alpha_{Y,m_Z}(1-2\eta_{\mathrm{source}}),
\]
the coherent transport couplings are
\[
\alpha_2'=\alpha_{2,*}+\delta\alpha_2,
\qquad
\alpha_Y'=\alpha_{Y,*}+\delta\alpha_Y^{\parallel}+\delta\alpha_Y^{\perp},
\]
and they emit one coherent D10 mass chart
\[
M_W^{(10)}=v\sqrt{\pi\alpha_2'},
\qquad
M_Z^{(10)}=v\sqrt{\pi(\alpha_2'+\alpha_Y')},
\qquad
v=v(P).
\]
\[
a_0(P)
:=
\frac{\alpha_{2,m_Z}(P)+\alpha_{Y,m_Z}(P)}
{\alpha_{2,m_Z}(P)\alpha_{Y,m_Z}(P)},
\qquad
s_0(P)
:=
\frac{\alpha_{Y,m_Z}(P)}
{\alpha_{2,m_Z}(P)+\alpha_{Y,m_Z}(P)}.
\]
Thus the certificate fixes the candidate mass pair together with the
running-family anchor \((a_0,s_0,v)\) from the D10 basis. The physical electromagnetic row
is read from the Ward-projected unbroken \(\U(1)_Q\) channel in
Theorem~\ref{thm:ward_u1q_transport}.
\end{theorem}

\begin{proof}
QT1 reduces every admissible response to the two displayed coordinates.  QT2
fixes their common primitive activity.  QT3 gives the charged and neutral path
characters, including the common \(\mathbb Z_6\) subtraction, and QT4 fixes
their signs.  The fibre functional
\[
\mathcal J_Y(t)
=\frac12(1+4(\tau_2^{\mathrm{exact}})^2)t^2
+(\tau_2^{\mathrm{exact}}+2\eta_{\mathrm{source}})t
\]
is strictly convex, so its unique minimizer is
\(-(\tau_2^{\mathrm{exact}}+2\eta_{\mathrm{source}})
 /(1+4(\tau_2^{\mathrm{exact}})^2)\).
Substitution gives the repaired couplings and the one-Higgs mass chart.  QT5
excludes every output-changing admissible deformation, which proves uniqueness
on the certified class.
\end{proof}

Using the same frozen electroweak validation basis
\[
\begin{aligned}
\alpha_{2,m_Z}&=0.03377843630219015,\\
\alpha_{Y,m_Z}&=0.010131601067241624,\\
\alpha_U&=0.04112498041477454,\\
\eta_{\mathrm{source}}&=0.022147000871961295,\\
v&=246.76711732749683~\mathrm{GeV},
\end{aligned}
\]
the conditional quotient-transport law evaluates to
\[
M_W^{(10)}=80.37700001539531~\mathrm{GeV},
\qquad
M_Z^{(10)}=91.18797807794321~\mathrm{GeV},
\]
\[
\begin{aligned}
a_0&=128.30576920234813,\\
s_0&=0.23073542347506173.
\end{aligned}
\]
This specialization verifies the candidate value-law algebra. It does not
derive QT1--QT5, establish that this calibration tuple belongs to the strict
source-audit pixel branch, or identify \(W,Z\) as complex propagator poles.
A machine-checkable promotion receipt must provide the finite path lists,
quotient canonicalizer, exact incidence and central-trace checks, fibre Gram
calculation, deformation enumeration or positive MAR gap, and a same-branch
no-target dependency DAG.

\begin{proposition}[Local nondegeneracy of the two-output D10 mass chart]
Define
\[
n_{\mathrm{fib}}(\tau_2)
:=
1+\frac{\alpha_2\tau_2+\alpha_Y\tau_Y(\tau_2)}
{\alpha_2+\alpha_Y}.
\]
On the physical domain
\(1+\tau_2>0\) and
\(n_{\mathrm{fib}}(\tau_2)+\delta n>0\),
\[
\det
\frac{\partial(M_W^{(10)},M_Z^{(10)})}
{\partial(\tau_2,\delta n)}
=
\frac{M_W^{(10)}M_Z^{(10)}}
{4(1+\tau_2)\,[n_{\mathrm{fib}}(\tau_2)+\delta n]}
>0.
\]
Holding \((\alpha_2,\alpha_Y,v,\eta)\) fixed, the two-output chart is locally
nondegenerate.  This determinant neither proves QT1 nor excludes an additional
physical response channel.  Within the displayed two-coordinate chart, the
missing content is a source selector rather than another fit coordinate.
\end{proposition}

\begin{proof}
\(M_W^{(10)}\) is independent of \(\delta n\), while
\[
\frac{\partial M_W^{(10)}}{\partial\tau_2}
=\frac{M_W^{(10)}}{2(1+\tau_2)},
\qquad
\frac{\partial M_Z^{(10)}}{\partial\delta n}
=\frac{M_Z^{(10)}}{2[n_{\mathrm{fib}}(\tau_2)+\delta n]}.
\]
The determinant is the product of these diagonal entries.
\end{proof}

\begin{theorem}[Ward-projected \texorpdfstring{\(\U(1)_Q\)}{U(1)_Q} transport law and Thomson-limit electromagnetic readout]\label{thm:ward_u1q_transport}
Let the D10 source-only running family be locked by the forward pixel solve on the declared
running/matching/threshold/scheme surface, with source basis
\[
Q_{\mathrm{src}}(P)=
\bigl(\alpha_U(P),\alpha_{2,m_Z}(P),\alpha_{Y,m_Z}(P),
\eta_{\mathrm{source}}(P),v(P)\bigr).
\]
Assume the realized D9 branch, so that
\[
Q=T_3+Y
\]
on the realized low-energy branch and color-singlet physical states carry integer \(Q\)-charge.
Assume further that the post-selector electroweak transport object is the populated quotient-local
kernel
\[
K^{\mathrm{EW}}_{D10}(q^2;P)
=
\bigl\{\Pi_{AA}(q^2;P),\Pi_{AZ}(q^2;P),\Pi_{ZZ}(q^2;P),\Pi_{WW}(q^2;P)\bigr\}
\]
on the declared physical observable algebra, together with a Ward projector
\[
\mathcal W_Q:
K^{\mathrm{EW}}_{D10}(q^2;P)\longrightarrow \Pi_Q(q^2;P)
\]
onto the unbroken electromagnetic channel satisfying at \(q^2=0\)
\[
\mathcal W_Q[\Pi_{AZ}(0;P)]=0,
\qquad
\mathcal W_Q[\Pi_{AA}(0;P)]\neq 0.
\]
Assume the abelian projected edge-sector probabilities obey the \(\U(1)\) heat-kernel law
\[
p_n(q^2;P)\propto e^{-t_Q(q^2;P)\lambda_n},
\qquad
\lambda_n=n^2,
\]
equivalently
\[
g_Q^2(q^2;P)=\frac{t_Q(q^2;P)}{2\pi},
\]
and that the scalar readout package satisfies the exact provenance lock
\[
\begin{aligned}
\text{family\_source\_id} &= \texttt{d10\_running\_tree},\\
\text{scheme\_id} &= \texttt{freeze\_once},\\
\text{origin\_kernel\_id} &= \texttt{EWTransportKernel\_D10}.
\end{aligned}
\]
Then the unique electromagnetic coupling readout on the Ward-projected D10 lane is
\[
\alpha_{\mathrm{em}}^{-1}(q^2;P)=\frac{8\pi^2}{t_Q(q^2;P)}.
\]
If
\[
a_0(P):=\alpha_{\mathrm{em}}^{-1}(m_Z^2;P),
\]
then the unique zero-normalized affine scalar on the same source-locked family is
\[
\delta_\alpha(q^2;m_Z^2;P)
:=
\frac{t_Q(m_Z^2;P)}{t_Q(q^2;P)}-1,
\]
so that
\[
\alpha_{\mathrm{em}}^{-1}(q^2;P)
=
a_0(P)\bigl(1+\delta_\alpha(q^2;m_Z^2;P)\bigr).
\]
Hence the Thomson-limit electromagnetic readout on that same
Ward-projected electromagnetic kernel:
\[
\alpha_{\mathrm{Th}}^{-1}(P)
:=
\lim_{q^2\to0}\alpha_{\mathrm{em}}^{-1}(q^2;P)
=
a_0(P)\,\frac{t_Q(m_Z^2;P)}{t_Q(0;P)}.
\]
Consequently, the low-energy electromagnetic row on this lane is read as the Thomson endpoint of
the D10 electromagnetic transport family instead of as a mass-chart byproduct.
\end{theorem}

\begin{proof}
The realized D9 branch fixes the electric charge operator by \(Q=T_3+Y\), and the
quotient-protected charge theorem fixes the physical color-singlet \(Q\)-lattice to be integral.
The abelian edge-sector theorem supplies the heat-kernel extraction rule
\[
p_n\propto e^{-t_Q\lambda_n},
\qquad
g_Q^2=\frac{t_Q}{2\pi},
\]
so the Ward-projected \(Q\)-channel determines a unique transport time \(t_Q(q^2;P)\). Because
the Ward projector kills \(A\) - \(Z\) mixing at \(q^2=0\), the selected channel is the physical
unbroken electromagnetic channel. Then
\[
\alpha_{\mathrm{em}}=\frac{g_Q^2}{4\pi}=\frac{t_Q}{8\pi^2},
\]
hence
\[
\alpha_{\mathrm{em}}^{-1}=\frac{8\pi^2}{t_Q}.
\]
Evaluating at \(q^2=m_Z^2\) defines the source-locked anchor \(a_0(P)\), and dividing by the same
expression at general \(q^2\) yields
\[
\delta_\alpha=\frac{t_Q(m_Z^2)}{t_Q(q^2)}-1.
\]
The provenance-equality clause forces one running-family source, one frozen scheme, and one origin
kernel for all scalar readouts, so no mixed-family or \(Z\)-only surrogate is promotable. The
Thomson formula is the \(q^2\to0\) limit of the same identity.
\end{proof}

\begin{corollary}[Conditional Ward-projected Maxwell normalization]\label{cor:ward_maxwell_normalization}
On the same Ward-projected \(\U(1)_Q\) lane, assume in addition the low-energy Maxwell action with
the displayed nonzero kinetic normalization. Then the electromagnetic field strength and current obey
\[
F_Q=dA_Q,\qquad dF_Q=0,\qquad d{*}F_Q=g_Q^2(q^2;P){*}J_Q,
\]
with
\[
g_Q^2(q^2;P)=\frac{t_Q(q^2;P)}{2\pi},
\qquad
\alpha_{\mathrm{em}}^{-1}(q^2;P)=\frac{8\pi^2}{t_Q(q^2;P)}.
\]
At the Thomson endpoint, this is the Maxwell normalization used by the public
fine-structure row.
\end{corollary}

\begin{proof}
The compact reconstruction supplies the unbroken abelian electromagnetic connection direction
\(\U(1)_Q\) with charge generator \(Q=T_3+Y\), but does not supply the Maxwell kinetic action.
On an abelian factor the compact-gauge
curvature restricts to \(F_Q=dA_Q\), so \(d^2=0\) gives \(dF_Q=0\). Varying the quadratic
electromagnetic action in the OPH coupling convention gives \(d{*}F_Q=g_Q^2{*}J_Q\). The theorem
above identifies the same branch coupling as \(g_Q^2=t_Q/(2\pi)\), hence
\(\alpha_{\mathrm{em}}^{-1}=8\pi^2/t_Q\).
\end{proof}

\begin{corollary}[Source-only admissibility criterion on the Ward-projected D10 lane]
The Ward-projected D10 lane supports a source-only fine-structure row only when the populated
Ward-projected transport kernel satisfies
\[
\lim_{q^2\to0}\alpha_{\mathrm{em}}^{-1}(q^2;P)=\alpha_{\mathrm{Th}}^{-1}(P),
\]
on the same source-locked family and without any separate endpoint value feeding that lane.
Concretely, the evidence packet must contain
\(\texttt{WARD\_Q\_SPECTRAL\_MEASURE\_RECEIPT}\),
\(\texttt{J24Q\_SOURCE\_JACOBI\_RECEIPT}\),
\(\texttt{OMEGAQ\_NORMALIZATION\_RECEIPT}\),
\(\texttt{XI\_SAME\_SCHEME\_REMAINDER\_RECEIPT}\),
\(\texttt{NO\_THOMSON\_TARGET\_LEAK\_DAG}\),
\(\texttt{FULL\_PIXEL\_CONTRACTION\_INTERVAL}\), and
\(\texttt{ALPHA\_PREDICTION\_BEATS\_DIRECT\_MEASUREMENT}\) under the declared
direct-measurement interval. Without those objects, the public endpoint is an empirical
hadron-closure row.
\end{corollary}

This criterion is the two-current marginal of the larger OPH-QCD hadronic precision
backend. It is adequate for the Thomson endpoint and HVP branch only if the same source law also
records the QCD quotient ensemble, source parameter map, Ward current normalization, systematics
ledger, and no-target-leak DAG. Full hadronic precision claims require the higher-point and
transition spectral exports from that same backend rather than reusing the scalar fine-structure
residual.

The endpoint table records the source-locked anchor from the runtime candidate,
\[
a_0(P)=128.3079654732862482099611087417567\ldots,
\]
and the OPH-plus-empirical hadron-closure value
\[
\alpha_{\mathrm{emp}}^{-1}(0)=136.2636589431201,
\qquad
\alpha_{\mathrm{emp}}^{-1}(0)\in[136.1583832834267,136.3690975239657].
\]
The repository's source-derivation witness path records the source-side audit trunk without
a built-in reference inverse-\(\alpha\) value. The measured endpoint
\(137.035999177(21)\) is excluded as a source-only closure-solve input and is a separate
comparison coordinate.
The source-side audit trunk emits the certified value
\(\alpha_{\mathrm{root}}^{-1}=136.994835177413\ldots\) (CL-6, closed), so its difference
from the endpoint readout is an endpoint-and-matching audit packet.
This certificate establishes a unique root of the incomplete declared numerical map. It does not
derive a relation between that root and the physical Thomson endpoint.

\begin{remark}
The empirical closure row is a Thomson-limit inverse fine-structure evaluation using external
hadron data. The source anchor
\(a_0(P)=128.3079654732862482099611087417567\ldots\) is the electroweak-scale running-family value at
\(m_Z^2\), while the measured row is separate. The 2022 CODATA/NIST reference value is
\(\alpha^{-1}=137.035999177(21)\)~\cite{codata2022}.
\end{remark}

The measured-reference inverse repair surface provides diagnostic validation and is excluded from
the prediction ledger.
The companion carrier-side pair
\[
\begin{aligned}
\alpha_{\mathrm{em},\dagger}^{-1}&=132.6344411210187,\\
\sin^2\theta_{W,\dagger}&=0.22333729871365943
\end{aligned}
\]
is mass-chart bookkeeping on that frozen validation surface. It is not the public electromagnetic
readout on the Ward-projected \(\U(1)_Q\) lane.

The electroweak quantitative-closure branch also contains several subordinate theorem objects: an
underdetermination theorem for the quadratic repair family, the smallest theorem route
through \(\mathrm{ColorBalancedQuadraticRepairDescent}_{EW}\), the candidate mass-emitter
contract beneath the quarantined \(W/Z\) audit rows, and the Ward-projected electromagnetic transport theorem above
the source-locked running-family anchor. The scalar-package consistency clause is the shared
provenance lock
\[
(\texttt{d10\_running\_tree},\texttt{freeze\_once},\texttt{EWTransportKernel\_D10}).
\]
A further forward transmutation certificate sits directly beneath the candidate readout and
records that the runtime source-side basis reconstructs the same \(\alpha_U(P)\), \(t_U(P)\), and
\(t_{\mathrm{tr}}(P)\) as the pixel-closure solve without reading them back from measured
couplings. This runtime separation does not exclude target dependence. These objects locate the
candidate repair formula inside the quantitative branch; they do not supply QT1--QT5 or a
physical pole receipt.

\subsection{Hierarchy interpretation on the local transmutation branch}

The selected branch certifies a dimensionless electroweak
hierarchy/naturality relation. The large
cosmic horizon ratio and the electroweak hierarchy belong to different OPH branches.
The proposed global capacity branch \(N_{\mathrm{CRC}}\), conditional on CP-1 to CP-3 and a
constructed \(F\), accounts for the cosmic area and horizon-length
hierarchy. The local pixel/transmutation readout is the ratio
\[
\frac{v}{E_\star}
=
P_\star^{-1/2}
\exp\!\left[-\frac{2\pi}{4\alpha_U(P_\star)}\right].
\]
The hierarchy certificate records the interval
\[
I_U=[0.041123336195630494,\;0.041125336195630496],
\]
the Krawczyk inclusion
\[
K(I_U)\subset
[0.0411243357185544983,\;0.0411243366727064662]
\subset \mathrm{int}(I_U),
\]
and a strictly negative derivative enclosure
\([-10.995768,-10.985284]\).  Thus the local source equation has a unique zero in
that enclosure.  On the public endpoint branch,
\[
\begin{aligned}
\alpha_U(P_\star)&=0.041124336195630495,\\
\frac{v}{E_\star}&=2.0199803239725553\times10^{-17}.
\end{aligned}
\]
The source-audit branch excludes the public Thomson endpoint upstream and has certified forward
pixel \(P_{\mathrm{fwd}}=1.630972095858897\ldots\). These two named branches
must not be merged.  The dimensionless ratio is generated by the ordinary
exponential transmutation factor controlled by the unified diffusion coupling;
a weak scale in GeV additionally requires an independently source-closed
\(E_\star\). The
twenty-four-round capacity contraction belongs to the global horizon-length ratio.

This is also the OPH reading of Higgs naturalness on the declared D10/D11 surface. The
observer-visible scalar mass is the normal-form readout
\[
m_H=H_{\mathrm{OPH}}(P_\star)
\]
on the selected branch. The split into a bare Higgs mass plus a cutoff correction is a regulator
coordinate split outside the observer-facing physical variable. A scalar relevant deformation absent
from the retained branch constraints is off-branch; a scalar deformation present on the selected
branch has its value fixed by the branch equations. The claim inherits the D10/D11 quantitative
surface and the \(\mathcal R_U\) precision policy. The selected exact source-to-Higgs
coarse-graining square has \(\epsilon_H=0\), with
\(\epsilon_H\in[0,0]\), as stated below.

\subsection{Local/global resonance continuation for the hierarchy}

The local/global hierarchy continuation compares the D10 transmutation exponent with a proposed
global capacity branch at the same joint pair \((P_\star,N_{\mathrm{CRC}})\). Conditional on
CP-1 to CP-3 and a constructed physical map \(F\), the target relation is
\[
t_{\mathrm{tr}}(P_\star)
=
\frac{P_\star}{12}
\log\!\left(\frac{N_{\mathrm{CRC}}}{\pi}\right),
\]
or equivalently
\[
\frac{v}{E_{\mathrm{cell}}}
=
\left(\frac{N_{\mathrm{CRC}}}{\pi}\right)^{-P_\star/12}.
\]
Writing
\[
|g_\star'|
=
\left(\frac{N_{\mathrm{CRC}}}{\pi}\right)^{-1/48}
\]
turns the same expression into
\[
\frac{v}{E_{\mathrm{cell}}}
=
|g_\star'|^{4P_\star}.
\]
This packages the mathematical electroweak hierarchy bridge as a resonance between the local
declared-map root and a conditional global screen-capacity coordinate. That screen branch supplies
twelve curvature ports and an oriented 24-slot register. Independently, the hierarchy exponent uses
the product-adjoint count \(m_{\rm rep}=24\). The equal counts are aligned on this branch; a physical
identification requires a separate receipt.

The exponent count is fixed by the observer-visible product adjoint. On the realized
product branch,
\[
m_{\rm rep}
=2\dim(\mathfrak{su}(3)\oplus\mathfrak{su}(2)\oplus\mathfrak u(1))
=2(8+3+1)
=24.
\]
The factor \(2\) orients each product-adjoint channel and does not follow from the screen register. The
\(\mathfrak{su}(5)\) adjoint has the same single-orientation integer for a different support:
its \(X/Y\) mixed gauge channels are absent from the OPH product branch.

\begin{lemma}[Conditional global repair-tick lemma]\label{lem:global-repair-tick}
Assume CP-1 to CP-3, a constructed physical map \(F\), and a value
\(N_{\mathrm{CRC}}=F(N_{\mathrm{CRC}})\) with
observed-branch readout \(N_{\mathrm{CRC}}=S_{\mathrm{dS}}\) and D6 normalization
\[
N_{\mathrm{CRC}}=\pi\left(\frac{r_{\mathrm{CRC}}}{\ell_\star}\right)^2.
\]
Work in the multiplicative screen-radius coordinate \(\rho=\ell/r_{\mathrm{CRC}}\), so the
horizon sits at \(\rho=1\) and the local cell at \(\rho_\star=\ell_\star/r_{\mathrm{CRC}}\).
Define the global repair cycle \(G_N\) at the fixed point as the scale-free readback transport
of the horizon record coordinate, with fixed-point closure (readback without deficit or slack)
meaning \(G_N(1)=\rho_\star\) exactly. The representation-to-spectrum theorem gives
\(m_{\rm rep}=24\), so the selected resonance branch decomposes as \(G_N=g_N^{24}\) with
\(g_N\) the positive root. Then
\[
G_N(\rho)
=
\left(\frac{N_{\mathrm{CRC}}}{\pi}\right)^{-1/2}\rho,
\qquad
g_N(\rho)
=
\left(\frac{N_{\mathrm{CRC}}}{\pi}\right)^{-1/48}\rho,
\qquad
|g_\star'|=\left(\frac{N_{\mathrm{CRC}}}{\pi}\right)^{-1/48}.
\]
More generally an \(m\)-tick decomposition gives
\(|g_\star'|=(N_{\mathrm{CRC}}/\pi)^{-1/(2m)}\); \(m_{\rm rep}=24\) gives the stated
exponent.
\end{lemma}

\begin{proof}
The D6 normalization gives
\(\rho_\star=\ell_\star/r_{\mathrm{CRC}}=(N_{\mathrm{CRC}}/\pi)^{-1/2}\). The closure
condition follows at the readback interface. The Observers synthesis proposes the formal readback
expression
\(F(N)=\operatorname{Cap}_{\mathrm{read}}(\operatorname{Obs}(\operatorname{nf}_{r,N}(\mathfrak
U_{r,N})))\) in the refinement limit. This expression does not construct the physical map or
establish CP-1 to CP-3. Under the lemma's assumptions, modeling
\(\operatorname{Cap}_{\mathrm{read}}\) by the D6 area
law at the effective delivery resolution \(\rho_{\mathrm{read}}\), the canonical
scale-covariant extension of \(N=\pi(r/\ell)^2\), gives the D6-consistent counting
identification. The corpus
defines \(\operatorname{Cap}_{\mathrm{read}}\) as reconstructed capacity; this certificate adds
only the D6-consistent counting rule. Under the stated assumptions, the readback reconstructs
\(F(N)=\pi/\rho_{\mathrm{read}}^2\), so the fixed-point equation
\(N_{\mathrm{CRC}}=F(N_{\mathrm{CRC}})=\pi/\rho_\star^2\) forces
\(\rho_{\mathrm{read}}=\rho_\star\) for positive roots: deficit
(\(\rho_{\mathrm{read}}>\rho_\star\)) gives \(F(N)<N\), slack
(\(\rho_{\mathrm{read}}<\rho_\star\)) gives \(F(N)>N\). This is exactly
\(G_N(1)=\rho_\star\). Homogeneity forces \(G_N(\rho)=Q\rho\) with \(Q>0\), hence
\(Q=(N_{\mathrm{CRC}}/\pi)^{-1/2}\). The homogeneous one-tick map with
\(g_N^{24}=G_N\) is the positive \(24\)th root, so its multiplier, and hence its absolute
derivative in this coordinate, is
\((N_{\mathrm{CRC}}/\pi)^{-1/(2\cdot24)}=(N_{\mathrm{CRC}}/\pi)^{-1/48}\).
\end{proof}

The emitted particle rows continue to use the \(\mathcal R_U\) certificate above.
Lemma~\ref{lem:global-repair-tick} closes the mathematical global tick side, conditional on its
physical-map premise, with the derived
representation-to-spectrum count \(m_{\rm rep}=24\). The proof derives the closure transport at
the interface level of the formal \(F\) under the D6 area-law counting model of
\(\operatorname{Cap}_{\mathrm{read}}\).

\begin{theorem}[Joint product fixed-point branch]\label{thm:joint-pn-product}
Let \(x=\log N\), let \(X=I_P\times I_x\), and equip \(X\) with
\[
d((P,x),(Q,y))=\max\{|P-Q|,\ |x-y|\}.
\]
On the product-separated source branch the joint map is
\[
\mathcal J(P,x)=(\Gamma(P),\widehat C(x)),
\qquad
\widehat C(x)=\log F(e^x),
\]
where \(\Gamma\) is the local pixel closure map and \(F\) is the global capacity readback map.
If
\[
|\Gamma(P)-\Gamma(Q)|\le q_P|P-Q|,
\qquad
|\widehat C(x)-\widehat C(y)|\le q_N|x-y|,
\qquad q_P,q_N<1,
\]
then \(\mathcal J\) has a unique fixed point \((P_\star,\log N_{\mathrm{CRC}})\), and every
iteration in \(X\) converges to it with contraction constant \(q=\max\{q_P,q_N\}\).
\end{theorem}

\begin{proof}
For any two points in \(X\),
\[
d(\mathcal J(P,x),\mathcal J(Q,y))
\le
\max\{q_P|P-Q|,\ q_N|x-y|\}
\le
\max\{q_P,q_N\}d((P,x),(Q,y)).
\]
Banach's theorem gives existence, uniqueness, and stability. A genuinely coupled source map
\(\mathcal J(P,x)=(\Gamma(P,x),\widehat C(P,x))\) requires derivative bounds \(a,b,c,d\) and
a weight \(r>0\) with
\[
\max\{a+b/r,\ d+rc\}<1
\]
in the weighted sup metric. Without that bound, the coupled branch has residual freedom and no
uniqueness theorem.
\end{proof}

\begin{theorem}[Electroweak tick-projection bridge]\label{thm:ew-tick-projection}
On the selected D10 source branch,
\[
\frac{v}{E_{\mathrm{cell}}}
=
\exp\!\left[-\frac{2\pi}{4\alpha_U(P)}\right].
\]
Together with the global repair tick
\[
|g_\star'|=\left(\frac{N}{\pi}\right)^{-1/48},
\]
define the electroweak projection exponent by
\[
\Pi_{\mathrm{EW}}(P,N)
:=
\frac{-\log(v/E_{\mathrm{cell}})}{-\log |g_\star'|}
=
\frac{24\pi}{\alpha_U(P)\log(N/\pi)}.
\]
Since \(0<|g_\star'|<1\), this exponent is unique. The bridge projection
\[
\Pi_{\mathrm{EW}}(P_\star,N_{\mathrm{EW}}(P_\star))=4P_\star
\]
is equivalent to the bridge residual
\[
\mathcal B_{\mathrm{EW}}(P,N)
:=
\alpha_U(P)\log(N/\pi)-\frac{6\pi}{P}
=0.
\]
Equivalently, the bridge coordinate is
\[
N_{\mathrm{EW}}(P)
=
\pi\exp\!\left[\frac{6\pi}{P\alpha_U(P)}\right].
\]
For the public endpoint branch,
\[
N_{\mathrm{EW}}(P_\star)
=
3.5323546226929906511187512962330547600462\times10^{122}.
\]
For the bridge-map coordinate, conventionally labeled
\[
N_{\mathrm{CRC}}^{\mathrm{EW}}
=
\pi\exp\!\left[\frac{6\pi}{P_\star\alpha_U(P_\star)}\right],
\]
one has \(\mathcal B_{\mathrm{EW}}(P_\star,N_{\mathrm{CRC}}^{\mathrm{EW}})=0\), and hence
\[
\frac{v}{E_{\mathrm{cell}}}
=
\left(\frac{N_{\mathrm{CRC}}^{\mathrm{EW}}}{\pi}\right)^{-P_\star/12}
=
|g_\star'|^{4P_\star}.
\]
The factor \(4\) is the electroweak transmutation multiplicity
\(\beta_{\mathrm{EW}}=N_c+1=4\). The factor \(48=2m_{\rm rep}\) uses the
product-adjoint count, hence \(12=48/4\). The notation
\(N_{\mathrm{CRC}}^{\mathrm{EW}}\) does not identify this mathematical coordinate with a cosmic
readback fixed point. That identification requires a constructed physical map \(F\) and premises
CP-1 to CP-3. The Planck-\(\Lambda\) central capacity
\(N_\Lambda\simeq3.313\times10^{122}\) is about \(6.6\) percent lower than the bridge coordinate.
\end{theorem}

\begin{proof}
The D10 relation gives
\[
-\log(v/E_{\mathrm{cell}})=\frac{2\pi}{4\alpha_U(P)}=\frac{\pi}{2\alpha_U(P)}.
\]
The global tick relation gives
\[
-\log|g_\star'|=\frac{1}{48}\log(N/\pi).
\]
Taking the ratio gives the displayed \(\Pi_{\mathrm{EW}}\). Setting
\(\Pi_{\mathrm{EW}}=4P\) and rearranging gives
\(\alpha_U(P)\log(N/\pi)=6\pi/P\), exactly
\(\mathcal B_{\mathrm{EW}}(P,N)=0\). Substitution gives the displayed
\(N_{\mathrm{EW}}(P)\) and the hierarchy identity.
\end{proof}

\begin{theorem}[Icosahedral screen sieve and gated capacity composition]\label{thm:screen-sieve-capacity-projection}
On the locally six-valent triangulated \(S^2\) branch, assume integer charges,
a feasible twelve-unit configuration, and an additive cost satisfying
\[
 h(0)=0,\quad h(1)>0,\quad h(k)\ge h(1)|k|,
\]
strictly for \(|k|\ge2\). Assume edge-center collar readout, inverse-port
antipodal pairing, and the D-optimal vector/quadrupole selector of
Ref.~\cite{ophmicrophysics}. Then the screen has twelve unit curvature ports on
the regular icosahedral orbit, and an equal-weight invariant load \(X\) is read
locally as \(X/12\). For
\(X=\log(N/\pi)\), local cell entropy \(P/4\), and electroweak multiplicity
\(\beta_{\mathrm{EW}}=4\), define the screen composition coordinate
\[
\Gamma_{\mathrm{scr}}
=
\beta_{\mathrm{EW}}\frac{P}{4}\frac{1}{12}\log(N/\pi)
=
\frac{P}{12}\log(N/\pi).
\]
If, in addition, \textsc{HIERARCHY-SCREEN-READOUT} identifies
\(\Gamma_{\mathrm{scr}}=\log(E_{\mathrm{cell}}/v)\) and matches it to the
source transmutation exponent \(\pi/(2\alpha_U(P))\), then
\[
\log(E_{\mathrm{cell}}/v)
=\frac{P}{12}\log(N/\pi),
\qquad
\alpha_U(P)^{-1}=\frac{P}{6\pi}\log(N/\pi),
\]
which is equivalent to the electroweak bridge residual
\(\mathcal B_{\mathrm{EW}}(P,N)=0\).
\end{theorem}

\begin{proof}
Euler and triangular incidence give
\[
\sum_v q_v=6V-2E=12.
\]
The cost inequality gives
\(\sum_vh(q_v)\ge h(1)\sum_v|q_v|\ge12h(1)\), with equality only for twelve
positive unit charges. The D-optimal equality conditions are
\(F_1=2I_3\) and \(F_2=(4/5)I_5\). They make the six antipodal axes the unique
real equiangular tight frame in three dimensions, hence the regular
icosahedral axes. Edge-center collars expose the twelve vertices as central
ports. Equal-weight additivity gives \(X/12\). Multiplying by \(P/4\) and
\(\beta_{\mathrm{EW}}=4\) gives
\(\Gamma_{\mathrm{scr}}=(P/12)\log(N/\pi)\). The additional identification
receipt and the source transmutation equation give the remaining displays.
\end{proof}

The unit-splitting cost, D-optimal selector, and
\textsc{HIERARCHY-SCREEN-READOUT} are physical branch premises. Their source
production is work in progress.

\begin{theorem}[EW bridge-map fixed point]\label{thm:ew-exact-capacity}
Define the log-capacity map
\[
\mathcal C_{\mathrm{EW}}(P,x)
=
(1-\lambda)x+\lambda\frac{6\pi}{P\alpha_U(P)},
\qquad 0<\lambda\le1.
\]
For fixed \(P=P_\star\) and \(\lambda=1/2\), this is a contraction with
Lipschitz constant \(1/2\). Its unique fixed point is
\[
x_\star=\frac{6\pi}{P_\star\alpha_U(P_\star)},\qquad
N_{\mathrm{CRC}}^{\mathrm{EW}}=\pi e^{x_\star}.
\]
Therefore
\[
\mathcal B_{\mathrm{EW}}(P_\star,N_{\mathrm{CRC}}^{\mathrm{EW}})=0.
\]
This theorem concerns the deliberately defined map \(\mathcal C_{\mathrm{EW}}\), not the physical
cosmic readback map \(F\). Identifying the output with cosmic capacity is conditional on CP-1 to
CP-3. The Planck-\(\Lambda\) central value \(3.313\times10^{122}\) differs by about \(6.6\)
percent.
\end{theorem}

\begin{proof}
The map has derivative \(1-\lambda\) in \(x\), hence is a contraction for
\(0<\lambda\le1\). Solving \(x=\mathcal C_{\mathrm{EW}}(P_\star,x)\) gives the
displayed \(x_\star\). Substitution into \(\mathcal B_{\mathrm{EW}}\) gives
zero.
\end{proof}

\begin{theorem}[Conditional RG/Higgs naturality identity]\label{thm:rg-higgs-naturality}
On a declared hierarchy branch satisfying the selected-surface premises, let \(Q_s\) be the source hierarchy quotient and
\(Q_H\) the Higgs/electroweak quotient. Define
\[
\rho_{sH}([x]_s)=
\left[
P(x),N(x),
\Theta(P(x),N(x)),
\Pi_{HT}F_{D11}F_{D10}(P(x),N(x)),0
\right]_H,
\]
with
\[
\Theta(P,N)=
\left(\frac{N}{\pi}\right)^{-P/12}
=|g_\star'|^{4P}.
\]
For the corresponding normal-form maps and obstruction maps,
\[
\rho_{sH}n_s=n_H\rho_{sH},
\qquad
\chi_{sH}h_s=h_H\rho_{sH}.
\]
Consequently,
\[
\varepsilon_H
=
\max\{\varepsilon^n_{sH},\varepsilon^h_{sH}\}
=0,
\qquad
\varepsilon_H\in[0,0].
\]
\end{theorem}

\begin{proof}
The two routes through the source and Higgs quotients preserve the same tuple
\[
(P,N,\Theta(P,N),\Pi_{HT}F_{D11}F_{D10}(P,N)).
\]
The equality with \(|g_\star'|^{4P}\) follows from
\(|g_\star'|=(N/\pi)^{-1/48}\). A hidden scalar relevant deformation is not an independent
physical input on the selected branch: if absent from the retained finite constraint family, it
is off-branch by refinement stability; if retained, its value is fixed by the branch equations
and belongs to the displayed Higgs/top coordinate. Both controlled squares commute exactly, so
the product defect is zero.
\end{proof}

The conditional RG/Higgs naturality map does not itself use the measured weak scale. Measured \(W/Z\),
Higgs/top, \(G\), Planck-area, and \(\Lambda\) data are excluded by its declared input contract. The local/global package
records a conditional global repair tick, joint product branch, electroweak projection bridge,
mathematical EW bridge-map fixed point, finite readback-resolution certificate, and RG naturality
square. Conditional on the declared branch, this certifies the dimensionless electroweak hierarchy/naturality
identities; it does not construct the physical cosmic map \(F\) or discharge CP-1 to CP-3. The strict
source-root certificate is open, and the result does not independently attach
\(E_\star\) to a physical GeV scale. The finite readback certificate defines
\[
F_r(N)=\operatorname{Cap}_{\rm read}
\left(\operatorname{Obs}(\operatorname{nf}_{r,N}(\mathfrak U_{r,N}))\right),
\qquad
\rho_{\rm read}(r,N)=\sqrt{\pi/F_r(N)},
\]
and proves one selected effective delivery resolution on the finite branch, with
\(\rho_{\rm read}(r,N_{\rm CRC})\to(N_{\rm CRC}/\pi)^{-1/2}\) in the
positive-root refinement limit. The product-adjoint theorem gives
\(m_{\rm rep}=2(8+3+1)=24\). The conditional screen separately has twelve ports and an
oriented 24-slot register; their equal counts lack a physical identification receipt. The
mathematical local/global hierarchy-resonance package closes on the selected branch. This finite map
is not a construction of the physical cosmic readback map \(F\).
The promoted hierarchy row does not use the
local/global resonance as an input.

\subsection{Higgs/top critical stage}

The Higgs/top critical stage is a conditional quantitative stage built on top of the electroweak gauge core.
It is described by a downstream split law together with a Jacobian readout map
on the declared D10/D11 running, matching, and threshold surface.  Write the
coupling asymmetry as \(\rho_{\mathrm{EW}}\), distinct from
\(b_{\mathrm{tr}}=N_c+1=4\). The candidate D10 tuple
\[
(\eta_{\mathrm{source}},\rho_{\mathrm{EW}},\lambda_{EW},\tau_{2,\mathrm{tree}}^{\mathrm{exact}},\delta n_{\mathrm{tree}}^{\mathrm{exact}})
\]
emits the shared scalar
\[
\rho_{HT}=\log\!\bigl(1+\tau_{2,\mathrm{tree}}^{\mathrm{exact}}\bigr)
\]
and the declared residual formulas
\[
R_T=
-\tau_{2,\mathrm{tree}}^{\mathrm{exact}}\eta_{\mathrm{source}}^2
+\Bigl(1+\frac{\rho_{\mathrm{EW}}}{28}\Bigr)\eta_{\mathrm{source}}^6
+\frac{\eta_{\mathrm{source}}^8}{14}
+\frac{\eta_{\mathrm{source}}^9}{27},
\]
\[
R_H=
\eta_{\mathrm{source}}^5
-\frac{3}{25}\eta_{\mathrm{source}}^6
+\frac{\lambda_{EW}\eta_{\mathrm{source}}^6}{18}
+\frac{\eta_{\mathrm{source}}^8}{2\rho_{\mathrm{EW}}}.
\]
The split coordinates are
\[
\pi_y=
\frac{\eta_{\mathrm{source}}+\left(\frac32+\frac{\rho_{\mathrm{EW}}}4\right)\rho_{HT}+R_T}{\sqrt{\pi}},
\qquad
\pi_\lambda=
\frac{\eta_{\mathrm{source}}-\left(\frac43-\frac{\rho_{\mathrm{EW}}}{54}\right)\rho_{HT}+R_H}{\sqrt{\pi}},
\]
so the declared D11 Jacobian reads out
\[
\delta y_t(\mu_t)=\pi_y\,y_t^{\mathrm{core}}(\mu_t),
\qquad
\delta\lambda(\mu_t)=-\frac{16}{9}\pi_\lambda\,\lambda^{\mathrm{core}}(\mu_t).
\]
These equations are promoted from a declared formula surface only if a finite
D11 split-character certificate passes DS1--DS5: one strict source branch and
frozen pre-split carrier; an exact two-coordinate quotient response; explicit
characters deriving \(\rho_{HT},R_T,R_H,\pi_y,\pi_\lambda\) and \(-16/9\);
oriented mismatch descent with canonical normalization and positivity; and an
exhaustive deformation quotient proving rigidity or a target-free positive-gap
selector.  The repository evaluates the displayed formulas without emitting
that certificate.
This gives
\[
m_H = 125.1995304097179~\mathrm{GeV},
\qquad
m_t^{D11} = 172.3523553288312~\mathrm{GeV}.
\]
This is a coordinate on the declared D10/D11 Jacobian surface, obtained by
back-solving from the measured pair through the synchronization-scale scan, so
it is an exact interpolation of the PDG pair, a target-anchored fit; it
carries no certified Higgs or top complex pole. The
target-free headline is the double-criticality family evaluated below. The same surface emits a companion top coordinate.
The selected-class quark support witness carries a
target-anchored running-top audit row using the PDG 2025 cross-section entry; it is not a separate
public source-only top prediction. The bridge to the auxiliary direct-top PDG row is closed as a
corpus-limited codomain no-go; the auxiliary row is compare-only.
The one-scalar companion seed
\[
\sigma_{D11,\mathrm{HT}}=\frac{\alpha_U\cos(2\theta_{W0})}{\sqrt{\pi}}
\]
is on disk only as the fixed-ray companion branch with \(\pi_y=\pi_\lambda=\sigma_{D11,\mathrm{HT}}\).
Separately, the same D11 Jacobian admits a compare-only inverse slice that hits the canonical
Higgs/top reference pair exactly: it is an exact interpolation of the PDG pair, and the landing
is the defining condition of the slice rather than a derived result. That exact sidecar is
useful for bookkeeping and does not define
the forward branch used for the public rows.
The Higgs boson belongs naturally in the boson discussion, while the top quark will be revisited
in the quark-family chapter as the third-generation up-type quark. The top coordinate reported
here is the conditional D11 split coordinate on this surface, not a separate public top-mass prediction
row. The selected-class quark support wrapper carries only a target-anchored audit witness under
the strict public-output policy.

The local implication from the candidate D10 tuple through the displayed D11
formulas is executable and single-valued, and this arithmetic grade is
certified from raw interval data: the repository publishes the raw interval
input box for the eleven declared branch inputs (the D10 tuple and the D11
core/Jacobian constants, each with units or dimensionless normalization and
provenance tags), the outward-rounded interval extension of every displayed
node, the Jacobian interval enclosure over the full box, and a non-singular
diagonal readout block with determinant bounded away from zero
(\texttt{R\_HT\_interval\_input\_box\_log.json}, verifier
\texttt{verify\_ht\_interval\_box.py}); the certified statement is scoped to
the declared surface and makes no criticality-system existence or uniqueness
claim. Strict Higgs promotion is
blocked by the source-root, independent-scale, QT1--QT5, RG/scheme,
split-rigidity, top/threshold, complex-pole, uncertainty, and target-independent
provenance gates.  The exact inverse slice is a compare-only validation
surface.

\begin{table}[h]
\centering
\scriptsize
\begin{tabular}{@{}>{\raggedright\arraybackslash}p{0.13\textwidth}>{\raggedright\arraybackslash}p{0.21\textwidth}>{\raggedright\arraybackslash}p{0.34\textwidth}>{\raggedright\arraybackslash}p{0.21\textwidth}@{}}
\toprule
Particle & Stage & Claim tier & Reported value \\
\midrule
\(W\) boson & D10 electroweak branch & no source-only physical mass; audit coordinates quarantined & withheld \\
\(Z\) boson & D10 electroweak branch & no source-only physical mass; audit coordinates quarantined & withheld \\
Higgs boson & Higgs/top critical stage & no source-only physical mass; conditional coordinate not promoted & withheld \\
Top quark & Higgs/top critical stage & companion coordinate on declared surface; direct-top auxiliary codomain no-go & not a separate public top-mass prediction row \\
\bottomrule
\end{tabular}
\end{table}

\subsection{Source-closure and physical-pole envelope}\label{wzh-source-closure-envelope}

The preceding formulas end at running or declared-surface coordinates.  A
physical \(W/Z/H\) theorem requires the following additional composition.

\begin{proposition}[Conditional physical clock attachment]
If the same strict source branch emits the dimensionless energy
\(\varepsilon_{\rm clk}>0\) of a declared clock transition whose physical
frequency \(\nu_{\rm clk}\) defines the operational unit, then
\[
E_\star=\frac{h\nu_{\rm clk}}{\varepsilon_{\rm clk}}.
\]
Equivalently, \(\gamma_\star=\ell_\star\nu_{\rm clk}/c\) gives
\(E_\star=\hbar\nu_{\rm clk}/\gamma_\star\).  An operator-norm perturbation
of at most \(\epsilon\) changes an isolated selected clock gap by at most
\(2\epsilon\).  This is a conversion and stability theorem; the required
source clock packet is absent from the source.
\end{proposition}

The atomic component of that packet has an exact scalar reduction.  Let
\(P=P_3+P_4\) project onto the isolated zero-field \(6S_{1/2}\), \(I=7/2\)
cesium ground manifold, \(\dim V_3=7\), \(\dim V_4=9\), and let \(Q=1-P\).
On every real interval where \(Q(\widehat H-z)Q\) is invertible, define the
Feshbach--Schur operator
\(\mathcal F_P(z)=P(\widehat H-z)P-P\widehat HQ\,[Q(\widehat H-z)Q]^{-1}Q\widehat HP\).

\begin{proposition}[Cesium channel scalarization and monotone clock roots]
\(z\in\operatorname{spec}\widehat H\) exactly when
\(0\in\operatorname{spec}\mathcal F_P(z)\), with matching multiplicities.
Rotational invariance and the multiplicity-free decomposition
\(V_3\oplus V_4\) force \(\mathcal F_P(z)=f_3(z)P_3+f_4(z)P_4\), and
\(\mathrm d\mathcal F_P/\mathrm dz\preceq-P\) gives \(f_F'(z)\le-1\).  One
sign bracket per channel therefore isolates exactly one level, and the clock
normal form is
\(H^{\mathrm{eff}}_{\mathrm{clock}}=E_0P+\widehat a_{\mathrm{Cs}}\,\mathbf I\cdot\mathbf J\)
with \(E_0=(9E_4+7E_3)/16\) and
\(\varepsilon_{\mathrm{Cs}}=E_4-E_3=4\widehat a_{\mathrm{Cs}}\).
\end{proposition}

The public atomic packet may therefore export two monotone scalar interval
evaluators with resolvent certificates and sign brackets in place of a
correlated 55-electron matrix.  The open clock parents are the atomic-scheme
electromagnetic convention, the absolute electron ratio \(m_ec^2/E_\star\),
the cesium nuclear mass/spin/current/Compton packet, the two scalar
evaluators, summable refinement tails, and the no-target provenance freeze.
A non-entailment theorem fixes the direction of effort: two source extensions
that agree on every named object and differ in one unfixed parent
produce different unique gaps, so no unique \(\varepsilon_{\mathrm{Cs}}\) is a
theorem emitted by the source, and further evaluation of the existing
formulas cannot close the clock lane.

\begin{proposition}[Operational scale metrology]
In SI units
\(E/\mathrm{GeV}=[E/(h\nu_{\mathrm{Cs}})]\,[h\nu_{\mathrm{Cs}}/\mathrm{GeV}]\),
and the second factor is an exact unit conversion.  A source-only numerical
GeV mass is therefore equivalent to a source-only dimensionless operational
clock ratio.  Replacing cesium by another clock species adds the burden of
predicting that clock's frequency ratio to the defining cesium transition.
\end{proposition}

The stored clock-gap candidate reproduces the displayed gravitational
coupling to a relative agreement of order \(10^{-49}\) through the
clock--gravity identity, so the decimal is a calibration checksum excluded from the prediction ledger
(\texttt{clock\_source\_energy\_closure\_audit.json}).  The selection
mechanism for source laws themselves is closed separately by a source-action
rigidity theorem: on a finite quotient with a faithful reference measure, an
affinely independent feature basis, and a source-emitted moment vector, the
maximum-entropy law is unique, every feasible competitor pays a positive
relative-entropy gap with the Pinsker lower bound, and the surviving action
freedom consists of additive constants, exact feature redundancies, and
BRST-exact terms
(\texttt{qme\_source\_action\_rigidity\_mechanism.json}).  The physical
Standard-Model operator basis and its moment vector
\(c_r(P_\star,N_\star)\) are the open inputs of that mechanism.

\begin{proposition}[Unique frozen RG and threshold transport]
Let the canonically normalized coefficient vector \(X(\mu)\) obey a locally
Lipschitz beta system on each interval of a frozen finite threshold list.
Assume deterministic matching maps, fixed loop order, field content,
renormalization scheme, threshold locations, and decoupling order.  Then one
source initial condition determines one low-energy coefficient vector.  If
\(L_j\) bounds the beta-system Lipschitz constant and \(K_j\) bounds threshold
map \(j\), perturbations obey the corresponding product of
\(e^{L_j\Delta t_j}\) and \(K_j\), plus the declared truncation and matching
remainders.
\end{proposition}

\begin{proof}
Picard--Lindel\"of gives existence and uniqueness between thresholds.
Gr\"onwall bounds perturbations on each interval; applying each deterministic
Lipschitz threshold map and inducting through the ordered list gives the
stated composition.
\end{proof}

The repository's printed running/matching packet is a declared-convention
contract, not this completed source theorem: its beta provenance, threshold
origins, matching interval composition, and truncation enclosure are
promotion gates.

\begin{theorem}[Conditional complex-pole promotion]
For \(B=W,Z,H\), let \(\Gamma_B^{\rm phys}(s,\xi)\) be the BRST-complete
physical inverse two-point block, including every declared mixing field after
the Ward/Slavnov--Taylor projection, analytically continued to a declared
Riemann sheet, and set
\[
D_B(s,\xi)=\det\Gamma_B^{\rm phys}(s,\xi).
\]
Assume a reference determinant has exactly one simple physical zero inside a
certified contour, no zero lies on the contour, and the full determinant obeys
\(|D_B-D_{B,0}|<|D_{B,0}|\) there.  In the neutral sector also assume the
Ward-protected photon line has been separated from the massive eigenvalue.
Then the enclosed pole is unique and stable.  With the exact convention
\[
s_B=\left(M_B-\frac{i}{2}\Gamma_B\right)^2,
\qquad M_B>0,\quad\Gamma_B\ge0,
\]
it determines one mass and width.  If a Nielsen identity has the form
\(\partial_\xi D_B=C_B^{(\xi)}D_B\), the simple pole is gauge-parameter independent.
\end{theorem}

\begin{proof}
Rouch\'e's theorem preserves the number of enclosed zeros under the certified
perturbation.  Simplicity gives the implicit-function stability bound and a
nonzero residue.  Differentiating \(D_B(s_B(\xi),\xi)=0\) and using the Nielsen
identity gives \(ds_B/d\xi=0\).  An invertible analytic field redefinition
multiplies \(D_B\) by a nonzero analytic factor and therefore preserves the
pole set.
\end{proof}

No source-emitted \(W/Z/H\) self-energy kernels, analytic-sheet receipt,
contour enclosures, widths, residues, or pole-convention uncertainty map are
present on the displayed mass-chart surface. Schema fields named
\texttt{MW\_pole}, \texttt{MZ\_pole}, or \texttt{mt\_pole} are labels only and do not
supply this theorem.

\begin{theorem}[Rigidity, dependency, and W/Z/H composition criterion]
Let \(\mathcal F_{\mathrm{EW}}\) be the fully enumerated class of source maps
that preserve the declared locality, symmetry, dimensional, refinement, and
provenance constraints, modulo gauge and trivial reparameterization.  Suppose
either the pole readout is constant on this class or a target-independent
selector has one minimizer with a positive winner gap.  Suppose further that
a hash-bound acyclic dependency manifest, committed before evaluation, records
the source artifacts, code, selectors, conventions, uncertainty model, and
outputs, and that measured \(W/Z/H\) data and calibrated proxies are absent
from every predicted ancestor set.

If, in addition, a unique source root and independently source-closed physical
\(E_\star\) emit the canonical D10/D11 coefficients, the hypotheses of
Definition~\ref{def:particle-d10-qt}, DS1--DS5, the frozen RG proposition, and the
complex-pole theorem hold, then OPH determines one source-separated pole
triple \((s_W,s_Z,s_H)\) with a certified uncertainty enclosure.
\end{theorem}

\begin{proof}
Rigidity or the positive-gap selector removes alternative source maps.  A
deterministic leaf depends only on its ancestors, so the committed DAG gives
formal source separation.  The unique root and scale emit one coefficient
packet, frozen RG and matching transport it uniquely, and the pole theorem
gives one stable pole per particle.  Composition proves the claim.
\end{proof}

This criterion is not satisfied by the displayed numbers. QT1--QT5 are
certificate assumptions, the public and source-audit pixel lanes differ, the
physical scale and concrete RG receipt are open, DS1--DS5 are not emitted, and
the complex-pole and target-independent precommitment gates are absent.

The missing object is a source-emitted certificate.
The structural theory admits target-free analytic counterfamilies
\(\tau_2=-c\eta^2\), \(\delta n=d(1-\rho_{\rm EW})\eta^2\) and
\((\pi_y,\pi_\lambda)\mapsto(\pi_y+a\eta^N,
\pi_\lambda+b\eta^N)\), all on the same open physical domain but with
different mass-chart outputs. Likewise, a fixed running mass is compatible
with both $s-m_R^2$ and $s-m_R^2-\epsilon$, which have different poles.
These countermodels prove that two-channel exhaustion and analyticity do not
emit the D10 character, D11 split, or physical kernel. Nor is a formal path
realization sufficient: any finite polynomial can be encoded by assigning one
weighted path to each monomial. A non-vacuous certificate must independently
fix its transitions, path measure, signs, quotient action, exhaustive census,
rigidity gap, and no-target DAG.

Once genuine source packets are supplied, no stochastic simulation is needed.
Summable clock-Hamiltonian refinements give a unique limiting gap; frozen
Lipschitz RG segments compose uniquely; Rouch\'e, Nielsen, and analytic
conjugacy isolate gauge-independent poles. If a determinant defect is bounded
by \(\epsilon\), \(a=D'(s_0)\ne0\), and \(|D''|\le K\), every $r$ with
\(\epsilon<|a|r-Kr^2/2\) encloses the unique displaced pole and propagates to
mass/width bounds through the selected square root. The four absent source
objects are therefore the factorized clock packet, independently weighted
D10 carrier, D11 split-character carrier, and BRST-complete pole-kernel packet.
A runtime DAG proves computational separation only; a blind claim also requires a disclosed,
target-independent frozen specification.

The electroweak chapter therefore presents four mathematically distinct
surfaces: the selected-carrier chart, the candidate value law, the
conditional QT1--QT5 selection theorem, and the
measured-reference inverse surface used for diagnostic \(W/Z\) values.  A physical
pole surface requires the additional source-root, scale, RG/scheme,
two-point-kernel, rigidity, uncertainty, and provenance receipts.

\subsection{Repair-tuple selection and the color amplitude/loop split}\label{ew-repair-tuple-selection}

The candidate D10 tuple sits inside a two-parameter freedom that the emitted
corpus does not close.  With
\(\rho_{\mathrm{EW}}=(\alpha_2-\alpha_Y)/(\alpha_2+\alpha_Y)\), the coherent
quadratic family is
\[
\tau_{2,\mathrm{tree}}^{\mathrm{exact}}=-c\,\eta_{\mathrm{source}}^2,
\qquad
\delta n_{\mathrm{tree}}^{\mathrm{exact}}=d\,(1-\rho_{\mathrm{EW}})\,\eta_{\mathrm{source}}^2,
\]
and an open neighborhood of distinct \((c,d)\) values preserves the positive
mass-chart domain and defines consistent D10 repairs. The charged leg is the
\(\SU(2)_L\) coupling correction \(\delta\alpha_2=\alpha_2\tau_2\); the neutral leg is the
hypercharge screening correction carried by \(\delta n\).

The compact-gauge reconstruction fixes the color count but does not fix the
channel assignment or normalization that would close this freedom.
Doplicher--Roberts/Tannaka reconstruction returns the color triplet sector with statistical
dimension \(d(\rho_3)=N_c=3\), and its standard conjugate intertwiner satisfies
\(R^*R=d(\rho_3)=N_c\), so \(R\) has norm \(\sqrt{N_c}\) while
\(R/\sqrt{N_c}\) is the isometry; a closed color loop carries the full
dimension \(N_c\). The two repair legs sit at different levels. The charged leg is the broken,
mass-generating \(\SU(2)_L\) channel, driven by the color-singlet condensate amplitude whose
large-\(N_c\) scaling is the decay-constant scaling \(\sqrt{N_c}\) assigned to one standard
conjugate intertwiner. The
neutral leg would be the unbroken hypercharge screening channel, a
vacuum-polarization loop that closes the color line and carries raw trace
weight \(N_c\).  Under these extra hypotheses the alternative quadratic model
has
\[
c=\frac{\sqrt{N_c}}{2}=\frac{\sqrt3}{2},
\qquad
d=\frac{N_c}{2}=\frac32,
\]
where the chart coefficient \(d=N_c/2\) represents the raw neutral trace
because
\((\alpha_2+\alpha_Y)(1-\rho_{\mathrm{EW}})=2\alpha_Y\).

The resulting comparison is a target-informed, fixed-slice diagnostic of that
alternative model. Profiling the charged coefficient against
the D10 running-tree \(W\) coordinate does not use the D11 Jacobian core, but
it does inherit the declared \(\alpha_2\), \(v\), running, and scheme surface. It selects
\(c=0.8670\), while profiling the neutral coefficient returns \(d=1.461\). Both values are
target-conditioned fit outputs. Their proximity to \(\sqrt3/2\) and \(3/2\) is an internal
calibration observation, not physical \(W/Z\) evidence, and no experimental confidence band is
assigned across the noncommensurate chart and mass conventions. That slice is not a test of the complete candidate
value law, which changes both coordinates and includes higher path characters;
it therefore neither excludes that law nor proves or excludes a separate
running-tree companion.

The color-balanced rule does not equal the complete candidate value law in
Theorem~\ref{thm:particle-conditional-d10-qt}; it gives a different \(W\)
coordinate.  A proof of its amplitude and trace hypotheses would establish
that alternative quadratic model, not QT1--QT5.  The complete value-law
promotion task is the finite quotient-path campaign: construct the two-channel
quotient, enumerate the charged and neutral path lists, verify the
\((1,2,1)\) and \((1,4,2)\) incidences with their \(1/3\) color measure and
\(1/6\) central trace, compute the fibre Gram form and residual pairing,
enumerate admissible deformations, and prove a positive MAR gap.  Measured
\(W/Z\) agreement cannot replace that certificate.  The first implementation
step is to freeze the actual pre-repair carrier independently of these desired
outputs as an observer-like self-reading patch with bounded local state,
typed ports, readback records, and admissible feedback/repair moves.  That
carrier may contain neither the target masses nor the desired
incidence and central-trace coefficients.  An exact path/orbit generator can
then be treated as untrusted and checked by a smaller independent verifier;
two canonicalizers should agree on orbit hashes, and the final formal checker
must recompute path legality, exhaustion, exact response rank and nullspace,
central action, oriented mismatch first variation, fibre Gram form, and the
deformation census.  Agreement between two canonicalizers is regression
evidence rather than a completeness proof.  The order of
\(\mathbb Z_6\) supplies the averaging coefficient \(1/6\), not automatically
a normalized trace \(1/6\): the checker must separate the trivial physical
matter-kernel action from the declared D10 center-label/transport trace space
and derive the trace, source amplitude, and subtraction sign in the latter.

\subsection{Selector rigidity and the discrete repair-law boundary}\label{ew-selector-rigidity}

The candidate D10 selector closes the continuous \((c,d)\) freedom of the
preceding subsection.  Writing \(x=\tau_{2,\mathrm{tree}}^{\mathrm{exact}}\),
\(y=\delta n_{\mathrm{tree}}^{\mathrm{exact}}\), and
\(\kappa=(\alpha_Y+\alpha_2)/\alpha_Y\), the selector
\[
J_{10}(x,y)=\frac{x^2\left[1+(2\eta_{\mathrm{source}}+x)^2+4x^2\right]}{1+4x^2}
+\frac{\kappa^2}{4}(1+4x^2)y^2
\]
admits the exact factorization
\[
J_{10}-\frac34x^2-\frac{\kappa^2}{4}y^2
=\frac{x^2}{4(1+4x^2)}
\left[8(x+\eta_{\mathrm{source}})^2+8\eta_{\mathrm{source}}^2+1
+4\kappa^2(1+4x^2)y^2\right],
\]
so \(J_{10}\ge\frac34x^2+\frac{\kappa^2}{4}y^2\) with equality exactly at the
origin, and \(x^2+y^2\ge r^2\) forces
\(J_{10}\ge\min\{3/4,\kappa^2/4\}r^2\).  The selector picks the zero
deformation with a quantitative gap. The candidate value-law tuple has
\(J_{10}\approx3.12\times10^{-7}>0\), so it is incompatible with the selector,
and the residual electroweak freedom is
one discrete choice between two incompatible laws rather than a continuous
chart family.

On the strict source-audit branch the two branch points give the tree/chart
coordinates
\[
\begin{aligned}
\text{zero-selector law:}\quad
&M_W/E_\star=6.579631\times10^{-18}, &
M_Z/E_\star&=7.463335\times10^{-18},\\
\text{carrier value law:}\quad
&M_W/E_\star=6.578870\times10^{-18}, &
M_Z/E_\star&=7.463750\times10^{-18},
\end{aligned}
\]
with unclosed-clock displays \((80.3301,\,91.1191)\) and
\((80.3208,\,91.1242)\)~GeV and discrete ambiguity widths of \(9.3\)~MeV on
\(M_W\) and \(5.1\)~MeV on \(M_Z\)
(\texttt{wzh\_residual\_elimination\_boundary.json}). These are running/chart
coordinates. The PDG mass targets use a mass-dependent-width Breit--Wigner convention, and
complex-pole masses use another convention. The coordinates therefore have no valid physical
pull or near-hit interpretation; the complex-pole convention diagnostic (\(W\) at \(0.5\)
propagated experimental standard deviations from the converted pole, \(Z\) at about \(17\))
carries no theory covariance. Numerical \(W/Z\) proximity cannot decide between the laws;
the decision object is a source-law
selection principle, and the quantitative mechanism of the source-action
rigidity theorem applies once the electroweak feature basis and moment vector
are emitted.

The two-loop audit returned \((79.115335,89.802735)~\mathrm{GeV}\) in its designated primary
cell. It combines an MSSM one-loop baseline with an SM two-loop-minus-one-loop increment, so it is
an inconsistent MSSM-1L+SM-2L hybrid prescription test. The pole audit returned approximately
\((79.53284,89.71232)~\mathrm{GeV}\) in the primary cell of a partial
PRTS/Feynman-gauge prescription. Its definition of \(v\), tadpole treatment, field-content and
threshold matching, scale choice, and higher-order corrections are open. The outputs establish
neither a unique scheme conversion nor a unique \(1\)--\(2\%\) defect, and they do not exhaust
the physical prescription family.

A conditional finite C10 root/Cartan carrier emits the value-law coefficients
exactly: the transitive \(C_3\) color action has invariant measure \(1/3\),
the regular \(\mathbb Z_6\) register has rank-one trace \(1/6\), the
transitive four-slot register has slot measure \(1/4\), the frozen
independent-product source law gives
\(\lambda_{\mathrm{EW}}=\eta_{\mathrm{source}}^2/(4\rho_{\mathrm{EW}})\), and
the depth-two path census reproduces the candidate quadratic response
polynomials.  A companion C11 \(\SU(3)\)-Casimir carrier emits the declared
D11 normalizations, \(\kappa_\lambda=C_F^2=16/9\), the response coefficients
\(3/2+\rho_{\mathrm{EW}}/4\) and \(4/3-\rho_{\mathrm{EW}}/54\), and the
logarithmic character \(\rho_{HT}=\log(1+\tau_2)\) from multiplicativity with
unit derivative. Both carriers are target-exposed candidates without a
source-independent QT1--QT5 selection or census certificate.

The declared D11 Jacobian entries follow exactly from the core,
\(\partial m_t/\partial y_t=m_{t,0}/y_0\) and
\(\partial m_H/\partial\lambda=m_{H,0}/(2\lambda_0)\), and the exact readout
is \(m_t=m_{t,0}(1+r_y)\) and \(m_H=m_{H,0}\sqrt{1-r_\lambda}\) with the
exact linearization error \(m_{H,0}r^2/(2(1+\sqrt{1-r})^2)\).  The
synchronized D11 core is target-conditioned: its synchronization scale
minimizes an objective containing the Higgs and top comparison values.
Without that switch, the literal candidate one-loop source equations give the target-free tree coordinates
\(m_t^{\overline{\mathrm{MS}}}/E_\star=1.267758\times10^{-17}\) and
\(m_H^{\mathrm{tree}}/E_\star=9.427653\times10^{-18}\), with unclosed-clock
displays \(154.78\)~GeV and \(115.10\)~GeV and a QCD-converted top display of
\(164.13\)~GeV.  These are one-loop tree coordinates on the declared surface;
they are not physical poles, and the pole gates of
\S\ref{wzh-source-closure-envelope} keep their status.

The candidate one-loop coordinates form a zero-continuous-%
parameter family.  The literal core derives its top boundary from the gauge
sector through the double-criticality condition
\(\lambda(\mu_b)=0\), \(\beta_\lambda(\mu_b)=0\), which fixes
\(y_t(\mu_b)=[(2g_2^4+(g_2^2+g_Y^2)^2)/16]^{1/4}\). The candidate branch
imposes this at the gauge-unification scale \(\mu_U\), while the electroweak
transmutation in the same model anchors \(v\) at the pixel energy
\(E_{\mathrm{cell}}=E_\star/\sqrt P\), so the candidate model carries two
different high-scale anchors.  Evaluating the same law at the named source
scales gives, at one loop, \((m_t,m_H)=(164.1,\,115.1)\)~GeV at \(\mu_U\),
\((170.4,\,127.8)\) at \(E_{\mathrm{cell}}\), and \((170.7,\,128.3)\) at
\(E_\star\); a benchmark-validated two-loop upgrade shifts these to
\((169.4,\,119.4)\), \((175.7,\,131.8)\), and \((175.9,\,132.3)\).  The
family brackets the measured pair in both channels at both loop orders, so
the candidate deficit decomposes into the boundary-scale choice plus loop
truncation, with no continuous freedom anywhere.  Along the fit-free curve
the Higgs coordinate at the measured top mass is \(125.7\)~GeV at two loops,
within \(0.5\) percent of measurement and inside the declared tree-to-pole
matching band; the implied boundary scale is \(4.8\times10^{17}\)~GeV,
between \(\mu_U\) and \(E_{\mathrm{cell}}\). No boundary-scale selection theorem is supplied
(\texttt{d11\_criticality\_boundary\_scan.json},
\texttt{d11\_criticality\_comparison.json}); numerical agreement cannot
select the scale, and the declared-surface values obtained by back-solving
from the measured pair stay classified as target-anchored fits.

The selection question admits a near closure.  The only scale-selecting
condition available inside the flow itself, the triple-criticality root
\(\mathrm d\beta_\lambda/\mathrm d\ln\mu=0\) at the critical point, has no
solution: along the family the flow derivative is dominated by
\(-24y_t^3\beta_{y_t}>0\) and stays bounded away from zero across the whole
window, so every family point is a clean \(\lambda\) minimum and the
boundary scale must be selected by the source structure.  That structure
supplies a variational selection principle with three premises: (AR1) the
criticality boundary is a record that reconciles the model's two anchor
records, the gauge-unification record at \(\mu_U\) and the transmutation
record at \(E_{\mathrm{cell}}\); (AR2) the reconciliation cost is quadratic
in renormalization time; and (AR3) the two anchor records carry equal
capacity.  These premises give a strictly convex cost with a unique minimizer
at the log-midpoint \(\sqrt{\mu_U E_{\mathrm{cell}}}=E_\star e^{-\pi}P^{-1/6}\),
hence \((m_t,m_H)=(172.63,\,125.77)\)~GeV at two loops; the implication is
proved exactly on rational sample points
(\texttt{d11\_boundary\_scale\_midpoint\_selection\_theorem.json}).

Two of the three premises reduce to the axioms.  The quadratic-cost premise
AR2 is a theorem under the canonical record model: the one-loop chart is
inverse-affine, so \(1/\alpha(t)\) is exactly linear in renormalization time,
and the Kullback--Leibler divergence between fixed-resolution Gaussian records
with an affine stored coordinate is exactly
\((s^2/2\sigma^2)(t-t_i)^2\), a quadratic cost with no leading-order
approximation.  The reconciliation placement in AR1 is also a theorem: the
mismatch functional is port-additive by construction and the repair move
settles a record at the cost minimizer, so a record with two scale-parents
settles at the capacity-weighted log-mean.  The equal-capacity premise AR3
reduces to equal slope and readback resolution, which same-class registers at
equal refinement depth supply. Two finite carrier facts are open for
the D11 carrier: CF1, that the criticality boundary record carries exactly two
parent ports, one to the gauge-unification register and one to the
transmutation register; and CF2, that those two anchor registers are the same
register class at equal refinement depth.  Given the canonical record model and
CF1 and CF2, the axioms force the boundary scale, and with it
\((m_t,m_H)=(172.63,\,125.77)\)~GeV, with no free choices. CF1 and CF2 are
statements about the same D11 carrier census certificate (of the QT1--QT5
class) required by the \(W/Z\) value law; one certificate closes the
Higgs-scale selection, the \(W/Z\) law, and the D10 two-law choice together.
The equal-capacity assumption is measurable: a capacity asymmetry moves the
selected scale off the midpoint by a computable amount, about \(2.1\)~GeV in
\(m_H\) per \(e\)-fold, and the registered three-loop implied scale reads it off
(\texttt{d11\_anchor\_reconciliation\_reduction\_theorems.json}).  The candidate
registry is frozen before any three-loop computation exists, so the three-loop
implied scale is the registered discriminating test.

A conditional selection theorem accompanies the registry.  If the boundary
record reconciles the two anchor records with port-additive quadratic cost
in renormalization time and equal anchor capacities, the boundary scale is
exactly the log-midpoint.  Two of the three premises reduce to the axioms'
record and repair structure: the one-loop chart is inverse-affine in
\(\ln\mu\), so the canonical Gaussian MaxEnt record model gives an exactly
quadratic reconciliation cost, and repair minimization of port-additive
mismatch places a two-parent record at the capacity-weighted log-mean.  The
surviving content is two finite carrier facts, the two-parent port structure
of the boundary record and the equal-class equal-depth status of the anchor
registers; a carrier certificate of the census class closes both
(\texttt{d11\_boundary\_scale\_midpoint\_selection\_theorem.json},
\texttt{d11\_anchor\_reconciliation\_reduction\_theorems.json}).

\subsection{Target-conditioned electroweak chart envelope}\label{ew-conditional-envelope}

Evaluating two target-conditioned repair candidates at the endpoint pixel and across
the empirical-closure interval gives the following comparison envelope.
It is not a source interval: its candidate set is not an exhaustive
source-derived deformation class, and it mixes the public endpoint lane with
the strict source-audit question.
\[
\begin{aligned}
m_H &\in [125.183,\,125.232]~\mathrm{GeV}, & &\text{measured } 125.13\pm0.11,\\
m_t &\in [172.278,\,172.352]~\mathrm{GeV}, & &\text{measured } 172.1\pm0.6,\\
c_W &\in [80.3692,\,80.3774]~\mathrm{GeV}, & &\text{PDG BW coordinate } 80.3692\pm0.0133,\\
c_Z &\in [91.1880,\,91.1983]~\mathrm{GeV}, & &\text{PDG BW coordinate } 91.1880\pm0.0020.
\end{aligned}
\]
The Higgs and top entries are target-conditioned comparisons on their declared mass scheme. The
\(W/Z\) entries are chart coordinates. The PDG targets use mass-dependent-width Breit--Wigner
parameters, and complex-pole masses use another convention, so no physical coverage or pull is
assigned. At the endpoint pixel the two repair selections give \(Z\)-chart coordinates
\(91.18798\) and \(91.18801\,\mathrm{GeV}\). The selection spread between the two displayed laws is
roughly eight \(\mathrm{MeV}\) in the \(W\) chart, thirty \(\mathrm{keV}\) in the \(Z\) chart, sixteen
\(\mathrm{MeV}\) in \(m_H\), and thirty \(\mathrm{MeV}\) in \(m_t\).  These
figures compare two ans\"atze; they do not bound the full QT5 deformation
class.

\subsection{Residual attribution across the electroweak rows}\label{ew-residual-attribution}

On this comparison envelope, the Higgs and top residuals are target-conditioned mass-scheme
diagnostics. No \(W/Z\) standard-deviation assignment is valid because the chart coordinates,
PDG mass-dependent-width parameters, and complex poles use different conventions.
The \(W\)-chart spread is repair-selection freedom: the two candidate repair laws differ by about
eight \(\mathrm{MeV}\). The \(m_H\) and \(m_t\) offsets track the same repair-selection width folded through the
declared Higgs and top Jacobian cores.

The \(Z\)-chart row is the most sensitive adapter case. The empirical-closure pixel inherits
the frozen electromagnetic anchor deficit through the fine-structure endpoint, and \(c_Z\) tracks
the pixel with \(\mathrm{d}c_Z/\mathrm{d}P\approx123~\mathrm{GeV}\) per unit \(P\). Two
propagation branches separate the mechanism. The direct branch injects the certified anchor shift
into the running couplings and re-evaluates the candidate \(W/Z\) chart. That branch shifts
\(c_Z\) by \(-66\) to \(-87~\mathrm{MeV}\) on the hypercharge line and by roughly \(-230\) to
\(-302~\mathrm{MeV}\) on the proportional line, exceeding the adapter displacement by an order of magnitude,
so it is excluded. The pixel branch moves the endpoint to the measured fine-structure value and
propagates the induced pixel shift through \(\mathrm{d}P/\mathrm{d}A_{\mathrm{Th}}\). That branch
moves the empirical pixel onto the calibration pixel to within \(4\times10^{-7}\) in \(P\),
equivalently \(-0.05~\mathrm{MeV}\) in the \(Z\) chart. This is an internal adapter closure check;
it supplies no physical mass or pole promotion.

For this comparison surface, a closed fine-structure anchor bridge would
collapse the endpoint pixel interval onto the calibration pixel and remove
this particular chart displacement. It would not close QT1--QT5, the independent
scale, the RG/scheme certificate, or the complex-pole gate. The anchor bridge is the object developed
in the fine-structure companion, and its source branch is the subject of Section~\ref{ew-anchor-source-branch}.

\subsection{Source branch of the fine-structure anchor bridge}\label{ew-anchor-source-branch}

The certified anchor gap \([0.649,\,0.855]\) in inverse fine-structure units names the shift the
source side supplies at the electroweak anchor for the empirical-closure endpoint to reach the
measured fine-structure value. The gap is defined by that measured requirement, so inserting it
back is circular. The non-circular source route continues the declared running convention to the
next order and tests whether the induced anchor shift lands inside the certified gap.

That test is executed. The one-loop anchor is reproduced from the declared coefficients and the
transmutation-certificate boundary data to a residual of \(10^{-11}\). The standard two-loop
running system, integrated down the same \(33.2\) e-folds with a top-Yukawa convention scan, shifts
the anchor by \([+1.62,\,+2.14]\) in inverse fine-structure units for every convention. The shift
has the correct sign and overshoots the certified gap by a factor of two to three. Restoring the
balance requires the threshold and scheme-conversion data, which the matching contract classifies
as hidden fit parameters when they are undeclared. The balance requirement is quantitative: the
threshold map removes about \(2.14\) inverse-alpha, equivalent under naive coefficient accounting
to a single effective threshold near \(458~\mathrm{GeV}\). The anchor bridge has no source
scheme-lock or threshold map; its endpoint is empirical.

\section{Flavor Transport, Generation Structure, and the Yukawa Dictionary}

The gauge branch explains why the realized low-energy sector has three generations, but it does
not by itself produce the full flavor dictionary. This paper therefore needs a
separate flavor chapter. The imported realized-branch input here is the D9 count package
\(N_c=3\) and \(N_g=3\). This chapter is the bridge between that structural result and the
matter-family chapters for quarks, charged leptons, and neutrinos.

The flavor derivation is deliberately constructive. It does not claim that the full
OPH flavor observable is theorem-level. Instead, it builds the object chain that a
closed theorem would have to pass through. The derivation chain has the following main
architecture: start from a refinement-indexed family transport kernel, derive a centered
generation-bundle branch generator, lift that to same-label transport data, derive the induced
overlap-edge transport cocycle, reduce the cocycle to a persistent flavor observable, and then
push that observable into common sector-response objects for the downstream quark, charged-lepton,
and neutrino derivations.

Mathematically, this derivation is where ``why this family exists'' becomes a concrete technical
question. The realized Standard Model branch gives three generations structurally; the flavor
derivation is where one tries to turn that three-generation existence statement into transport,
splitting, suppression, phase, and excitation data. The relevant objects are the intermediate
transport and spectral structures, not the final fermion masses themselves, that the mass
readouts consume.

The active chain, in slightly compressed form, is:
\begin{enumerate}[leftmargin=2em]
\item normalize a refinement-indexed family transport kernel;
\item derive a centered generation-bundle branch generator on the realized three-generation
charged bundle;
\item lift this to same-label edge-line transport and then to an overlap-edge cocycle with
explicit defect and gap bookkeeping;
\item reduce those data to projectors, spectral gaps, pair suppressions, and cycle phases;
\item push the resulting family object into sector-response objects and then into
suppression/phase tensors for the downstream matter derivations.
\end{enumerate}

Three lane-specific claim boundaries are attached to this region of the derivation. The shared
excitation dictionary is the common proof-facing base. Above that base, the charged lane carries
exact centered readback, a closed common-shift no-go, the declared same-label \(q_e\) readback, a
source-side determinant character for any fixed formal exponent vector
\[
S_M=\sum_e M_e^{\mathrm{ch}}\log q_e,
\]
a determinant-line lift on theorem-grade physical charged data, and an algebraic mass
readout from theorem-grade \(A_{\mathrm{ch}}(P)\). The theorem lane does not emit a theorem-grade
sector-isolated charged determinant exponent vector, and it does not identify a source-side
determinant character with the physical charged determinant line. The charged-lepton theorem gap is the
determinant trace-lift attachment \(3\mu(r)=S_M(r)\) on the charged determinant channel. The charged determinant channel has a corpus-limited no-go boundary: the uncentered trace lift is not emitted. The neutrino lane carries a target-informed weighted-cycle
candidate above a template family kernel, with compare-only bridge and absolute-attachment diagnostics. The quark
lane carries a theorem-grade source-spread obstruction: its ordered profile shapes leave a free
\((\mathbb R_{>0})^2\) fiber. Selected-class descent proves representative independence but does
not choose either positive modulus. Target-anchored mixed-convention mass coordinates and
GeV-valued mass textures are audit data; they are neither a source-only running sextet nor
physical dimensionless Yukawa matrices. A separate compare-only microphysical bridge records
edge-statistics transport on a diagnostic surface. Together these boundaries mark where the flavor
derivation leaves the common transport backbone and passes to lane-specific closure contracts.

\section{Quark Family Derivation}

The quark lane carries a source-only non-identifiability theorem on the public physical quark frame
class \(f_P\) chosen by \(P\). The source equations determine the ordered up- and down-sector
profile rays but leave their endpoint spans as two independent positive moduli. Selected-class
descent does not remove this freedom, so no numeric mass row is emitted. Same-family and restricted
common-refinement artifacts reproduce their chosen targets only after those moduli are obtained by
target inversion. Their rows also mix renormalization conventions, while their GeV-valued matrices
are mass textures rather than physical dimensionless Yukawa matrices. This theorem does not classify
all public quark frame classes.

\subsection{What quarks are in this derivation}

Quarks are the color-charged elementary constituents of hadronic matter. In nature, the up quark
and down quark dominate protons and neutrons, while the strange quark, charm quark, bottom
quark, and top quark appear in progressively heavier and more unstable sectors. In the OPH
particle derivation, the quark lane has a closed downstream algebraic readout conditional on a
source-only physical spread datum. The target-free source equations do not emit that datum: they
fix two ordered profile rays and leave their positive endpoint spans independent. The compatible
fiber is \((\mathbb R_{>0})^2\), so the six numeric rows are not source-only predictions.

The quark route starts from the shared flavor excitation dictionary, internalizes the target-free
mass bridge on the D12 mass ray, proves selected-bridge-fiber representative independence for the
attached physical spread datum on \(f_P\), and applies the affine mean law once both positive
moduli are supplied. The stored target-anchored matrices have GeV-valued singular values drawn
from several comparison conventions. They are mass textures, not physical dimensionless Yukawa
matrices.
The same-label left-handed selector surface closes to the singleton
\(\sigma_{\mathrm{ref}}\). That lower object is a negative sheet-selector statement
on the selected D12 sheet. It does not break the independent positive-rescaling action on the two
sector profiles.

\subsection{Emitted quark rows}

\begin{table}[h]
\centering
\small
\begin{tabular}{@{}llll@{}}
\toprule
Quark & Claim tier & Public value & Derivation stage \\
\midrule
Up quark & source spread non-identifiable & withheld & \(\overline{\mathrm{MS}}\), \(\mu=2\,\mathrm{GeV}\), comparison chart \\
Down quark & source spread non-identifiable & withheld & \(\overline{\mathrm{MS}}\), \(\mu=2\,\mathrm{GeV}\), comparison chart \\
Strange quark & source spread non-identifiable & withheld & \(\overline{\mathrm{MS}}\), \(\mu=2\,\mathrm{GeV}\), comparison chart \\
Charm quark & source spread non-identifiable & withheld & \(\overline{\mathrm{MS}}\), \(\mu=m_c(\mu)\), comparison chart \\
Bottom quark & source spread non-identifiable & withheld & \(\overline{\mathrm{MS}}\), \(\mu=m_b(\mu)\), comparison chart \\
Top quark & separate extraction coordinate & withheld & cross-section pole-mass chart \\
\bottomrule
\end{tabular}
\end{table}

Running quark masses are coordinates on a renormalization-group trajectory. Finite
renormalizations and scale changes alter those coordinates while preserving physical amplitudes,
so OPH can emit an RG-covariant trajectory or invariant but cannot derive the human choice of an
\(\overline{\mathrm{MS}}\) chart. The chart must be declared after source emission. The top row is
a pole extraction coordinate rather than a sixth member of one common running-mass packet. All six
rows are withheld because the spread pair is non-identifiable from the source corpus. A physical
Yukawa construction would additionally transport every running coordinate to one common scale,
convert the top coordinate, emit the running Higgs expectation value in the same scheme, and apply
\(y_q(\mu)=\sqrt2\,m_q(\mu)/v(\mu)\).

\subsection{Target-free mass bridge and selected-class public theorem}\label{quark-family-strongest-overlap-branch-primitives}

The same-label left-handed solver surface closes to the singleton
\(\sigma_{\mathrm{ref}}\) (\ophid{sigma_ref}). This is a negative same-sheet selector statement:
same-sheet rephasing preserves CKM moduli, so it does not move that selected sheet to the
physical CKM shell. Same-sheet overlap scans, chirality-swapped basis diagnostics, and other
non-sector-attached orbit improvements are compare-only and do not change the selected-class
theorem.

A separate target-free mass bridge is internalized on the emitted D12 ray. On the minimal light branch
\[
y_u=c_u\,\varepsilon^6,
\qquad
y_d=c_d\,\varepsilon^6,
\qquad
\varepsilon=\frac16,
\]
the light-quark overlap-defect theorem emits
\[
\Delta_{ud}^{\mathrm{overlap}}
=
\frac16\log\frac{c_d}{c_u}.
\]
The emitted same-family D12 mass object is the ray
\[
D12_{ud}^{\mathrm{mass}}
\equiv
\mathcal R_{D12}^{ud}
\qquad
\text{(paper id \ophid{D12_ud_mass_ray})},
\]
and on that ray the same scalar is equivalently the one-scalar law
\[
\Theta_{ud}^{\mathrm{mass}}
:=
\text{\ophid{quark_same_family_value_law}}.
\]
The exact scalar identities are
\[
\Delta_{ud}^{\mathrm{overlap}}=\frac{t_1}{5},
\qquad
\log\frac{c_d}{c_u}=\frac65\,t_1,
\qquad
t_1=5\,\Delta_{ud}^{\mathrm{overlap}}=\frac56\log\frac{c_d}{c_u}.
\]
The D12 mass-side package is therefore functorial:
\[
\begin{aligned}
\eta_Q^{\mathrm{centered}}&=-\frac{1-x_2^2}{27}\,t_1,\\
\kappa_Q&=-\frac{t_1}{54},\\
x_2&=-0.5175863354681689.
\end{aligned}
\]
The odd source package is also forced:
\[
\begin{aligned}
\beta_{u,\mathrm{diag},B}^{\mathrm{source}}&=\frac{t_1}{10},\\
\beta_{d,\mathrm{diag},B}^{\mathrm{source}}&=-\frac{t_1}{10},
\end{aligned}
\]
\[
\begin{aligned}
\texttt{source\_readback\_u\_log\_per\_side}
&=
\left(-\frac{t_1}{10},\,0,\,+\frac{t_1}{10}\right),\\
\texttt{source\_readback\_d\_log\_per\_side}
&=
\left(+\frac{t_1}{10},\,0,\,-\frac{t_1}{10}\right),
\end{aligned}
\]
For any supplied positive spread pair one obtains
\[
\tau_u=\frac{\sigma_d}{10(\sigma_u+\sigma_d)}\,t_1,
\qquad
\tau_d=\frac{\sigma_u}{10(\sigma_u+\sigma_d)}\,t_1.
\]
These identities do not select \(\sigma_u\) or \(\sigma_d\). The mass bridge is not part of the
selected-class theorem boundary. The builder-facing pure-\(B\)
payload pair is a lower implementation object beneath the exact theorem surface, and
\ophid{intrinsic_scale_law_D12} is the derived wrapper above \(\Theta_{ud}^{\mathrm{mass}}\).

Write the same-label left-handed physical carrier as
\[
\Sigma_{ud}^{\mathrm{phys}}
:=
\left\{
(\sigma_{\mathrm{id}},\tau,U_{u,L},U_{d,L},V_{\mathrm{CKM}},I_{\mathrm{CKM}})
:
V_{\mathrm{CKM}}=U_{u,L}^\dagger U_{d,L}
\right\}/\!\sim,
\]
where
\[
(U_{u,L},U_{d,L},V)\sim
(U_{u,L}D_u,\ U_{d,L}D_d,\ D_u^\dagger V D_d)
\]
for diagonal \(D_u,D_d\in U(1)^3\). On the selected public quark frame class \(f_P\),
represented on the realized lane by the explicit common-refinement transport-frame class
\([F_0^\dagger F_1]\), the exact \(\Sigma_{ud}^{\mathrm{phys}}\) element and the attached exact
sigma datum are independent of the declared representative inside the selected bridge fiber.
Therefore the exact sigma readout descends uniquely to the selected audit/support surface on
\(f_P\). It is not a public source-only mass prediction because the physical sigma datum is
target-derived.

\paragraph{Sigma descent non-selection.}
Let \(R_{\mathrm{decl}}(f_P)\) be the declared selected bridge fiber over the public quark
frame class. If a map
\[
\Sigma:R_{\mathrm{decl}}(f_P)\to\mathbb R^4
\]
is constant on that fiber, then it descends to a well-defined public datum
\(\overline\Sigma(f_P)\). This proves representative independence only. It does not select the
value of \(\overline\Sigma(f_P)\), because every constant vector in \(\mathbb R^4\) would also
descend. Therefore a value obtained from the declared running-quark target surface is
target-derived after descent.

\paragraph{Common-scale physical rejection of the reciprocal-ray candidate.}
A physical comparison first converts all six quark eigenvalues to dimensionless Yukawa singular
values in one effective theory, scheme, and scale. Using the Standard Model
\(\overline{\mathrm{MS}}\) values at \(M_Z\) tabulated by Antusch, Hinze, and
Saad~\cite{antusch2025running}, the two ordered log-gap ratios are
\[
\rho_u=1.1108888543,
\qquad
\rho_d=0.7519410008,
\qquad
\rho_u\rho_d=0.8353228768.
\]
The reciprocal-ray law requires the last product to equal one. It is therefore falsified on the
properly typed target surface. Even after granting the four endpoint Yukawas
\((y_u,y_t,y_d,y_b)\), the minimax reciprocal shape misses the held-out \(y_c\) and \(y_s\) by
\(21.556\%\) at \(M_Z\). The result persists under running: the products are
\(0.836092\), \(0.837752\), \(0.845227\), and \(0.843611\) at
\(10^3\), \(10^5\), \(10^{12}\), and \(10^{16}\,\mathrm{GeV}\), with best held-out errors
between \(19.93\%\) and \(21.42\%\). The stored \(\rho=1.294285\) formula misses the charm
coordinate by \(56.5\%\) at \(M_Z\), even under the endpoint grant.

The generic physical interface is therefore
\[
(\mu_u,\sigma_u,\rho_u,\mu_d,\sigma_d,\rho_d)\in\mathbb R^6,
\]
with an exact inverse to the six ordered positive Yukawas. The three-scalar theorem is exact only
inside its imposed reciprocal-ray, affine-mean, fixed-shared-scale subfamily. The two-spread
counterfamily below is a valid restricted non-identifiability lower bound, but it
is not a parameterization of the whole physical interface.

A separate selector audit sharpens the missing source object. Generation-blind flavor-singlet
scalars cannot emit a fixed nonzero bifundamental Yukawa matrix. Simultaneous \(S_3\) invariance
restricts a fixed matrix to the span of \(I\) and \(J\), which has a degenerate eigenspace and
aligned up/down frames. A physical construction therefore needs a source-derived flavor-orbit
selector, such as a spontaneous invariant functional or genuine bifundamental boundary data,
together with a quark--Higgs carrier. The calibration simulator contains no quark, Higgs, or
Yukawa coupling and its enumerated observables have zero Fisher information for these coordinates.
Its \(64\mathrm{k}\) runs are at the calibration null: the two fusion-tower readings
are \(1.495245\) and \(1.484301\), while the dense first/full readings are \(1.498275\) and
\(1.497823\). Those values are not quark-mass evidence.

\paragraph{Restricted source-spread non-identifiability theorem.}
Remove every running-mass target, exact-witness, fitted-spread, and compare-only ancestor. Grant
the remaining source packet its strongest ordered three-point shape law. For
\(\rho_{\mathrm{ord}}>0\), define
\[
v_u=\frac{1}{3(1+\rho_{\mathrm{ord}})}
\bigl(-(2\rho_{\mathrm{ord}}+1),\rho_{\mathrm{ord}}-1,
\rho_{\mathrm{ord}}+2\bigr),
\]
\[
v_d=\frac{1}{3(1+\rho_{\mathrm{ord}})}
\bigl(-(\rho_{\mathrm{ord}}+2),1-\rho_{\mathrm{ord}},
2\rho_{\mathrm{ord}}+1\bigr).
\]
Both vectors have zero trace and unit endpoint span. Their adjacent-gap ratios are
\(\rho_{\mathrm{ord}}\) and \(\rho_{\mathrm{ord}}^{-1}\), respectively. Conversely, those
three conditions determine each profile up to one positive endpoint span. Hence every compatible
pair is
\[
E_u=\sigma_u v_u,\qquad E_d=\sigma_d v_d,
\qquad (\sigma_u,\sigma_d)\in(\mathbb R_{>0})^2.
\]
The group \((\mathbb R_{>0})^2\) acts freely and transitively by independent rescaling of the two
spans while fixing the source shape data. The requested four-tuple
\[
\left(\frac{\sigma_u+\sigma_d}{2},
\frac{\sigma_u-\sigma_d}{2},\sigma_u,\sigma_d\right)
\]
is not invariant under this action. No unique source-only spread package follows from the stated
corpus.

\paragraph{No-extra-axiom MAR non-definability theorem.}
This ambiguity is not merely a missing implementation.  Fix any generic MAR-admissible one-Higgs
three-generation package and write
\[
Y_q=U_{q,L}\,\operatorname{diag}
\!\left(e^{\mu_q+\sigma_qv_{q,1}},e^{\mu_q+\sigma_qv_{q,2}},
e^{\mu_q+\sigma_qv_{q,3}}\right)U_{q,R}^{\dagger},
\qquad q\in\{u,d\},
\]
with \(\sum_i v_{q,i}=0\).  For every \(\lambda_u,\lambda_d>0\), replace
\(\sigma_q\) by \(\lambda_q\sigma_q\) while holding the frames fixed.  Gauge representations,
anomaly cancellation, hypercharges, one-Higgs Yukawa completability, the CKM matrix, intrinsic CP
capability, and the weak-sector counting clause are unchanged.  In addition
\[
\det\exp(\lambda_q\sigma_qv_q)
=\exp\!\left(\lambda_q\sigma_q\sum_i v_{q,i}\right)=1,
\]
so determinant normalization does not select either modulus.

MAR assigns every member the same complexity vector
\[
C(\mathfrak S)=(\chi_{\mathrm{cpl}},N_{\mathrm{nonab}},N_c,N_g),
\]
which contains no Yukawa eigenvalue.  The members are nevertheless physically inequivalent because
the declared physical-equivalence relation preserves Yukawa invariants.  Axiom~3 also does not
remove the family: its gauge-invariant local constraint values and associated multipliers are not
numerically specified and no map from \(P\) to quark Yukawa multipliers is emitted.  Thus Axioms
1--5 plus fixed \(P\) admit a free \((\mathbb R_{>0})^2\) family of equal-MAR-score quark spectra.
There is no smallest positive rescaling; its infimum is zero, which would give the massless or
degenerate boundary rather than the observed spectrum.  Therefore the stated axioms do not define
a unique quark mass spectrum.  This conclusion uses no additional axiom and can be overturned only
by deriving, from the existing axioms, a source functional that is nonconstant on this explicit
counterfamily.

\paragraph{The kernel interface: the exact content of the missing derivation.}
Two companion theorems make that boundary executable. First, the admissibility battery recorded
for the family-transport lane (positive-semidefinite hermitian descendants, open three-cluster
gaps at every level, a simple centered spectrum, the conjugacy-Riesz margin, persistent projector
labeling, and overlap-edge amplitudes above the floor) places no constraint on the spectrum of
the centered compressed branch generator: for every pair \((r,s)\in(\mathbb R_{>0})^2\) there is a
certificate-passing two-level kernel whose generator has raw gap ratio exactly \(r\) and spectral
span exactly \(s\). The construction places the target spectrum by unitary conjugation,
\(T=Q\,\operatorname{diag}(\sqrt{\mu})\,Q^{\dagger}\), so the descendant spectrum is exact, and
shrinks the refinement drift geometrically until the Riesz margin passes. The persistence
certificates are therefore a filter, not a generator: a kernel derivation that merely passes them
cannot emit the ordered ratio constant, the mean-law coordinate, or the spans.

Second, the three-scalar interface on record is exact only on its imposed reciprocal-ray
subfamily.  At fixed shared scale \(g_{\mathrm{ch}}\), that subfamily factors through
\[
(r,\sigma_u,\sigma_d)\in(\mathbb R_{>0})^3,
\qquad
\rho_{\mathrm{ord}}=\frac{3}{2+r},
\qquad
x_2=\frac{r-1}{r+1},
\]
with the rays, the affine mean coefficients \(A_{ud},B_{ud}\), the sector means, and the six
coordinates closed-form in the triple, and with the explicit left inverse
\(\sigma_u=\tfrac12\ln(m_t/m_u)\), \(\sigma_d=\tfrac12\ln(m_b/m_d)\), and \(r=3/\rho-2\) at
\(\rho=\ln(m_c/m_u)/\ln(m_t/m_c)\); the forward map is injective and the executed round trip
closes at machine precision inside that subfamily.  It is not the general interface of two
ordered three-point spectra.

Indeed, for any \(x\ne\pm1\), let
\[
L(x)=\operatorname{ctr}(-1,x,1),
\qquad
Q(x)=\operatorname{ctr}(1,x^2,1).
\]
Their Gram determinant is
\[
\det\operatorname{Gram}(L,Q)=\frac43(1-x^2)^2.
\]
Thus \(L,Q\) form a basis of the centered three-vector plane, and every centered sector spectrum
has unique coordinates \(E_q=a_qL+b_qQ\).  Once both \(a_q\) and \(b_q\) are free, \(x\), and
hence \(r\), is a basis choice rather than an additional invariant of the spectrum.  The general
common-scale eigenvalue interface has six scalar coordinates: two centered coordinates and one
mean for each sector.  A proposed \(r\) plus four centered coordinates plus two means is a
redundant seven-coordinate chart. The reciprocal-ray law removes one centered coordinate per sector and
thereby recovers its special three-scalar chart.

By the freedom theorem and the non-definability theorem above, the corpus selects neither
the ray triple nor the six-scalar physical interface. The executable acceptance harness is
only a fail-closed score of the reciprocal-ray subfamily
(\texttt{family\_transport\_kernel\_admissibility\_freedom\_theorem.json},
\texttt{quark\_kernel\_three\_scalar\_interface\_theorem.json},
\texttt{quark\_kernel\_normalization\_acceptance\_current.json}).

The edge data do not remove the ambiguity. Even after granting source values
\(S_{13},S_{23},\delta_{21}>0\), every positive pair can be written as
\[
\sigma_u=S_{13}+c_u\delta_{21},\qquad
\sigma_d=S_{23}+c_d\delta_{21}
\]
for suitable \((c_u,c_d)\). The source equations emit no rule fixing these two coefficients. The
specific coefficients used by the numerical candidate are therefore a declared ansatz from a
hand-written family-transport template, not an OPH-derived kernel.

Define
\[
\sigma_{\mathrm{seed}}^{ud}:=\frac{\sigma_u+\sigma_d}{2},
\qquad
\eta_{ud}:=\frac{\sigma_u-\sigma_d}{2},
\]
\[
\begin{aligned}
A_{ud}&:=\frac{1}{2(1+\rho_{\mathrm{ord}}-x_2^2)},\\
B_{ud}&:=\frac{1}{2\!\left(1-x_2^2-\frac{x_2^2}{1+\rho_{\mathrm{ord}}}\right)},\\
\rho_{\mathrm{ord}}&=1.294284936377706.
\end{aligned}
\]
Then
\[
g_u
=
g_{\mathrm{ch}}
\exp\!\bigl(-(A_{ud}\sigma_{\mathrm{seed}}^{ud}-B_{ud}\eta_{ud})\bigr),
\qquad
g_d
=
g_{\mathrm{ch}}
\exp\!\bigl(-(A_{ud}\sigma_{\mathrm{seed}}^{ud}+B_{ud}\eta_{ud})\bigr).
\]
The Jacobian from \((\sigma_u,\sigma_d)\) to
\((\log(g_u/g_{\mathrm{ch}}),\log(g_d/g_{\mathrm{ch}}))\) has determinant
\(-A_{ud}B_{ud}\neq0\). The two free moduli therefore change the mass readout; they are not a
gauge redundancy. Given a source-only physical spread datum, the affine mean law emits
\((g_u,g_d)\) algebraically on \(f_P\). The target-derived packet is retained as a mixed-convention
audit witness. Its dimensionful matrices are mass textures. The conditional selected-class route
can be written as
\[
\!f_P+\Sigma_{ud}^{\mathrm{source}}(P)
\Longrightarrow
\bigl(\sigma_u,\sigma_d,\sigma_{\mathrm{seed}}^{ud},\eta_{ud}\bigr)_{\mathrm{phys}}
\Longrightarrow
(g_u,g_d)
\Longrightarrow
\bigl(m_u,m_d,m_s,m_c,m_b,m_t\bigr),
\]
followed, only after common-scale RG transport and division by the running Higgs expectation
value, by dimensionless Yukawa matrices. The first arrow is obstructed by the free
\((\mathbb R_{>0})^2\) action on the stated source corpus.

\paragraph{Target-anchored \(S_3/D12\) two-mode witness.}
An ansatz built with visible target coordinates reproduces the six mixed-convention comparison
coordinates numerically. Its exact mathematical component starts
from the transposition Cayley graph of \(S_3\).  The adjacency spectrum is
\(3,0,-3\) with multiplicities \(1,4,1\), so the Laplacian spectrum is \(0,3,6\) and
\[
\frac{e^{-3\tau}-e^{-6\tau}}{1-e^{-3\tau}}=e^{-3\tau}.
\]
This identity is a finite-group theorem. No OPH theorem identifies the three distinct
isotypic heat values with three generations or supplies the proposed heat time
\[
\tau_f=\frac P4-\frac{\pi\alpha_U}{5}.
\]

The ansatz sets \(w=\pi\alpha_U\), \(r=e^{-3\tau_f}\),
\(\rho=3/(2+r)\), and \(x=(r-1)/(r+1)\), and then uses
\[
\begin{aligned}
e_u&=S_{13}+\frac{\rho\delta_{21}}{1+\rho},&
e_d&=S_{23}+\frac{\delta_{21}}{2(1+\rho-x^2)},\\
\bar\sigma_u&=e_u-w\Delta S_{13},&
\bar\sigma_d&=e_d-w(1-\Delta S_{13}),\\
a_u&=e_u+\frac w2,&
a_d&=e_d-\frac w5,\\
b_u&=b_u^{\mathrm{ray}}-\frac{w\rho}{10},&
b_d&=b_d^{\mathrm{ray}}-\frac w4.
\end{aligned}
\]
The ray coordinates have the symbolic form
\[
b_u^{\mathrm{ray}}
=a_u\frac{-\rho x+\rho-x-1}{(1+\rho)(x^2-1)},
\qquad
b_d^{\mathrm{ray}}
=a_d\frac{-\rho x-\rho-x+1}{(1+\rho)(x^2-1)}.
\]
The ansatz feeds \(\bar\sigma_u,\bar\sigma_d\) into the candidate affine mean law above and
sets \(E_q=a_qL+b_qQ\).  With the frozen repository numbers, this gives
\[
\begin{gathered}
r=0.3179975211,\qquad \rho=1.2942205385,\qquad x=-0.5174535369,\\
(a_u,b_u)=(5.6430129216,-4.9927828217),\qquad
(a_d,b_d)=(3.3952170977,-1.8369623926).
\end{gathered}
\]

\paragraph{Conditional normalized-trace lemma.}
Let \(\ell:M_d(\mathbb C)\to\mathbb C\) be complex linear, invariant under all unitary
conjugations, and normalized by \(\ell(I)=1\).  Unitary conjugacy of rank-one projectors and
additivity over an orthonormal resolution of \(I\) give
\(\ell(A)=\operatorname{Tr}(A)/d\), so every rank-one channel has weight \(1/d\).
Likewise, a width operator on \(M_d(\mathbb C)\) invariant under the full
\(U(d)\times U(d)\) left-right action is scalar by Schur's lemma; total Hilbert--Schmidt width
\(t\) therefore assigns \(t/d^2\) to each unit slot.  Consequently, \emph{if} the physical heat,
up-odd, down-odd, up-even, and down-even responses are respectively identified with one isotropic
slot in \(M_5(\mathbb C)\) of total width \(5w\) and rank-one modules of dimensions
\(2,5,10,4\), the coefficients \(w/5,w/2,w/5,\rho w/10,w/4\) follow.
The full invariance and normalization of the response law, those five module assignments, their
signs and orientation, and exclusion of competing assignments are hypotheses outside OPH.
Thus this lemma closes the denominator arithmetic conditional on the physical channel
functor; it does not remove target dependence from the formula.

The per-particle comparison table is excluded from the public prediction ledger. It is a
target-anchored diagnostic: the evaluator emits dimensionless coordinates, inserts
an unproved one-GeV unit, and compares them to a table that mixes light-quark
\(\overline{\mathrm{MS}}\) values at \(2\,\mathrm{GeV}\), heavy-quark self-scale values, and a
separate top extraction.

The raw diagonal residual sum against the supplied central values and nominal standard deviations
is \(1.1653\), with a maximum relative residual of \(0.2946\%\).  It is not a likelihood or
goodness-of-fit statistic: discovery used these targets, there is no source-theory covariance, and
the rows are not one common-scale spectrum. In the target-conditioned \(219{,}615\)-member denominator
grammar, the selected tuple \((5,2,5,10,4)\) is the unique minimum of that raw residual sum, but
eight formulas tie its best maximum error. The grammar contains the target-exposed inputs, graph,
entries, assignments, signs, and functional forms, so it is not a global look-elsewhere correction.

The evaluator is runtime-target-separated but target-dependent.
\(S_{13},S_{23},\delta_{21},\Delta S_{13}\), and \(g_{\mathrm{ch}}\) descend from
the explicitly hand-written \texttt{family\_transport\_kernel} template.  The last quantity is the
dimensionless template-eigenvalue mean plus its minimum gap; treating it as GeV supplies the
missing unit by hand.  The selected pixel branch depends on an
internal quark-spectrum continuation.  The edge laws add log-overlap
suppression to a linearly scaling \(TT^\dagger\) gap, and \(\Delta S_{13}\) is a selected
basis-dependent matrix entry. The manifest classifies the ansatz as target-informed and disables promotion.
It neither contradicts the non-identifiability theorem nor supplies a physical mass law.

\paragraph{Sufficient Flavor Source Closure contract.}
The proof audit isolates a sufficient conditional completion without instantiating it.
A physical theorem would require all of the following on one acyclic, target-free source DAG:
\begin{enumerate}[leftmargin=2em]
\item a unique interval-certified source root for \(P,\alpha_U(P)\), with no dependency on an
internal quark continuation;
\item a physical \(S_3\) family-carrier attachment, the heat-time law
\(\tau_f=P/4-\pi\alpha_U/5\), the \(\rho,x\) dictionary, and the two edge-response laws;
\item a source-labeled simple-spectrum family generator and charged seed with exact
common-refinement intertwiners, monomial transport of the label rays, positive common
normalization, and selected edge magnitudes bounded away from zero;
\item a physical response-channel functor that proves the normalized-trace hypotheses, module
assignments, signs, orientation, competing-channel exclusion, and the affine mean law on that
carrier, on a domain where the \(A,B\) denominators do not vanish;
\item a common-scale physical readout with
\(m_q(\mu)=y_q(\mu)v(\mu)/\sqrt2\), including field normalization, units, and the running Higgs
expectation value in the same scheme;
\item a source-only RG packet with existence and non-blowup across the required intervals, a
locally Lipschitz beta system, frozen threshold ordering and matching maps, top conversion, and
the declared comparison charts.
\end{enumerate}
Given these hypotheses, composition of the single-valued maps makes the final sextet unique.
This is a sufficient proof-obligation contract, not a proved-minimal OPH theorem.
The numerical table is not its corollary: the evaluator's coordinates are dimensionless and no
instance of the final two receipts exists. What the
counterfamily proves as a necessary condition is narrower and decisive: any successful source
functional must be nonconstant on the independent
\((\lambda_u,\lambda_d)\) centered-spread rescaling orbit.

\paragraph{RSCC as a target-conditioned specification.}
Representation-Slot Cumulant Closure (RSCC) declares
\[
F=\mathbb C^3_{\mathrm{perm}}\oplus V_{\mathrm{std}}
\cong\mathbf1\oplus2V_{\mathrm{std}},
\qquad \dim F=5,
\qquad \dim F_0=4,
\]
and the composite response dimensions
\[
(29,432,22,32,840,1008,432,1584).
\]
These integers are arithmetically correct for the declared direct sums and tensor products.  If a
continuous Gaussian source space, full unitary isotropy, reversible independent sheets, the listed
effect ranks, and the listed signs are additionally assumed, the second-cumulant identity gives
\[
\log\mathbb E[e^X]=\frac12\operatorname{Var}X
=\frac{r}{d}w^2.
\]
This conditional identity does not construct the physical effects or select the modules.
In particular,
\(\dim\operatorname{End}_{S_3}(\mathbf1\oplus2V_{\mathrm{std}})=5\), so
\(S_3\)-invariance does not force one scalar covariance on the two-copy multiplicity space.
Pooling heterogeneous sums such as \(\operatorname{End}(F)\oplus F_0\) into a single
\(1/29\) response requires an additional symmetry that mixes physically distinct blocks.
The proposed \(F\) also has heat spectrum \(0^{(1)},3^{(4)}\), not the regular representation's
additional sign-sector value \(6^{(1)}\); the physical regular-heat-to-family attachment is open.
Finally, a genuine Gaussian covariance term is nonnegative.  The negative \(w^2\) entries require a
separate signed response or subtraction law, not orientation reversal alone.

With \(w=\pi\alpha_U\), the declared RSCC formulas are
\[
\begin{aligned}
\bar\sigma_u&=3P+\left(5+\frac4{15}\right)w+\frac{w^2}{29},&
\bar\sigma_d&=2P+\frac4{15}w-\frac{w^2}{432},\\
a_u&=3P+\left(5+\frac45\right)w+\frac{w^2}{22},&
a_d&=2P+\left(1+\frac1{32}\right)w+\frac{w^2}{420},
\end{aligned}
\]
and
\[
\frac{g_{\mathrm{ch}}}{v}
=2\exp\!\left[-2\pi+\frac{P}{1008}
+\frac{w^2}{432}-\frac{w^2}{1584}\right].
\]
Conditional on these formulas and the inherited heat-time, ray, even-response, affine \(A,B\),
and exponentiation laws, the downstream map is deterministic.  Using the D10
\(v\)-display gives a maximum residual of \(0.2943587\%\) and a raw diagonal nominal-residual sum
of \(1.1613966\) against the same mixed-convention comparison table.  Neither number is a
likelihood: the ledger uses visible target-derived effective coordinates,
there is no common-scale RG packet or theory covariance, and the supplied RSCC comparison changes
the nominal \(u\)-row uncertainty relative to the preceding bundle.

Target dependence and ablation fix the status. The target-inferred denominators for four
visible effective coordinates are approximately
\[
29.19665,\qquad428.18579,\qquad21.90604,\qquad836.47734,
\]
close to the selected \(29,432,22,840\), while no target-independent grammar of alternative
module expressions is frozen. The zero-\(w^2\), zero-\(\delta_g\) ablation has the lower
maximum residual \(0.2142293\%\) and raw residual sum \(0.6694007\). This does not make
the ablation a physical theory; it shows that the detailed covariance ledger is not selected by
the numerical agreement. The pixel input depends on an internal quark model, and the
D10 artifact marks the selector candidate-only with mixed sources.

RSCC is an explicit, falsifiable target-conditioned specification for parts of receipts F2--F5
and discharges none of F1--F6. The claimed rescaling-orbit statement must also be scoped correctly:
\[
\|\lambda E_q-E_q\|^2=(\lambda-1)^2\|E_q\|^2
\]
has its minimum at one only on the orbit through the declared RSCC vector.  For a generic base
vector \(C_q\), the minimizer of \(\|\lambda C_q-E_q\|^2\) is
\(\langle C_q,E_q\rangle/\|C_q\|^2\), and no OPH dynamics requires minimizing this
postulated residual. RSCC is a target-conditioned module-ledger ansatz; it does not satisfy the
Flavor Source Closure contract or define a physical quark-mass law.

\paragraph{Further-theorem audit and conditional QFRC rigidity.}
The finite-MaxEnt premise is false. On a finite spectrum, maximizing
entropy under mean and covariance constraints gives a discrete exponential-quadratic Gibbs law,
not a Gaussian density.  Explicitly, on support \(\{-2,-1,0,1,2\}\),
\[
p_A=(1/12,1/6,1/2,1/6,1/12),\qquad
p_B=(1/16,1/4,3/8,1/4,1/16)
\]
both have mean zero and variance one, while
\(\kappa_4(p_A)=0\) and \(\kappa_4(p_B)=-1/2\).  The actual finite-support MaxEnt law at
variance one has a nonzero fourth cumulant (approximately \(-0.46815\)).  Hence neither finite
MaxEnt nor the first two moments entails RSCC's Gaussian two-cumulant truncation.  Gaussianization
can be recovered only conditionally from an exported triangular array with a Lindeberg or adequate
mixing/bounded-dependency condition, normalization, and covariance convergence; no such array is
emitted by the source.

The Quark Flavor Register Closure (QFRC) proposal instead states an exact primitive-path
certificate.  Conditional on QF1--QF9, normalized trace fixes each path weight as rank over the
declared register dimension, including
\[
-\frac15,\frac4{15},\frac1{29},-\frac1{432},\frac45,\frac1{22},
\frac1{32},\frac1{420},-\frac1{10},-\frac14,
\frac1{1008},\frac1{432},-\frac1{1584}.
\]
Exact record projectors also allow only the zero or unit scalar multiplier, and inert tensor
refinement preserves normalized trace weights.  This is a useful rigidity theorem, but its
hypotheses contain the required physical content: the typed register construction, exhaustive
primitive-path catalogue, effect ranks, structural multiplicities, signs, winding character,
refinement intertwiners, implementation invariance, and positive-gap branch selector.  A neutral
finite register can be changed while preserving the broad structural signature and changing the
normalized response, so the present broad axioms do not select QFRC.  That countermodel is scoped
to the broad signature; it is not a realization of every stronger carrier condition one might add.

The selector results sharpen the same boundary.  An equivariant section requires a
stabilizer-fixed candidate in each orbit fiber; an ambiguous transitive fiber need not have one.
Conversely, a complete finite invariant candidate class with a unique positive-gap minimizer has a
robust selector when refinement error is below half the gap.  QFRC supplies neither a physically
complete candidate class nor a source-derived cost and gap.  The absolute-scale rescaling no-go,
finite-renormalization scheme ambiguity, interval-root schema, and piecewise-RG uniqueness theorem
are likewise exact obstruction or conditional well-posedness statements.  No actual target-free
source map with strict interval inclusion, operational clock, or frozen beta/threshold/matching
packet accompanies the proposal.  Therefore QFRC composes algebraically to the frozen RSCC packet
but closes none of F1--F6 and authorizes no numerical quark-mass claim.

\paragraph{Targeted POFT carrier-emission assay.}
The Primitive Oriented Family Transport proposal supplies explicit complex $3\times3$ matrices
$T_0,T_1$, but the simulator test must not manufacture them through a fitted readout.  The frozen
direct observable is therefore the mean of the natural three-label permutation matrices carried by
the saved oriented $S_3$ edges.  Its singular-value ratios are invariant under family permutations,
unitary row/column phases, and an overall scale, so they give a necessary comparison before any
entrywise match is considered.  POFT requires
\[
s(T_0)/s_1=(1,0.55147\ldots,0.25517\ldots),\qquad
s(T_1)/s_1=(1,0.54625\ldots,0.27451\ldots).
\]
A fresh $4{,}096$-patch BW run, an independent $4{,}096$-patch fusion run, a dense
$65{,}536$-patch population run, and a $65{,}536$-patch BW run contain $830{,}066$ edges in
total.  Their corresponding triples are
\[
(1,0.00522,0.00134),\quad(1,0.00552,0.00361),\quad
(1,0.00179,0.00002),\quad(1,0.00071,0.00028).
\]
Source-node block bootstraps retain the same near-rank-one result.  Every state lies within the
predeclared $0.02$ Haar-null tolerance and more than $0.5407$ from both POFT targets under the
predeclared $0.05$ POFT tolerance.  The saved state schema contains only edge endpoints,
integer $S_3$ labels, and geometry: it exports no complex oriented family amplitude and no paired
coarse/fine edge intertwiner.  Hence all direct $T_0$, refined $T_1$, and joint physical-emission
receipts are false.  The conclusion is scoped to the natural direct carrier.  A source-derived
complex lift could define a different carrier.  Choosing its phases or paths to reconstruct the
proposed matrices would make the simulator test circular.

\paragraph{Maximal theorem-emitted quark package theorem.}
Let
\[
\begin{aligned}
P:={}&p_a\\
&+\text{(Axioms 1--5)}\\
&+\text{Assumption 6}\\
&+\text{the emitted OPH nodes D1--D10}\\
&+\text{the listed public D12 quark objects}.
\end{aligned}
\]
Then the quark-side package can be written in four layers:
\begin{enumerate}[leftmargin=2em]
\item the emitted D12 mass ray
\[
D12_{ud}^{\mathrm{mass}}
=
\mathcal R_{D12}^{ud}
=
\left\{
\lambda\left(\frac15,-\frac{1-x_2^2}{27}\right):\lambda\ge0
\right\},
\qquad
\lambda=\mathrm{ray\_modulus}=t_1;
\]
\item the same-label left-handed selector value
\[
\sigma_{ud}=\sigma_{\mathrm{ref}},
\]
with canonical token \(\texttt{D12::same\_label\_left::reference\_sheet}\), which is a negative
sheet-selector statement;
\item the separate target-free mass bridge
\[
P\vdash
\Delta_{ud}^{\mathrm{overlap}}=\frac16\log\frac{c_d}{c_u},
\qquad
P\vdash
\Theta_{ud}^{\mathrm{mass}}=\text{\ophid{quark_same_family_value_law}};
\]
\item the selected-class support wrapper on \(f_P\), which proves representative independence for
the attached physical spread datum. It does not select either modulus. Conditional mass readout
would also require an RG-covariant trajectory and a declared comparison chart; physical Yukawa
matrices require the additional common-scale dimensionless conversion.
\end{enumerate}

Separate same-family and common-refinement artifacts reproduce their chosen target rows exactly.
Their sigma datum is obtained by inversion of those rows. The associated GeV-valued matrices have
the same status: target-anchored mass-texture audits. A continuation sidecar also backreads a
mass-side scalar after identifying coefficient ratios with target mass ratios. None of these
surfaces breaks the free source-spread action.

\subsection{Quark Theorem Boundary}\label{quark-family-boundary}

The conditional downstream closure route is
\[
f_P
\quad+\quad
\Sigma_{ud}^{\mathrm{source}}(P)
\Longrightarrow
\bigl(\sigma_u,\sigma_d,\sigma_{\mathrm{seed}}^{ud},\eta_{ud}\bigr)_{\mathrm{phys}}
\Longrightarrow
\bigl(g_u,g_d\bigr)
\Longrightarrow
\bigl(m_u,m_d,m_s,m_c,m_b,m_t\bigr),
\]
The first arrow is not identified by the source theory. The last arrow denotes scheme-labelled
mass coordinates only after an RG trajectory and comparison chart are specified. Physical
dimensionless Yukawa matrices lie one common-scale normalization beyond it. Selected-fiber
descent and global frame classification are logically separate from both obstructions.

\subsection{Strong-CP branch}\label{quark-family-strong-cp-boundary}

The selected-class quark wrapper carries target-anchored mass textures on the public quark frame
class \(f_P\). Strong CP is a separate phase-side invariant. The available corpus does not derive the
bare QCD angle \(\theta_{\mathrm{QCD}}\), does not emit the physical anomaly-invariant
combination \(\bar\theta\), and does not prove that the physical strong-CP phase vanishes.

This phase problem is independent of the source-spread and common-scale Yukawa obstructions. A
closure would require a source-emitted quark mass matrix at one declared scale, its physical
determinant-line phase contribution, and a theorem fixing the topological-angle contribution on the
realized branch. The present GeV mass textures do not supply that input.

\section{Charged-Lepton Family Derivation}

The charged-lepton derivation differs from the quark derivation in a precise way. It does not emit public charged masses from \(P\). It emits an exact same-family witness, a closed common-shift no-go, the declared same-label \(q_e\) readback, a determinant-line lift on theorem-grade physical charged data, and a downstream algebraic mass readout from theorem-grade \(A_{\mathrm{ch}}(P)\). For any fixed formal source exponent vector \(M_\bullet^{\mathrm{ch}}\), the same-label readback defines a source-side determinant character
\[
S_M=\sum_e M_e^{\mathrm{ch}}\log q_e.
\]
The theorem lane does not emit a theorem-grade sector-isolated charged determinant exponent vector, and it does not identify a source-side determinant character with the physical charged determinant line. The electron, muon, and tau rows are therefore reported as \(n/a\) on the public theorem lane.

The mathematical split is equally precise. The charged-lepton derivation starts from
the ordered charged package, proves that the realized support is a one-dimensional linear subray,
exposes the canonical quadratic support-extension direction, maps that into the charged
excitation gaps, closes a two-scalar support-extension law shell, isolates the smaller eta
source-readback primitive on that same carrier, and then builds the log-spectrum and forward
shape/scale surface. Those centered objects are common-shift invariant. The determinant line fixes the physical affine scalar once theorem-grade physical charged data are present, and the mass readout is then algebraic. The charged theorem boundary has two distinct pieces. One piece is promotion of the latent charged sector-response candidate to theorem-grade \(\widehat C_e\). The other is source-to-physical determinant attachment. For a fixed formal source exponent vector \(M_\bullet^{\mathrm{ch}}\), that attachment is the identity
\[
3\mu(r)=\sum_e M_e^{\mathrm{ch}}\log q_e(r),
\]
equivalently zero determinant-normalization defect
\[
N_{\det}(P)=s_{\det}(P)-\sum_e M_e^{\mathrm{ch}}\log q_e(P),
\]
on the charged determinant channel.

Physically, the electron is the light charged lepton that makes atoms and chemistry possible, the
muon is its heavier unstable cousin, and the tau lepton is the heaviest charged lepton. A promoted
charged-lepton derivation would have to explain how those three rows emerge from the shared flavor
dictionary without Koide-assisted fitting. The theorem-grade lane has to derive promotion of the
latent charged sector-response candidate \(\widehat C_e^{\mathrm{cand}}\) to theorem-grade
\(\widehat C_e\) by closing the branch-generator splitting theorem. On theorem-grade physical
\(Y_e\), a refinement-stable uncentered lift collapses the determinant-line section and affine
anchor to one descended physical affine scalar \(\mu_{\mathrm{phys}}(Y_e)\), with
\[
\widetilde C_e(Y_e)=\widehat C_e(Y_e)+\mu_{\mathrm{phys}}(Y_e)\,\mathbf 1,
\qquad
s_{\det}(Y_e)=3\mu_{\mathrm{phys}}(Y_e),
\qquad
A_{\mathrm{ch}}(Y_e)=\mu_{\mathrm{phys}}(Y_e).
\]

\begin{theorem}[Charged same-carrier source-pair readback]
Fix the ordered charged carrier
\[
\begin{gathered}
(-1,x_2,1),\\
x_2=-0.5175863354681689.
\end{gathered}
\]
If the charged source pair
\[
(\eta_{\mathrm{ext}},\sigma_{\mathrm{ext}})
=
(\eta_{\mathrm{source\_support\_extension\_log\_per\_side}},
\sigma_{\mathrm{source\_support\_extension\_total\_log\_per\_side}})
\]
is emitted on that carrier, then the centered charged logs are
\[
e_{\log,\mathrm{centered}}
=
-\frac{(3+x_2)\sigma_{\mathrm{ext}}-\eta_{\mathrm{ext}}}{6},
\]
\[
\mu_{\log,\mathrm{centered}}
=
\frac{x_2\sigma_{\mathrm{ext}}-\eta_{\mathrm{ext}}}{3},
\]
\[
\tau_{\log,\mathrm{centered}}
=
\frac{(3-x_2)\sigma_{\mathrm{ext}}+\eta_{\mathrm{ext}}}{6},
\]
and therefore the charged masses are
\[
m_e = g_e\,e^{e_{\log,\mathrm{centered}}},
\qquad
m_\mu = g_e\,e^{\mu_{\log,\mathrm{centered}}},
\qquad
m_\tau = g_e\,e^{\tau_{\log,\mathrm{centered}}}.
\]
\end{theorem}

This theorem is exact on the same-carrier shell, but the absolute values are not emitted.
The visible scalar order is
\[
\eta_{\mathrm{ext}}
\quad\text{then}\quad
\sigma_{\mathrm{ext}},
\]
and the charged absolute scale \(g_e\) is unresolved. Builder-facing code therefore
exposes \(\eta_{\mathrm{ext}}\) and then \(\sigma_{\mathrm{ext}}\) as the first same-carrier
residuals, but the theorem-grade boundary lies above that surface. If the latent
candidate \(\widehat C_e^{\mathrm{cand}}\) is promoted, then \(\eta_{\mathrm{ext}}\) and
\(\sigma_{\mathrm{ext}}\) become charged spectral invariants instead of separate primitive goals,
and the absolute-scale burden is pushed to one affine-covariant absolute charged anchor
\(A_{\mathrm{ch}}\). In the
local chain, \(\widehat C_e\) itself is undeclared: only the centered compressed
generation-bundle branch-operator candidate \(\widehat C_e^{\mathrm{cand}}\) is on disk, and the
operator-side gate for that package is the upstream promotion theorem
\ophid{oph_generation_bundle_branch_generator_splitting}, not a new ad hoc charged
operator ansatz. Its exact gate clause is the compression-descendant commutator statement
\ophid{compression_descendant_commutator_vanishes_or_is_uniformly_quadratic_small_after_central_split}.
The centered common-shift quotient is closed negatively, so centered data alone do not
emit the affine anchor \(A_{\mathrm{ch}}\).

\begin{theorem}[Charged absolute-scale underdetermination]
Let
\[
E_e^{\mathrm{centered}}
=
\bigl(e_{\log,\mathrm{centered}},\mu_{\log,\mathrm{centered}},\tau_{\log,\mathrm{centered}}\bigr)
\]
be the centered charged log triple emitted from the charged source pair
\((\eta_{\mathrm{ext}},\sigma_{\mathrm{ext}})\). Then for every \(c\in\mathbb R\),
\[
Y_e(c):=\exp(c)\,\mathrm{diag}\!\bigl(e^{E_e^{\mathrm{centered}}}\bigr)
\]
has the same charged spectral invariants
\[
\eta_{\mathrm{ext}},\qquad \sigma_{\mathrm{ext}},\qquad \gamma_{21},\qquad \gamma_{32},
\]
the same centered log vector, and the same ratio data. Only the absolute masses scale:
\[
(m_e,m_\mu,m_\tau)\longmapsto e^c\,(m_e,m_\mu,m_\tau).
\]
Hence the charged OPH chain determines only the quotient class
\[
E_e^{\mathrm{centered}}\in \mathbb R^3/\langle(1,1,1)\rangle,
\]
not the absolute scalar \(g_e=e^c\).
\end{theorem}

The proof is immediate from the centered sum rule
\[
e_{\log,\mathrm{centered}}+\mu_{\log,\mathrm{centered}}+\tau_{\log,\mathrm{centered}}=0,
\]
which implies
\[
\det Y_e^{\mathrm{shape}}=1,
\qquad
\det Y_e = g_e^3.
\]
All emitted charged invariants depend only on differences of log entries, so a common shift in the
\((1,1,1)\) direction leaves them unchanged. This is the exact reason the available corpus closes the
\(P\)-driven charged lane as a no-go instead of as a mass theorem: the theorem surface fixes the
centered shape, while the determinant-line landing fixes the absolute normalization.

The charged absolute-scale
lane is explicitly typed. The charged scale is a linear quantity
\[
g_e = e^{\mu_e^{\mathrm{abs}}},
\]
so the log-coordinate \(\mu_e^{\mathrm{abs}}\) must not be mixed directly with centered log gaps.
The same-family writeback therefore records the type-consistent shell
\[
\begin{aligned}
\mu_{e,\mathrm{seed}}^{\mathrm{abs}}
&= \log(0.9231656602589082)\\
&= -0.07994658034676537,
\end{aligned}
\]
\[
\mu_{e,\mathrm{cand}}^{\mathrm{abs}}
=
\mu_{e,\mathrm{seed}}^{\mathrm{abs}}-\gamma_{\min}
=
-0.38231224060567365,
\]
\[
g_e^{\mathrm{cand}}
=
e^{\mu_{e,\mathrm{cand}}^{\mathrm{abs}}}
=
0.6822819838027987.
\]
This is a representation-consistency shell only, not a charged-mass theorem. It fixes the
linear-vs-log coordinate discipline. This surface does not emit a theorem-grade value law for \(g_e\).
The audited shortcut
\[
\begin{aligned}
\Delta_e^{\mathrm{abs}}&=0.30236566025890826,\\
g_e&=0.6822819838027987
\end{aligned}
\]
is not a theorem-grade closure. It merely chooses one representative on the common-shift orbit and
does not land on the physical charged masses.

A physical completion requires a theorem-grade affine-covariant section
\(A_{\mathrm{ch}}\) of the quotient map, satisfying
\[
A_{\mathrm{ch}}(\ell + c\,\mathbf 1)=A_{\mathrm{ch}}(\ell)+c.
\]
The clean realization of that section is an uncentered charged response lift carrying a
determinant line, in which
\[
A_{\mathrm{ch}}=\tfrac13 \log \det(Y_e)=\tfrac13 \operatorname{tr}(\log Y_e).
\]
Once such an anchor exists,
\[
g_e = e^{A_{\mathrm{ch}}(\ell)},
\qquad
\Delta_e^{\mathrm{abs}} = \log g_{\mathrm{ch}}^{\mathrm{shared}} - A_{\mathrm{ch}}(\ell).
\]

\paragraph{Compare-only charged continuation bridge.}
For comparison, there is a sharper charged continuation bridge under three extra continuation
assumptions:
\begin{enumerate}[leftmargin=2em]
\item a uniform \(\mathbb Z_6\) center-label ensemble, so \(\varepsilon = 1/6\);
\item the balanced positive Hermitian circulant branch, so \(\rho/a = 1/\sqrt2\);
\item the phenomenological charged-family phase choice \(\delta = 2/9\), the fitted phase of the
Brannen parametrization of the Koide relation~\cite{koide1983,brannen2006}.
\end{enumerate}
These are continuation inputs, not consequences of the OPH axioms.  For the positive-semidefinite
Hermitian square-root-mass carrier
\[
C=aI+\rho\left(e^{i\delta}R+e^{-i\delta}R^2\right),
\qquad a,\rho\geq0,
\]
where \(R\) is the regular cyclic shift, \(R^3=I\).  The commutant of the regular \(C_3\)
action is spanned by \(I,R,R^2\), which makes \(C\) the general Hermitian
\(C_3\)-equivariant carrier. Fourier diagonalization gives
\[
r_k=a+2\rho\cos\!\left(\delta+\frac{2\pi k}{3}\right),
\qquad k=0,1,2.
\]
In a positive spectral chamber, \(m_k=s r_k^2\) and
\(\sqrt{m_k}=\sqrt{s}\,r_k\).  The roots-of-unity identities
\[
\sum_{k=0}^2\cos\!\left(\delta+\frac{2\pi k}{3}\right)=0,
\qquad
\sum_{k=0}^2\cos^2\!\left(\delta+\frac{2\pi k}{3}\right)=\frac32
\]
give the Koide invariant
\[
Q=\frac{\sum_km_k}{(\sum_k\sqrt{m_k})^2}
 =\frac{1+2(\rho/a)^2}{3},
\]
so \(\rho/a=1/\sqrt2\) is exactly equivalent to \(Q=2/3\).  At balance, positivity holds for
\(|\delta|\leq\pi/12\pmod{2\pi/3}\), and throughout that chamber Koide is independent of
\(\delta\). Outside it, physical square roots are absolute eigenvalues and the physical Koide
ratio differs from the signed trace expression. Equivalently, if
\(E_0=3a^2\) and \(E_c=6\rho^2\) are the singlet and charged-plane Hilbert--Schmidt powers, then
\[
Q=\frac{1+E_c/E_0}{3},
\qquad
Q=\frac23\ \Longleftrightarrow\ E_c=E_0.
\]
The positive square-root-mass vector lies on the \(45^\circ\) Koide cone about the democratic
direction. Circulant symmetry leaves the relative norm of the singlet and charged plane free.

The finite OPH response packet sets that norm. Let
\[
\mathcal V_K=\mathbf1\oplus\chi\oplus\bar\chi,
\qquad \mathcal H_{\rm or}=\mathbb C^2,
\qquad
\mathcal A_K=B(\mathcal V_K\otimes\mathcal H_{\rm or})\simeq M_6(\mathbb C).
\]
A one-state record cannot distinguish orientation reversal, so \(\mathbb C^2\) is the minimal
faithful orientation record. Application of Minimal Admissible Realization (MAR) to this
response-local register is an explicit hypothesis. With \(q_+\) one oriented outcome, define
\[
E_+=P_0\otimes I_{\rm or}+P_c\otimes q_+,
\qquad
Z_0=P_0\otimes I_{\rm or},
\qquad
Z_c=P_c\otimes q_+.
\]
The blocks \(Z_0,Z_c\) both have rank two. Born--L\"uders conditioning of the unique MaxEnt
state \(I_6/6\) on \(E_+\) gives
\[
p_0=p_c=\frac12.
\]
The orientation-blind singlet has event action zero, while the selected charged orientation has
action \(\ln2\). Charged multiplicity two times its weight \(1/2\) equals the singlet weight.
The unique normalized trace on \(M_6(\mathbb C)\) fixes these probabilities without the free trace
parameter present on \(\mathbb C\oplus M_2(\mathbb C)\).

In \(L^2(B(\mathbb C^3),\tau_3)\), the operators \(I,R,R^2\) are orthonormal, and
\[
\mathcal J(e_0)=I,
\qquad
\mathcal J(e_+)=R,
\qquad
\mathcal J(e_-)=R^2
\]
is unitary and intertwines the \(C_3\) action, orientation reversal, and the two block projections.
The canonical GNS square-root amplitude is
\[
\xi=\sqrt{p_0}\,e_0+\sqrt{\frac{p_c}{2}}
\left(e^{i\delta}e_++e^{-i\delta}e_-\right),
\]
so its response has \(a=\sqrt{p_0}\) and \(|b|=\sqrt{p_c/2}\).  Hence
\[
\frac{|b|^2}{a^2}=\frac{p_c}{2p_0}=\frac12,
\qquad
\frac{|b|}{a}=\frac1{\sqrt2},
\qquad Q=\frac23.
\]
An orientation action gap \(\Delta S\) gives the general form
\[
\frac{p_c}{p_0}=2e^{-\Delta S},
\qquad
\frac{|b|^2}{a^2}=e^{-\Delta S},
\qquad
Q(\Delta S)=\frac{1+2e^{-\Delta S}}{3}.
\]
The physical OPH Koide theorem is conditional on faithful trace-preserving u.c.p. maps \(\Phi\)
and \(\Psi\) from the source process to a physical response process and back, with
\(\Psi\Phi=\mathrm{id}\). Kadison--Schwarz applied to both maps puts every source element in the
multiplicative domain of \(\Phi\), so \(\Phi\) is a \(*\)-monomorphism and an exact \(L^2\)
isometry. If it also intertwines \(Z_0,Z_c\) with physical response records and sends the
conditioned source state to the normalized positive square-root-mass response
\(C=(Y_e^\dagger Y_e)^{1/4}\), then it transports \(p_0=p_c\) to \(E_0=E_c\) and proves physical
Koide. For an accepted finite chiral checkpoint with three left and three right family modes,
positive kinetic metrics, and a committed neutral Higgs direction, the reversible-checkpoint
conditions force \(\Phi=\operatorname{Ad}_{J_L\oplus J_E}\), preserving chirality and the
regular \(C_3\) action. If \(M_F=X_F^2\) is the positive source response, then
\[
\widehat Y_e=\frac{\sqrt2}{v}J_LM_FJ_E^\dagger,
\qquad
\mathcal M_L=J_LM_FJ_L^\dagger.
\]
The source and canonical physical square-root responses have equal block powers. This closes the
finite source-scale mass response. The extended branch \(\mathrm{OPH}^{+}_{\rm ch}\) adds graded
physical completion and quotient source-law selection. Given an exhaustive MAR carrier class with
a positive winner gap, a source-closed BV/BRST continuum, and an interval or contraction
certificate for the charged QFT self-map, it selects one accepted charged sector and one dressed
mass readout. The stable electron mass is a dressed spectral lower edge; the muon and tau masses
are real parts of specified dressed resonance roots. A balanced \(C_3\) fixed point gives
\(Q=2/3\). A \(C_3\)-symmetric fixed point with charged attenuation \(\chi_\star\) gives
\[
Q=\frac{1+e^{-2\chi_\star}}{3}.
\]
An off-plane response requires the full response operator. These extended-branch principles and
the QFT certificate are additional conditions beyond OPH5.
The relation \(\operatorname{Hom}_{C_3}(\mathbf1,\chi\oplus\bar\chi)=0\) prevents the neutral
source grammar from producing a labeled charged-family vector through symmetry. A charged source
tensor, connection, or selected orbit must supply the phase and individual ratios.

The count \((N_c+1)/(2N_cN_g)=2/9\) is arithmetic, and the value \(2/9\) is the
empirical Brannen--Koide fitted phase~\cite{koide1983,brannen2006}; the integer
factorization is conditioned on that value. A physical phase theorem requires a
regular-\(C_3\) family bundle, an oriented family connection
whose declared loop has holonomy
\(\exp[i\beta_{\rm EW}/(2N_cN_g)]\), and an attachment identifying that holonomy with the phase of
the charged Fourier coefficient. Hypercharge acts as a common scalar on the generation copies,
while the shift \(R\) permutes them. Pure finite \(A_5/C_3\) holonomy yields cube-root phases.
Construction of the required continuous family connection and its source selector is work in
progress. Under the stated connection and attachment hypotheses, \(\delta=2/9\). The phase of one
equal link differs by a factor of three from the total triangular holonomy.
With the scale-free normalization \(a=1\), the ordered roots
\[
r_k = a + 2\rho\cos\!\left(\delta + \frac{2\pi k}{3}\right)
\]
become
\[
\begin{aligned}
r_e&=0.040349908219207475,\\
r_\mu&=0.5802119201475368,\\
r_\tau&=2.3794381716332555.
\end{aligned}
\]
Under those assumptions the bridge selects the same-carrier pair
\[
\begin{aligned}
\eta_{\mathrm{ext}}&=-6.729586682888832,\\
\sigma_{\mathrm{ext}}&=8.154061112725994,
\end{aligned}
\]
with ordered gap values
\[
\begin{aligned}
\gamma_{21}&=5.33160859254774,\\
\gamma_{32}&=2.822452520178255,\\
\kappa_{\mathrm{ext}}&=-4.59605680397234,
\end{aligned}
\]
and centered logs
\[
E_{\log,\mathrm{centered}}
=
[-4.495223235091244,\ 0.836385357456495,\ 3.6588378776347503].
\]
Against the charged references, this continuation-centered shape has residual norm
\[
\left\|E_{\log,\mathrm{centered}}^{\mathrm{cont}}
-E_{\log,\mathrm{centered}}^{\mathrm{ref}}\right\|
\approx
2.13\times10^{-5},
\]
so this compare-only branch is a near-exact centered-shape closure up to one common absolute
scale. The ppm-scale shape agreement carries a large pull in measurement units: the same
\(\delta=2/9\) ansatz misses the measured \(m_\mu/m_e\) ratio by \(9.83\)~ppm, approximately
\(440\sigma\) of the measurement precision, so the row is excluded as a prediction and stands as
an approximate fit. The theorem lane is unpromoted because the public affine absolute anchor is external
to this near-exact centered charged-shape branch. The compare-only common scale
required for exact absolute masses is
\[
g_e^\star = 0.04577885783568762.
\]
Equivalently, relative to the stored shared-budget seed
\[
g_{\mathrm{ch}}^{\mathrm{shared}} = 0.9231656602589082,
\]
the missing affine absolute anchor on the determinant-line route would have to take the target value
\[
\Delta_e^{\mathrm{abs},\star}
=
\log\!\frac{g_{\mathrm{ch}}^{\mathrm{shared}}}{g_e^\star}
=
3.003986333402356.
\]
This identifies the public charged theorem boundary: the compare-only branch nearly solves
the centered charged shape, but the theorem-grade derivation lacks promotion of the latent
candidate \(\widehat C_e^{\mathrm{cand}}\) to theorem-grade \(\widehat C_e\), and then one
affine-covariant absolute anchor \(A_{\mathrm{ch}}\) that would turn that centered readback into
public charged masses on the theorem lane.

\paragraph{Icosahedral face-incidence carrier.}
The screen geometry supplies a sharper candidate carrier than a scalar
twelve-port moment.  The twenty outward-oriented icosahedral faces form
\(A_5/C_3\), and the face stabilizer cyclically permutes the three corners.
Thus each face has a canonical local regular-\(C_3\) fiber, and an equivariant
Hermitian circulant has a face-representative-independent unordered spectrum.
This is an exact geometric lemma.  It is not a physical generation
attachment: the sixty face-corner flags form the regular \(A_5\) torsor, not
one canonical three-dimensional matter-family space.

A conditional theorem package declares the diagonal affine repair map
\[
T(\kappa,\chi_\rho,\zeta_\delta)
=
(s_\kappa-q_\kappa\kappa,\,
  s_\chi+q_\chi\chi_\rho,\,
  s_\zeta+q_\zeta\zeta_\delta).
\]
For that declared map it proves
\(\lVert DT\rVert_\infty=0.0015356510519\ldots<1\), so Banach's theorem closes
existence, uniqueness, and iterative stability of its fixed point.  This is a
conditional closure, not a derivation of the repair dynamics.  A
source-multiplier family \(T_{\boldsymbol\lambda}\) with the
same displayed symmetry, analytic degree, Jacobian, and contraction but
different masses exists.  It proves scoped selector non-identifiability
under those properties, not non-entailment from every OPH axiom.  Exact
zeroth-order block balance gives
\(|b|/a=1/\sqrt2\), whereas the endpoint operator uses
\(|b|/a=e^{-\chi_\rho}/\sqrt2\); a source-derived bare-to-endpoint repair bridge
is work in progress.

\begin{proposition}[Finite CFQ schema satisfiability]
\label{prop:charged-cfq-digital-model}
The stipulated CFQ register-and-path class is nonempty.  There is an explicit
finite algebraic model with register dimensions
\[
(50,31,10,512,77,21,27,5),
\]
connected transition graphs, normalized tracial states, rank-one event
projectors, and the eight declared path weights
\[
\left(\frac1{50},-\frac1{31},-\frac1{310},\frac1{512},
\frac1{77},\frac1{21},\frac1{27},\frac1{135}\right).
\]
Each noncentral event in \(B(H_r)\) admits an accepted/rejected central record
dilation in \(B(H_r)\otimes D_2\).  The graph family is covariant under the
sixty proper icosahedral rotations, and normalized path weights are unchanged
under inert ancillary stabilization \(X\mapsto X\otimes I_k\).
\end{proposition}

\begin{proof}
For each declared connected graph, the diagonal matrix units and both directed
units on every edge generate \(E_{ij}\) along graph paths and hence generate
the full matrix algebra.  The normalized trace of a rank-one event is the
reciprocal register dimension, and tensor products multiply the two depth-two
weights.  For an event \(P\), the pinching map
\[
\mathcal E_P(X)=PXP+(I-P)X(I-P)
\]
is a unital trace-preserving conditional expectation.  Its instrument writes
\(P\rho P\) and \((I-P)\rho(I-P)\) into orthogonal accepted/rejected pointers in
the central \(D_2\) factor.  Explicit permutation intertwiners preserve graph
incidence and ranks for all sixty proper rotations.  Finally,
\(\tau_{dk}(P\otimes I_k)=\tau_d(P)\), and the partial-trace coarse map retracts
the inert embedding.  These facts construct one model of the declared schema.
\end{proof}

Proposition~\ref{prop:charged-cfq-digital-model} proves formal nonemptiness.
It supplies no derivation of the charged source.  The eight dimensions, path
automaton, coupling degrees, grading, signs, clock, and scalar response are
model inputs.  The verifier exhausts that automaton, not the physically
admissible OPH path space.  The construction resolves the abstract
event-versus-central-record issue at fixed cutoff.  Physical response
selection, cofinal screen refinement, and charged-sector attachment are work
in progress.

\begin{proposition}[Conditional charged nature and singularity transport]
\label{prop:charged-nature-pole-transport}
Let \(M_F\) be the positive face endpoint.  Suppose a physical chiral
three-family carrier supplies a natural unitary \(J_L\), canonical kinetic
metrics, and
\[
 \frac{v^2}{2}\widehat Y_e\widehat Y_e^\dagger
 =J_LM_F^2J_L^\dagger.                                  \tag{NI6}
\]
Then its positive left charged response is \(J_LM_FJ_L^\dagger\).  Suppose in
addition that the exact renormalized charged kernel is supplied, its singular
lines satisfy the declared infrared and regularity conditions, and a
CFQ--Dyson map has singularity readout \(J_LM_FJ_L^\dagger\).  Regular
invertible field changes preserve the zero set, and a regular Nielsen
factorization \(\partial_\xi K=A_\xi K+KB_\xi\) makes each simple zero gauge
independent.
\end{proposition}

\begin{proof}
The first statement follows from uniqueness of the positive square root.
Multiplication of \(K\) on the left and right by analytic invertible matrices
multiplies its determinant by nonvanishing factors.  Jacobi's identity gives
\[
 \partial_\xi\det K=(\operatorname{tr}A_\xi+
 \operatorname{tr}B_\xi)\det K,
\]
so a simple zero cannot move while the Nielsen factors are regular.  These are
the standard mixed-fermion and gauge-independence implications used in the
pole literature~\cite{kniehl2013,gambino1999}; the stable charged-line
infrared qualification follows the infraparticle boundary
\cite{buchholz1986,duch2023}.
\end{proof}

Proposition~\ref{prop:charged-nature-pole-transport} closes the downstream
logic after its physical premises are supplied.  It does not supply them.
NI6 is the desired operator attachment squared, and RP4
premise sets the Dyson singularity readout equal to the face response.
The explicit kernel \(K_0(s)=sI-M_F^2\) has zero self-energy; it is a
free algebraic existence witness rather than the interacting charged-lepton
1PI kernel.  The physical NI1--NI8 and RP1--RP8 receipts are work in progress.
The verifier reconstructs the fixed point and response shape but does not check
\(M_F=g_{\rm end}S_F\): a coherent change
of the ordered spectrum by factors \((2,1/2,1)\), together with the matching
mass matrix, Yukawa matrix, and pole roots, preserves the determinant and
satisfies the tested mass relations.  The package is therefore a conditional
gate theorem, not physical nature identification, an interacting pole theorem,
or a promoted mass prediction.

Public promotion is false.  It requires a quotient-visible face-to-charged-family
intertwiner, the \(\ln2\) MaxEnt and ensemble-to-Yukawa bridge, a charged
connection deriving both the \(2/9\) base phase and its correction, a normalized
\(\mathbb Z_6\) determinant character, an OPH dynamical theorem selecting the
physical CFQ registers, state, exhaustive path category, grading, clock, and
response rather than declaring them, cofinal refinement naturality,
one receipt-bound coherent source tuple frozen before comparison, and a mass-
scheme map.  The exact receipts are
\texttt{charged\_icosahedral\_face\_carrier\_frontier.json} and
\texttt{charged\_face\_incidence\_conditional\_theorem.json}, together with the
conditional CFQ weight receipt; the compare-only
checks are \texttt{charged\_icosahedral\_completion\_v1\_review.json} and
\texttt{charged\_face\_incidence\_conditional\_theorem\_review.json}, the CFQ
rigidity check, and \texttt{charged\_cfq\_carrier\_v1\_review.json}.  The
candidate does not change the theorem mass rows.

\subsection{Absolute charged-lepton mass intervals from electromagnetic transport}\label{charged-kappa-interval}

The missing absolute anchor is one real scale. Fixing \(A_{\mathrm{ch}}\) is equivalent to fixing
the single family rescaling \(Y_e\mapsto e^{\kappa}Y_e\), under which every charged mass scales as
\(m_i\mapsto e^{\kappa}m_i\) while the centered shape and every mass ratio are invariant. The
declared charged antecedents are invariant under that rescaling, which is the content of the
common-shift no-go: the gauge representations, the anomaly structure, the complexity vector, the
determinant algebra, and the same-label readback take the same value on the whole family
\(\{e^{\kappa}Y_e\}\), so no source object built from those antecedents alone selects \(\kappa\).

The rescaling is not a symmetry of the declared electromagnetic transport. The charged leptons
enter the photon vacuum polarization, and the leptonic packet in the fine-structure endpoint map
carries
\[
P_{\mathrm{lep}}(\kappa)=\frac{1}{3\pi}\sum_{i}\left[2\log\frac{M_Z}{m_i}-\frac53\right],
\qquad
\frac{\mathrm{d}P_{\mathrm{lep}}}{\mathrm{d}\kappa}=-\frac{2}{\pi}.
\]
The declared endpoint pixel map \(P=\varphi+\sqrt\pi/A_{\mathrm{Th}}(P)\) consumes \(P_{\mathrm{lep}}\),
so the emitted endpoint and the pixel move with \(\kappa\). The common-shift no-go omits this
transport from its antecedent list. Source-only, the system is the stiff curve \((P(\kappa),\kappa)\)
with \(|\mathrm{d}P/\mathrm{d}\kappa|\approx6\times10^{-5}\), and the absolute masses are work in
progress on the source branch.

On the empirical-closure surface the transport identifies \(\kappa\). Inverting the on-shell
decomposition
\[
a_0+g=A^{-1}(0)\left(1-P_{\mathrm{lep}}(\kappa)-\Delta_{\mathrm{had}}-\Delta_{\mathrm{top}}\right)
\]
with the frozen anchor \(a_0=128.308\), the certified anchor-gap interval \(g\in[0.649,\,0.855]\),
the empirical hadronic payload \(\Delta_{\mathrm{had}}=0.02676\pm0.00075\), the frozen
continuation mass ratios,
and the measured \(A^{-1}(0)\) supplied as a compare-only exclusion inside the interval solve path
gives
\[
\kappa\in[-0.346,\,+0.337],\qquad \kappa_{\mathrm{central}}=-0.004.
\]
The charged masses then carry certified intervals,
\[
\begin{aligned}
m_e&=0.5089~[0.362,\,0.716]~\mathrm{MeV}, & &\text{measured } 0.5110~\mathrm{MeV},\\
m_\mu&=105.22~[74.8,\,148.0]~\mathrm{MeV}, & &\text{measured } 105.66~\mathrm{MeV},\\
m_\tau&=1.7695~[1.258,\,2.489]~\mathrm{GeV}, & &\text{measured } 1.7769~\mathrm{GeV},
\end{aligned}
\]
with central values within about \(0.4\%\) of measurement and the physical triple inside every
interval. The interval width reduces to the empirical hadronic payload uncertainty and the anchor
gap, which are the two open objects of the fine-structure lane.

The determinant prescription is a consistency test on this scale. It reproduces
the physical leptonic packet at \(|\kappa|\approx0.006\), so it is a tested
ansatz on the absolute scale rather than a free parameter. The
construction adds no axiom. The interval lane carries absolute masses only on the
empirical-closure surface; the source-only charged no-go is in force.

\subsection{Conditional response candidate and symmetry-breaking boundary}\label{charged-mcpr-frontier}

A conditional package defines the charged response by
stipulating a typed, incidence-complete response architecture on an oriented
icosahedral flag: eight finite registers of dimensions
\(50,31,10,512,77,21,27,5\), eight primitive path classes with normalized
rank-one trace weights, a public-block amplitude, an oriented process phase,
and a \(\mathbb Z_6\) determinant character.  Conditional on that
architecture, the three-channel response map is a Banach contraction with
contraction constant \(1.5\times10^{-3}\); its unique fixed point, the
regular-\(C_3\) shape functional, and the determinant character emit one
dimensionless charged triple on the strict source-audit branch with zero
runtime charged reference input:
\[
\frac{(m_e,m_\mu,m_\tau)}{E_\star}
=(4.18511\times10^{-23},\,8.65348\times10^{-21},\,1.45532\times10^{-19}),
\]
with mass ratios \(m_\mu/m_e=206.7683\), \(m_\tau/m_e=3477.3655\),
\(m_\tau/m_\mu=16.8177\), and unclosed-clock displays
\((0.51096,\,105.649,\,1776.78)\)~MeV.  A segregated compare-only audit
places this coherent branch about \(84\) ppm below the comparison
triple; the candidate computation has no dependency on the comparison file
(\texttt{charged\_mcpr\_completion\_conditional.json}).

The register dimensions, path table, signs, amplitude, phase, and determinant
exponent are declared model inputs, so the triple is a declared-model
coordinate. The Koide subsection derives the finite
public-block amplitude \(1/\sqrt2\) from the connected \(M_6\) response register.
The candidate supplies no recoverable attachment between that event and a physical chiral mass
channel, so its mass claim is conditional. The arithmetic identity
\((N_c+1)/(2N_cN_g)=2/9\) supplies no oriented phase; a holonomy
reading requires link variables, a loop, and a nonlocal readout map.
A product of fourteen normalized rank-one traces is neither a
\(\mathbb Z_6\) character nor a determinant, and the canonical determinant
norm carries the kinetic factor
\(\det\mathcal M_L=(v/\sqrt2)^3\,|\det Y_e|/\sqrt{\det Z_L\det Z_E}\); the
declared \(6^{-14}\) weight is a candidate positive path weight.

The source-side problem is a finite
symmetry-breaking program on the icosahedral carrier.  The vertex permutation
module decomposes as
\(\mathbb R^{12}\cong W_1\oplus W_3\oplus W_3'\oplus W_5\); the traceless
quadrupole map \(Q(w)=\sum_iw_i(p_ip_i^T-I/3)\) satisfies \(Q=QP_5\) and
\(Q^*Q=\frac85P_5\), so it sees exactly \(W_5\), and the two adjacent gaps of
its spectrum span the centered family plane.  A unique invariant MaxEnt state
has zero expectation on every nontrivial irreducible module, and the
homogeneous twelve-port branch has the unique uniform minimizer, so charged
family shape requires a source-selected \(W_5\) orbit.  Because
\(\operatorname{Sym}^2_0(\mathbf3)\cong\operatorname{Sym}^2_0(\mathbf3')\cong W_5\),
the family attachment is multiplicity-one up to one sign once the chiral
family fibers carry a three-dimensional \(A_5\) representation.  Work in
progress comprises an \(A_5\)-invariant effective action on \(W_5\) with a unique
simple-spectrum minimizing orbit (the invariant algebra has one quadratic,
two cubic, and two quartic coefficients, so the first decisive quantity is
the \(W_5\)-restricted Hessian scalar \(h_5\) at the homogeneous state), the
physical \(A_5\) family lift, the normed determinant-line descent with
kinetic factors, the interacting charged kernel packet, the operational
scale, and a precomparison source receipt.

The affine scale has a second admissible closure through electromagnetic
transport.  With the exact one-loop Ward kernel
\(I(z)=\int_0^1x(1-x)\log(1+zx(1-x))\,\mathrm dx\) in closed form, the
charged response
\(\mathcal W_Q(\mu;\ell)=\frac2\pi\sum_iI(q^2e^{-2(\mu+\ell_i)})\) is a
smooth, strictly decreasing bijection of the common log-scale \(\mu\) onto
\((0,\infty)\) at fixed centered shape \(\ell\), with high-energy slope
\(-2/\pi\).  A source-complete spacelike Ward endpoint pair with a
source-complete nonleptonic subtraction therefore fixes a unique
\(\mu_{\mathrm{ch}}\), hence the determinant line
\(|\det(M_e/E_\star)|=e^{3\mu_{\mathrm{ch}}}\) and the electron ratio
\(m_ec^2/E_\star\) consumed by the cesium clock packet.  The common-shift
no-go of the preceding subsections holds on the reduct without
mass-dependent electromagnetic transport; a Ward endpoint response and a
determinant-line basepoint are the two admissible closures of that orbit.
The Ward inversion, its interval enclosure, and its fail-closed source-packet
schema are encoded in \texttt{charged\_ward\_determinant\_line.json}; the
endpoint pair, the nonleptonic subtraction, the higher-order monotonicity
certificate, and the atomic transport factor are its open parents.

\section{Neutrino Family Derivation}

The neutrino lane contains one exact no-go and one failed comparison candidate. The intrinsic
isotropic branch is excluded by its spectral cap. The weighted-cycle construction descends from
two hand-written family-transport matrices and declared label, orientation, cycle, and exponent
choices selected with target information. Its bridge and absolute attachment are compare-only and
carry no prediction status.

For the isotropic matrix \(M=aI+\rho C\) with unit-modulus off-diagonal entries, the general
Gershgorin estimate applied to \(M^\dagger M\) gives
\[
\max_{i,j}|\Delta m_{ij}^2|\leq 8a\rho+4\rho^2.
\]
The certified bound is \(1.52304\times10^{-6}\,\mathrm{eV}^2\), about three orders of magnitude
below the atmospheric scale.

The weighted-cycle comparison candidate gives
\[
\begin{aligned}
\theta_{12}&=34.225904631810025^\circ,\\
\theta_{23}&=49.72282845058266^\circ,\\
\theta_{13}&=8.686355527700156^\circ,
\end{aligned}
\]
\[
\begin{aligned}
\delta_{\mathrm{PMNS}}&=305.58061231449796^\circ,\\
J&=-0.02753115613565372,
\end{aligned}
\]
and the dimensionless hierarchy invariant
\[
\begin{aligned}
\frac{\Delta m_{21}^2}{\Delta m_{32}^2}
&=0.030721110097966534.
\end{aligned}
\]
The official NuFIT 6.1 normal-ordering profile gives
\[
\Delta\chi^2_{23,\delta}=20.11955
\quad\text{with the tabulated atmospheric likelihood},
\qquad
\Delta\chi^2_{23,\delta}=18.43528
\quad\text{without it},
\]
at \((\sin^2\theta_{23},\delta_{\mathrm{CP}})=(0.582056,-54.419^\circ)\)~\cite{nufit61}.
Both values exceed the two-parameter \(3\sigma\) contour value \(11.83\). Separate marginal
intervals do not test this correlation. The published profiles overlap and are not summed;
their maximum is a lower bound on the unavailable full fixed-candidate \(\Delta\chi^2\).

The declared shared-basis construction sets
\[
M_{\mathrm{shared}}=U_{e,\mathrm{left}}^*\,M_{\mathrm{wc}}\,U_{e,\mathrm{left}}^\dagger,
\qquad
U_{\nu,\mathrm{shared}}=U_{e,\mathrm{left}}\,U_{\mathrm{wc}},
\]
the identity
\[
U_{e,\mathrm{left}}^\dagger U_{\nu,\mathrm{shared}}=U_{\mathrm{wc}}.
\]
This cancellation is a tautology because \(U_{\nu,\mathrm{shared}}\) was defined as
\(U_{e,\mathrm{left}}U_{\mathrm{wc}}\). It does not independently identify the physical
charged-lepton basis or validate the PMNS point. No theorem in the theorem stack places the
\(f\)-labelled weighted-cycle operator in the charged-lepton mass basis.

The charged-basis artifact has
\texttt{closure\_state: open}, and its three normalized shape singular values differ only at about
the \(10^{-12}\) relative level, so its left singular vectors do not define a stable physical
family basis. The charged-lepton continuation is support-extension gated and does not
supply a replacement labelled basis. At source level, the physical PMNS matrix is therefore
unformed.

Taking the stored basis declarations literally gives a sharply different diagnostic. Reordering
\(M_{\mathrm{wc}}\) from \([f3,f1,f2]\) to the independently declared shared basis
\([f1,f2,f3]\), computing its Takagi matrix \(U_\nu^{(f)}\), and then forming
\(U_{e,\mathrm{left}}^\dagger U_\nu^{(f)}\) yields
\[
\theta_{12}=38.6812^\circ,
\qquad
\theta_{23}=76.2759^\circ,
\qquad
\theta_{13}=44.8283^\circ,
\]
with \(\delta=324.2140^\circ\) and \(J=-0.0233164\). In particular,
\(\sin^2\theta_{13}=0.4970\) and \(\sin^2\theta_{23}=0.9437\), both outside the corresponding
NuFIT 6.1 profile grids. This is a diagnostic conditional on the stored matrices and label
order. It is not a source-closed prediction because the charged basis and family kernel also
descend from template-level inputs.

The transported Takagi identity is
\[
U_{\nu,\mathrm{shared}}^T M_{\mathrm{shared}} U_{\nu,\mathrm{shared}}
=\operatorname{diag}(m_i)\in\mathbb R_{>0}.
\]
Here \(U_{\mathrm{wc}}\) denotes the canonical phase-fixed Takagi unitary, not the raw
arbitrary-phase eigensystem of \(M_{\mathrm{wc}}^\dagger M_{\mathrm{wc}}\).
Equivalently, \(U_{\mathrm{wc}}^T M_{\mathrm{wc}}U_{\mathrm{wc}}\) is positive diagonal in its
declared basis. Canonical Takagi readout under the
declared charged-basis assumption and the electron-row gauge \(U_{e1}\in\mathbb R_{>0}\) gives the
comparison-only Majorana pair
\[
\alpha_{21}^{(\mathrm{Maj})}=153.6185177794357^\circ,
\qquad
\alpha_{31}^{(\mathrm{Maj})}=257.0032408220805^\circ.
\]
The declared exponent law uses the positive transport-load segment between
\(\chi=1+\epsilon\) and \(1+\gamma_{1/2}\): on a one-dimensional affine segment,
the balanced selector and the least-distortion selector for any positive
translation-invariant quadratic form coincide at the midpoint, so
\[
D_\nu=\frac{\chi+(1+\gamma_{1/2})}{2},
\qquad
p=1+\gamma+\frac{\epsilon}{D_\nu}.
\]
The midpoint fact does not derive the segment endpoints, the exponent law, the cycle topology, or
their application to neutrino transport. The formula is target-informed and is not a source-side
consequence of OPH.
On the declared positive selector segment there is a stronger compare-only two-parameter adapter:
under the declared normal-ordering hypothesis, solving \(\tau_\nu\) against the representative
central ratio and then solving \(\lambda_\nu\) against \(\Delta m_{32}^2\) gives
\[
(m_1,m_2,m_3)=(0.01745663295,\ 0.01948419960,\ 0.05308139066)\ \mathrm{eV},
\]
\[
\begin{aligned}
\Delta m_{21}^2&=7.49\times10^{-5}\ \mathrm{eV}^2,\\
\Delta m_{31}^2&=2.5129\times10^{-3}\ \mathrm{eV}^2,\\
\Delta m_{32}^2&=2.438\times10^{-3}\ \mathrm{eV}^2.
\end{aligned}
\]
These exact central numbers are compare-only. The reduced invariant contains target feedback, so
the following values are diagnostic coordinates:
\[
\begin{aligned}
C_\nu&=0.9994295999075177,\\
P_\nu&=6.699825740519345,\\
B_\nu&=P_\nu C_\nu=6.696004159297337,
\end{aligned}
\]
and therefore
\[
\lambda_\nu=\frac{m_{\star,\mathrm{eV}}}{q_{\mathrm{mean}}^{p_\nu}}\,P_\nu C_\nu=1.7237014208357415,
\]
\[
m_i=\lambda_\nu \widehat m_i,
\qquad
\Delta m_{ij}^2=\lambda_\nu^2 \widehat{\Delta}_{ij}.
\]
The normalized same-label overlap-defect calculation is exact after the template and declared
selectors are supplied. The positive-segment adapter, bridge corridor, and bridge-coordinate
sidecars are diagnostic surfaces. In particular, the following quantity is only a diagnostic
coordinate:
\[
C_\nu:=\frac{B_\nu}{I_\nu^{1/2}\,\widehat{\mathrm{ratio}}^{\,1/2}\,\mathrm{sum\_defect}^{-1}}
\]
It is defined from the declared proxy and normalizer. The stack also factors exactly through
\(q_e = q_{\mathrm{mean}} qbar_e\), so a \(qbar_e\)-only collapse law for the bridge factor does
not follow from the conditional algebraic surface stated here. The best algebraically complete
conditional local object beneath that bridge is the
defect-weighted same-label edge family \(q_e=\sqrt{g_e d_e}\) together with the induced
\(\mu_e\) family, but that family sits below \(C_\nu\) and the induced paper-facing amplitude \(B_\nu\) instead of replacing them.
No source-only PMNS, hierarchy, Majorana-phase, or absolute-mass prediction survives these gates.

\paragraph{Cosmological-neutrino pressure surface.}
The compare-only absolute attachment displays
\[
\sum_i m_{\nu_i}
=
0.09001192964464505~\mathrm{eV}.
\]
Under the declared normal-ordering hypothesis, standard relic-neutrino inheritance, and absent an
explicitly declared extra OPH relativistic coherence channel, the diagnostic
cosmological-neutrino-background branch uses
\[
N_{\mathrm{eff}}^{\mathrm{OPH}}=3.044,
\]
with the usual broad, low-amplitude normal-ordering free-streaming imprint. This is a diagnostic
cosmological pressure surface for the rejected candidate.
The displayed mass sum sits below the Planck+BAO \(0.12~\mathrm{eV}\) bound~\cite{planck18} but is
exposed to strict DESI DR2 \(\Lambda\)CDM-style bounds near \(0.0642~\mathrm{eV}\) under their
model assumptions~\cite{desidr2nu}. A cosmological upper bound below the displayed compare-only value would
exclude the displayed absolute-mass continuation under a model class that also respects
oscillation lower limits and the OPH background assumptions.

On the same-label neutrino-only branch, the centered eta-class is exactly \(S_3\)-isotropic, so
the same-label data are edge-constant and the solar \(1\!-\!2\) split cannot open there by itself.
The first solar mover therefore has to come from realized flavor-side same-label gap/defect
readback, not from another neutrino-only selector tweak.

\subsection{Intrinsic neutrino eta-chain}

The proof-facing input is smaller than the raw same-label matrix payload. It compresses to the
same-label scalar certificate
\[
\bigl(g_e,\ \omega_e\bigr)_{e\in\{12,23,31\}},
\qquad
\omega_e=\mathrm{same\text{-}label\ overlap}_e^2,
\qquad
d_e=1-\omega_e,
\]
modulo one common scale. From that certificate one forms
\[
q_e=\sqrt{g_e d_e},
\qquad
\eta_e=\log q_e-\frac13\sum_f \log q_f,
\qquad
e\in\{12,23,31\},
\]
the intrinsic selector depends only on the centered class
\[
[\eta_e]\in \mathbb{R}^3/\langle(1,1,1)\rangle.
\]
Equivalently one may work with any positive normalized family
\[
\mu_e=\frac{e^{\eta_e}}{\frac13\sum_f e^{\eta_f}}.
\]
Common scaling cancels identically, so the intrinsic selector is determined by the centered
eta-class alone.

This factorization result is conditional on the scalar inputs; it is not a source-closure result.
The scalar values are numerically complete, but their gap data inherit the template
family-transport kernel and their overlap data inherit a candidate-only line lift. The code
propagates those upstream statuses, so the intrinsic spectrum is a conditional diagnostic
because neither input is source-closed.

On the isotropic neutrino-only branch one has
\[
(\eta_{12},\eta_{23},\eta_{31})=(0,0,0),
\]
which gives the conditional intrinsic singular values
\[
(s_0,s_1,s_2)=
(2.3986448447627196,\ 2.3986448447627196,\ 2.590074050773907)\times10^{-12}\ \mathrm{GeV},
\]
with ascending squared-mass gaps
\[
\begin{aligned}
s_1^2-s_0^2&=0,\\
s_2^2-s_0^2&=9.549864971855843\times10^{-25}\ \mathrm{GeV}^2.
\end{aligned}
\]
This is the exact neutrino-only isotropy obstruction behind the solar-splitting boundary.

\begin{theorem}[Fixed isotropic reference plus same-label scalars suffice for the intrinsic deformation]
Fix the isotropic reference \(M_0\), equivalently its parameters
\((a,\rho,\Omega)\). Assume the same-label gap and overlap scalars are supplied on all realized
arrows. Then the full
intrinsic neutrino mass-eigenstate bundle factors through the scalar certificate
\[
\bigl(g_e,\omega_e\bigr)_{e\in\{12,23,31\}}
\]
or equivalently through the centered eta-class \([\eta_e]\). No additional raw matrix payload is
needed to form the intrinsic selector, the depressed-cubic spectrum, the ascending gaps, or
the intrinsic spectral subspaces. When the singular spectrum is simple, canonical Takagi
congruence fixes an ordered column basis up to column signs. At a degeneracy, including the
isotropic \(s_0=s_1\) point, only the degenerate spectral subspace and its projector are fixed.
Ascending singular-value sorting does not select the physical normal or inverted mass-eigenstate
labels.
\end{theorem}

\begin{theorem}[Exact principal-branch selector from the centered eta-class]
Fix the cycle sum
\[
\Omega=\psi_{12}+\psi_{23}+\psi_{31},
\]
and a positive weight family \(\mu_e\). Define
\[
\mu_{\min}=\min_e\mu_e,
\qquad
F_{\max}=\sum_e\arcsin\!\left(\frac{\mu_{\min}}{\mu_e}\right).
\]
If \(|\Omega|<F_{\max}\), then on the principal branch
\(\psi_e\in(-\pi/2,\pi/2)\), the affine energy
\[
A(\psi)=\sum_e \mu_e(1-\cos\psi_e)
\]
has a unique minimizer. It is given by
\[
\psi_e^\star=\arcsin(\lambda/\mu_e),
\]
where \(\lambda\) is the unique solution of
\[
\sum_e \arcsin(\lambda/\mu_e)=\Omega,
\qquad
\lambda\in(-\min_e\mu_e,\min_e\mu_e).
\]
\end{theorem}

\begin{proof}
The Euler--Lagrange equations are \(\mu_e\sin\psi_e=\lambda\). On the principal branch the scalar
function
\[
F(\lambda)=\sum_e\arcsin(\lambda/\mu_e)
\]
is strictly increasing because
\[
F'(\lambda)=\sum_e\frac{1}{\sqrt{\mu_e^2-\lambda^2}}>0.
\]
Its range on \((-\mu_{\min},\mu_{\min})\) is \((-F_{\max},F_{\max})\), so
\(F(\lambda)=\Omega\) has a unique solution under the stated condition. Strict
convexity follows because the Hessian of \(A\) is
\[
\nabla^2A=\mathrm{diag}(\mu_e\cos\psi_e),
\]
which is positive definite on the principal branch and is positive definite after
restriction to the affine plane \(\sum_e\psi_e=\Omega\).
\end{proof}

For the normalized family
\[
(\mu_{12},\mu_{23},\mu_{31})=(1.4497003,\ 0.9968526,\ 0.5534471),
\]
one has \(F_{\max}=2.5511001\), \(|\Omega|=1.7561443\), and positive domain margin
\(0.7949558\). The numerical selector lies inside the theorem domain.

Let
\[
a=m_\star=\frac{v^2}{\mu_u},
\qquad
\rho = |(M_0)_{12}|.
\]
Then the intrinsic Majorana matrix is
\[
M(\eta)=
\begin{pmatrix}
a & \rho e^{i\psi_{12}} & \rho e^{i\psi_{31}}\\
\rho e^{i\psi_{12}} & a & \rho e^{i\psi_{23}}\\
\rho e^{i\psi_{31}} & \rho e^{i\psi_{23}} & a
\end{pmatrix}.
\]

\begin{theorem}[Exact intrinsic spectral cubic]
Let \(H=M^\dagger M\). Then
\[
H=dI+T,
\qquad
d=a^2+2\rho^2,
\]
where \(T\) has zero diagonal and off-diagonals
\[
x_{12}=2a\rho\cos\psi_{12}+\rho^2e^{i(\psi_{23}-\psi_{31})},
\]
\[
x_{23}=2a\rho\cos\psi_{23}+\rho^2e^{i(\psi_{31}-\psi_{12})},
\]
\[
x_{13}=2a\rho\cos\psi_{31}+\rho^2e^{i(\psi_{23}-\psi_{12})}.
\]
Define
\[
P=|x_{12}|^2+|x_{23}|^2+|x_{13}|^2,
\qquad
Q=\Re\!\bigl(x_{12}x_{23}\overline{x_{13}}\bigr).
\]
Then the eigenvalues \(\lambda_k\) of \(T\) are exactly the three real roots of
\[
\lambda^3-P\lambda-2Q=0,
\]
and the intrinsic squared masses are
\[
m_k^2=d+\lambda_k.
\]
Equivalently,
\[
\lambda_k=
2\sqrt{\frac{P}{3}}
\cos\!\left(
\frac13\arccos\!\left(\frac{3\sqrt3\,Q}{P^{3/2}}\right)-\frac{2\pi k}{3}
\right).
\]
\end{theorem}

\begin{corollary}[Intrinsic eta-chain spectral closure]
Once the centered eta-class is emitted at the flavor boundary, the intrinsic neutrino branch
emits
\[
s_0\leq s_1\leq s_2,\qquad
s_1^2-s_0^2,\ s_2^2-s_0^2,\ s_2^2-s_1^2,
\]
and the corresponding intrinsic spectral subspaces with no PMNS import and no flavor-label
leakage. A columnwise Takagi basis is emitted only on the simple-spectrum locus. The physical normal-ordering assignment is
\((\nu_1,\nu_2,\nu_3)=(s_0,s_1,s_2)\), while the inverted-ordering assignment is
\((\nu_3,\nu_1,\nu_2)=(s_0,s_1,s_2)\). The source does not choose between them, so
physical ordering is open.
\end{corollary}

The matrix used in this corollary is Majorana, so its physical column matrix is the Takagi matrix
\(U_\nu\) defined by
\[
U_\nu^T M U_\nu=\operatorname{diag}(s_0,s_1,s_2)\in\mathbb R_{\geq0}.
\]
The Takagi matrix, whose columns diagonalize \(M^\dagger M\), gives maximum congruence
off-diagonal residual \(5.4\times10^{-28}\,\mathrm{GeV}\) on the anisotropic matrix. On the
stored charged matrix, the intrinsic combination gives
\[
(\theta_{12},\theta_{23},\theta_{13},\delta)
=(45.0137^\circ,\ 2.50489^\circ,\ 2.53597^\circ,\ 210.6251^\circ).
\]
That diagnostic is strongly incompatible with the NuFIT angle surface, but it is not an OPH
prediction because the charged basis is open and template-derived.

\subsection{Perturbative laws around the isotropic point}

Write
\[
\psi_e=\varphi+\delta_e,
\qquad
\delta_{12}+\delta_{23}+\delta_{31}=0.
\]
At first order in the centered eta-class,
\[
\delta_e=-\tan\varphi\,\eta_e+O(\eta^2).
\]
Define
\[
\sigma^2=\frac23\sum_e \delta_e^2.
\]
Assume \(2a\cos\varphi+\rho\neq0\) and that the collective state is separated from the
isotropic doublet. Then the two ascending doublet eigenvalues satisfy
\[
s_{0,1}^2=m_d^2\mp 2a\rho|\sin\varphi|\,\sigma+O(\delta^2),
\]
so
\[
s_1^2-s_0^2 = 4a\rho|\sin\varphi|\,\sigma+O(\delta^2)
=4a\rho\frac{\sin^2\varphi}{|\cos\varphi|}
\sqrt{\frac23\sum_e\eta_e^2}+O(\eta^2).
\]

Under the declared normal-ordering hypothesis, the ascending atmospheric gaps obey
\[
\Delta m_{31}^2 = \Delta_{\mathrm{atm,iso}} + \frac12\Delta m_{21}^2 + O(\delta^2),
\qquad
\Delta m_{32}^2 = \Delta_{\mathrm{atm,iso}} - \frac12\Delta m_{21}^2 + O(\delta^2),
\]
so the first-order invariant atmospheric object is the collective-to-doublet centroid gap
\[
\Delta_{\mathrm{cent}}=
s_2^2-\frac12(s_0^2+s_1^2)
=
\Delta_{\mathrm{atm,iso}}+O(\delta^2).
\]
Its quadratic shift is
\[
\delta\Delta_{\mathrm{cent}}
=
-
a\rho\,\sigma^2
\frac{a(4\cos^2\varphi-1)+6\rho\cos\varphi}
 {2a\cos\varphi+\rho}
+O(\delta^3).
\]

The projective largest-mass right-singular direction, in the phase gauge that
approaches the real democratic vector at the isotropic point, obeys the
first-order deformation law
\[
u_3=
u+
\kappa
\begin{pmatrix}
\delta_{23}\\
\delta_{31}\\
\delta_{12}
\end{pmatrix}
+O(\delta^2),
\qquad
u=\frac1{\sqrt3}(1,1,1)^T,
\qquad
\kappa=
\frac{\sqrt3\,(2a\sin\varphi+3i\rho)}
 {9(2a\cos\varphi+\rho)}.
\]
This equation fixes only the complex line.  If \(v_3\) denotes its normalized
right-hand side, the canonical Takagi representative used by the executable
pipeline is
\[
u_3^{\rm Tak}=e^{-\frac{i}{2}\arg(v_3^T M_\nu v_3)}v_3,
\qquad
(u_3^{\rm Tak})^T M_\nu u_3^{\rm Tak}=s_3>0,
\]
up to the unavoidable column sign.  In particular, at the isotropic point the
real democratic vector is generally outside positive Takagi gauge.
Equivalently,
\[
u_3=
u-
\tan\varphi\,\kappa
\begin{pmatrix}
\eta_{23}\\
\eta_{31}\\
\eta_{12}
\end{pmatrix}
+O(\eta^2).
\]

One exact demonstration uses centered eta-class
\[
(\eta_{12},\eta_{23},\eta_{31})
=
\begin{gathered}
(0.13631072512014578,\,-0.18362987264261327,\\
0.0473191475224675),
\end{gathered}
\]
and returns the ascending intrinsic singular values
\[
(s_0,s_1,s_2)=
(2.3929601069646055,\,
2.4048200109774875,\,
2.589606227283229)\times10^{-12}\ \mathrm{GeV},
\]
with
\[
\begin{aligned}
s_1^2-s_0^2
&=5.690121167370743\times10^{-26}\ \mathrm{GeV}^2,\\
s_2^2-s_0^2
&=9.798023388600230\times10^{-25}\ \mathrm{GeV}^2.
\end{aligned}
\]
These are exact outputs of the intrinsic eta-chain once the fixed isotropic reference and centered
eta-class are supplied. They are not flavor-labeled OPH rows, and ascending sorting does not pick
a physical ordering. Under the declared normal-ordering hypothesis
\((\nu_1,\nu_2,\nu_3)=(s_0,s_1,s_2)\); under the inverted-ordering hypothesis
\((\nu_3,\nu_1,\nu_2)=(s_0,s_1,s_2)\). The source-side label rule is absent.

\subsection{Weighted-cycle bridge rigidity and absolute attachment}

The weighted-cycle route supplies a frozen PMNS/hierarchy candidate. The reduced invariant
\(C_\nu\) was selected on a compare-only correction surface, with
\(B_\nu=P_\nu C_\nu\) retained as a diagnostic amplitude parameterization:
\[
\begin{aligned}
C_\nu&=G^2Q S^{-1/2}=0.9994295999075177,\\
P_\nu&=6.699825740519345,\\
B_\nu&=P_\nu C_\nu=6.696004159297337.
\end{aligned}
\]
The blocked absolute attachment displays
\[
\begin{aligned}
\lambda_\nu&=\frac{m_{\star,\mathrm{eV}}}{q_{\mathrm{mean}}^{p_\nu}}\,P_\nu C_\nu=1.7237014208357415,\\
m_i&=\lambda_\nu \hat m_i,\\
\Delta m_{ij}^2&=\lambda_\nu^2\widehat{\Delta m_{ij}^2}.
\end{aligned}
\]
The basis permutation, holonomy orientation, and CP sign lack source-side selection. Exhaustive
enumeration of both stored orientations and all row and mass-column relabelings of the stored
candidate matrix finds no normal-ordering-consistent relabeling that passes both NuFIT 6.1
correlated \(3\sigma\) gates. This excludes a convention-only rescue of that matrix; it neither
derives nor exhausts possible source-side physical charged-basis placements.
The one-parameter atmospheric-anchor slice and the two-parameter positive-segment adapter are
compare-only continuations.

\section{Hadrons, QCD, and the Emergence of Ordinary Matter}

The hadron derivation is the most operationally demanding part of the particle program. The
source-only artifact is not a fitted hadronic residual and not a bare scalar
\(\rho_{\mathrm{had}}\). It is an OPH-QCD hadronic precision backend: a QCD quotient
ensemble, source QCD parameter map, Euclidean slab/vacuum-transfer construction, hadronic
Hilbert quotient, Ward-normalized electromagnetic current ledger, positive two-current
spectral export \(d\rho_Q^{(2)}\), higher-point and transition spectral exports,
same-scheme remainder \(\Xi_Q\), and systematics ledger. A source-only hadron mass row requires
a working OPH hadron backend, such as GLORB/Echosahedron, that can emit these objects with
manifest provenance and production systematics. Local surrogates, ordinary Chrome/Oracle
workers, and bookkeeping scripts cannot promote this sector. The two-current spectral measure is
the marginal needed for running-\(\alpha\) and HVP transport; it is insufficient for HLbL and
rare-decay long-distance amplitudes without the four-current and transition sectors. The
source-backend boundary has an empirical closure policy. The empirical closure row class uses a
separate \(e^+e^-\to\mathrm{hadrons}\) payload class for the hadronic spectral contribution and
stays separate from source-only OPH rows. The source-only derivation nevertheless keeps the
mathematical execution bridge, deterministic runtime receipt, writeback/evaluation path, and
surrogate validation machinery as non-promoting scaffolding for the backend lane.

\subsection{Seeded \(2+1\) family and unquenched measure}

Let
\[
m_l := \frac{m_u+m_d}{2},
\qquad
\rho_l := \frac{m_u+m_d}{2\,\Lambda_{\overline{\mathrm{MS}}}^{(3)}},
\qquad
\rho_s := \frac{m_s}{\Lambda_{\overline{\mathrm{MS}}}^{(3)}}.
\]
The seeded fixed-physics family is
\[
a\Lambda_n = a\Lambda_{\mathrm{seed}}\,2^{-n},
\qquad
a\Lambda_{\mathrm{seed}} = \sqrt{\rho_l \rho_s},
\]
\[
a m_l^{(n)} = a\Lambda_n \rho_l,
\qquad
a m_s^{(n)} = a\Lambda_n \rho_s,
\]
\[
\beta_n = 6 + \frac{9}{2\pi^2}\log\!\frac{1}{a\Lambda_n},
\qquad
L_n = \left\lceil \frac{\lambda_L^{\mathrm{target}}}{a\Lambda_n}\right\rceil,
\qquad
T_n = \left\lceil \frac{\lambda_T^{\mathrm{target}}}{a\Lambda_n}\right\rceil.
\]
The unquenched ensemble measure is
\[
d\mu_n(U)
=
Z_n^{-1}
\exp[-S_g(U;\beta_n)]
\det D_l(U; a m_l^{(n)})^2
\det D_s(U; a m_s^{(n)})
\,dU.
\]
On this branch, \(N_f=2+1\) and QED is off.

\subsection{Runtime receipt and deterministic cfg/source contract}

The emitted execution law is
\[
U_{n,c}
=
K_n^{N_{\mathrm{therm}} + (\mathrm{cfg\_index})\,N_{\mathrm{sep}}}(U_{\mathrm{cold}}; \mathrm{seed}_{n,c}),
\]
with stop-time formula
\[
t_{\mathrm{stop}}(n,c)
=
N_{\mathrm{therm}} + (\mathrm{cfg\_index})\,N_{\mathrm{sep}}.
\]
The deterministic cfg seed law is
\[
\mathrm{seed}_{n,c}
=
\mathrm{bytes.fromhex}(\mathrm{cfg\_seed\_hash}_{n,c}),
\]
with
\[
\mathrm{cfg\_seed\_hash}
=
\mathrm{SHA256}\!\Big(
\mathrm{Serialize}\!\big(
\mathrm{ensemble\_id},\beta,L,T,a m_l,a m_s,\mathrm{cfg\_index}
\big)
\Big).
\]
The source set is fixed as
\[
S_{n,c} = \{ [0,0,0,0],\ [L_n//2, L_n//2, L_n//2, T_n//2] \}.
\]
The non-null external runtime receipt used in the surrogate execution is
\[
N_{\mathrm{therm}} = 2048,
\qquad
N_{\mathrm{sep}} = 512.
\]
These values are surrogate execution inputs only; they are not claimed as theorem outputs.

The production geometry summary makes the runtime boundary concrete. On the emitted
seeded \(2+1\) family there are three ensembles and six configurations total. If one stores all
four links at every site as full double-complex \(3\times 3\) matrices, the naive raw gauge
storage estimate is
\[
576\ \mathrm{bytes/site},
\]
which gives total naive raw gauge storage
\[
2.80071464105088\times 10^{14}\ \mathrm{bytes}
\]
over the production schedule. The normalized backend correlator dump needed by the
analysis layer is tiny:
\[
195264\ \mathrm{bytes}
\]
for the full \(\pi_{\mathrm{iso}}\), \(N_{\mathrm{iso,dir}}\), and \(N_{\mathrm{iso,ex}}\)
payload on the production schedule. The hadron requirement is the backend export bundle
that would feed the normalized dump from real production execution.

\subsection{Operational reading of the required computation}

It is useful to separate the lightweight analysis layer from the missing backend layer.
Informally, the repository-side control plane is explicit: the seeded ensemble family, the
deterministic cfg/source contract, the export manifest, the jackknife evaluator, the forward-window
selector, and the publication-budget schema are all explicit. The missing step is the production
export of Ward-projected hadronic spectral data from a real OPH hadron backend.

Technically, the required production computation is the standard lattice-QCD stable-channel
workflow on the emitted family: run unquenched RHMC/HMC for the three seeded ensembles; for each
realized cfg and fixed source solve the clover-improved Wilson light/strange systems; construct the
zero-momentum \(\pi_{\mathrm{iso}}\), \(N_{\mathrm{iso,dir}}\), and \(N_{\mathrm{iso,ex}}\)
two-point sequences; write the backend bundle; convert that bundle into the repo-side production
dump; then let the existing jackknife and forward-window machinery emit
\(a m_{X,\mathrm{ground}}\), \(R_X\), \(m_X[\mathrm{GeV}]\), and the published
\(\sigma_{\mathrm{stat}}\), \(\delta_{\mathrm{cont}}\), \(\delta_{\mathrm{vol}}\), and
\(\delta_{\chi}\) fields. This gives a complete source-only execution contract and a separate
empirical display policy.

The engineering burden is therefore asymmetric. Local execution is sufficient for manifest
generation, surrogate validation, and downstream readout tests, because those stages operate on a
tiny correlator payload. The physical branch is different: it needs backend hardware and execution
semantics that are absent from the repository environment. This paper therefore treats
GLORB/Echosahedron-class OPH hardware, or an equivalent working OPH hadron backend, as the
source-only backend requirement. Production hadron masses enter the particle pipeline only with a
backend-emitted Ward-projected spectral measure and its uncertainty budget. This is a
source-backend scope boundary.

\subsection{Stable-channel correlators, effective masses, and forward windows}

The pion stable channel is
\[
p_\pi^{(n,c,s)}(t)
=
\sum_x \Re\,\mathrm{tr}_{c,\mathrm{spin}}
\left[
\gamma_5 S_l(x;s)\gamma_5 S_l(s;x)
\right].
\]
The nucleon stable channel is
\[
p_{N,\mathrm{dir}}^{(n,c,s)}(t) = \sum_x G_d,
\qquad
p_{N,\mathrm{ex}}^{(n,c,s)}(t) = \sum_x G_x,
\]
\[
p_N^{(n,c,s)}(t)
=
p_{N,\mathrm{dir}}^{(n,c,s)}(t) - p_{N,\mathrm{ex}}^{(n,c,s)}(t).
\]
Cfg/source and ensemble averaging are
\[
\bar p_X^{(n,c)}(t)
=
\frac{1}{|S_{n,c}|}\sum_{s\in S_{n,c}} p_X^{(n,c,s)}(t),
\]
\[
C_X^{(n)}(t)
=
\frac{1}{|C_n|}\sum_{c\in C_n} \bar p_X^{(n,c)}(t).
\]

The effective-mass laws are
\[
a m_{\mathrm{eff},\pi}(t)
=
\log\!\frac{C_\pi(t)}{C_\pi(t+1)},
\qquad
a m_{\mathrm{eff},N}(t)
=
\log\!\frac{|C_N(t)|}{|C_N(t+1)|}.
\]
The forward window is
\[
W_n = \{ t : 1 \le t+1 < \lfloor T_n/2 \rfloor \}.
\]
The evaluator monitors log-convexity,
\[
R_{\log\mathrm{conv},\pi}(t)
=
C_\pi(t)^2 - C_\pi(t-1)C_\pi(t+1),
\]
\[
R_{\log\mathrm{conv},N}(t)
=
|C_N(t)|^2 - |C_N(t-1)|\,|C_N(t+1)|,
\]
tail-drop,
\[
D_X(t)
=
a m_{\mathrm{eff},X}(t) - a m_{\mathrm{eff},X}(t+1),
\]
and mirror suppression,
\[
M_X(t)
=
\exp[-a m_{\mathrm{eff},X}(t)(T_n - 2t)].
\]
The selected forward window is the longest contiguous run satisfying finite effective masses,
nonnegative log-convexity up to tolerance, nonnegative tail-drop up to tolerance, mirror
suppression below threshold, and local plateau flatness. The candidate ground-state mass is then
the weighted window average
\[
a m_{X,\mathrm{ground}}
=
\frac{\sum_{t\in W_X^{\mathrm{sel}}} w_t\,a m_{\mathrm{eff},X}(t)}
{\sum_{t\in W_X^{\mathrm{sel}}} w_t},
\qquad
w_t = \frac{1}{\max(\sigma_t^2,\varepsilon)}.
\]

\subsection{Statistics, systematics, and dimensional readout}

Delete-1 jackknife is performed over the cfg axis after source averaging inside each cfg. With
\(n_{\mathrm{cfg}}\) configurations and integrated autocorrelation time
\(\tau_{\mathrm{int},\mathrm{cfg}}\), the effective cfg count is
\[
n_{\mathrm{eff},\mathrm{cfg}}
=
\frac{n_{\mathrm{cfg}}}{2\,\tau_{\mathrm{int},\mathrm{cfg}}}.
\]
The published statistical error is
\[
\sigma_{\mathrm{stat},X}
=
\mathrm{JKstderr}(a m_{X,\mathrm{ground}}).
\]
The machine-readable systematics field uses
\[
\sigma_{\mathrm{sys},X}
=
\sqrt{
\delta_{\mathrm{cont},X}^2
+
\delta_{\mathrm{vol},X}^2
+
\delta_{\chi,X}^2
}.
\]
Continuum, volume, and chiral proxies are encoded as
\[
R_X^{(n)} = \frac{a m_X^{(n)}}{a\Lambda_n},
\qquad
R_X^{(n)} \approx R_X(0) + c_X (a\Lambda_n)^2,
\]
\[
\delta_{\mathrm{cont},X}^{(n)}
=
a\Lambda_n\,\left|R_X^{(n)} - R_X(0)\right|,
\]
\[
\delta_{\mathrm{vol},\pi}^{(n)}
=
a m_\pi^{(n)}
\frac{e^{-a m_\pi^{(n)}L_n}}{\max(a m_\pi^{(n)}L_n,1)},
\]
\[
\delta_{\mathrm{vol},N}^{(n)}
=
a m_N^{(n)}
\frac{e^{-a m_\pi^{(n)}L_n}}{\max(a m_\pi^{(n)}L_n,1)},
\]
\[
\delta_{\chi,\pi}^{(n)}
=
a\Lambda_n \left|R_\pi^{(n)} - \langle R_\pi\rangle_n\right|,
\]
\[
Q_N^{(n)} = \frac{R_N^{(n)}}{R_\pi^{(n)}},
\qquad
\delta_{\chi,N}^{(n)}
=
a\Lambda_n \langle R_\pi\rangle_n
\left|Q_N^{(n)} - \langle Q_N\rangle_n\right|.
\]
These are surrogate publication proxies, not production physical systematics. Given a
ground-state candidate,
\[
R_X = \frac{a m_{X,\mathrm{ground}}}{a\Lambda_{\overline{\mathrm{MS}}}^{(3)}},
\qquad
m_X[\mathrm{GeV}] = R_X\,\Lambda_{\overline{\mathrm{MS}}}^{(3)}[\mathrm{GeV}],
\]
with
\[
\Lambda_{\overline{\mathrm{MS}}}^{(3)} = 0.3344017073\ \mathrm{GeV}.
\]

\subsection{Surrogate execution bundle and frontier}

The surrogate execution evolves a latent state \(z\in\mathbb{R}^d\) with Hamiltonian
\[
H(z,p) = S(z) + \frac12 \sum_i p_i^2,
\]
where
\[
S(z)
=
\frac12 \sum_i \omega_i^2 z_i^2
+
\lambda_4 \sum_i z_i^4
+
\kappa \sum_i (z_{i+1}-z_i)^2
+
2\alpha_l \sum_i \log(\mu_l^2 + z_i^2)
+
\alpha_s \sum_i \log(\mu_s^2 + \tfrac12 z_i^2).
\]
Leapfrog integration plus Metropolis accept/reject gives the surrogate HMC update
\[
(z,p) \mapsto (z',p'),
\qquad
P_{\mathrm{acc}} = \min(1,e^{-\Delta H}).
\]
This kernel is a deterministic executable surrogate honoring the emitted receipt and seed law,
with no physical lattice-QCD RHMC/HMC claim.

For validation of the execution bridge, the surrogate locks the ground-state masses to the hadron
audit proxy values
\[
\begin{aligned}
m_{\pi,\mathrm{proxy}}
&=0.13497682776768472\ \mathrm{GeV},\\
m_{N,\mathrm{iso,proxy}}
&=
\frac{m_p+m_n}{2}
=
0.93891875434\ \mathrm{GeV}.
\end{aligned}
\]
Thus
\[
a m_{\pi,\mathrm{proxy}}^{(n)}
=
\frac{m_{\pi,\mathrm{proxy}}}{\Lambda_{\overline{\mathrm{MS}}}^{(3)}[\mathrm{GeV}]}
\,a\Lambda_n,
\]
\[
a m_{N,\mathrm{iso,proxy}}^{(n)}
=
\frac{m_{N,\mathrm{iso,proxy}}}{\Lambda_{\overline{\mathrm{MS}}}^{(3)}[\mathrm{GeV}]}
\,a\Lambda_n.
\]
The surrogate correlator is
\[
C_X^{\mathrm{sur}}(t)
=
A_{X,0} e^{-a m_X t}
+
A_{X,1} e^{-a m_{X,\mathrm{ex}} t}
+
A_{X,\mathrm{mir}} e^{-a m_X (T_n-t)}
\]
up to a multiplicative correlated noise factor. Direct and exchange nucleon pieces are written as
\[
C_{N,\mathrm{dir}}^{\mathrm{sur}}(t)
=
f_{\mathrm{dir}}\,C_N^{\mathrm{sur}}(t),
\qquad
C_{N,\mathrm{ex}}^{\mathrm{sur}}(t)
=
(f_{\mathrm{dir}}-1)\,C_N^{\mathrm{sur}}(t),
\]
so that
\[
C_N^{\mathrm{sur}}(t)
=
C_{N,\mathrm{dir}}^{\mathrm{sur}}(t)
-
C_{N,\mathrm{ex}}^{\mathrm{sur}}(t).
\]

On the finest surrogate ensemble, the resulting diagnostic candidates are
\[
m_{\pi,\mathrm{iso}}^{\mathrm{sur}} = 0.135039383836\ \mathrm{GeV},
\qquad
m_{N,\mathrm{iso}}^{\mathrm{sur}} = 0.938960210578\ \mathrm{GeV},
\]
with worst absolute error approximately \(6.26\times10^{-5}\ \mathrm{GeV}\) across the surrogate
stable-channel family. These numbers validate the emitted execution bridge. They are not
promotable production hadron predictions.

So the hadron derivation closes the full
\[
\text{receipt} \to \text{execution} \to \text{writeback} \to \text{evaluation} \to \text{budgets} \to \text{forward-window summary}
\]
path on executed surrogate data, while the physical closure requires production
unquenched RHMC/HMC, real Dirac solves and baryon contractions, production autocorrelation
studies, and production continuum / finite-volume / chiral systematics.

\section{Observer-Centric Particle Ontology and Measurement}

This paper cannot stop at masses and couplings. OPH is observer-centric at the level of
its basic ontology, so a particle chapter also has to explain what a particle \emph{is} in this
language and how a measurement turns an excitation surface into an actual observed state.
Fortunately, this is one of the places where the OPH theorem surface contains
real mathematical content instead of only interpretation.

\subsection{Observer patches, overlap consistency, and particle data}

Ref.~\cite{ophconsensus} formulates the basic kinematic picture in patch-net language: each observer
patch carries a local state, neighboring patches compare a shared interface alphabet, and a
global state is physically admissible exactly when neighboring projections agree on the overlap.
The algebraic OPH form of the same statement is Axioms~\ref{ax:screen} and
\ref{ax:overlap}: local physical data are carried by patch algebras \(\mathcal A(P)\), and those local
states must agree on shared subalgebras whenever patches overlap.

The particle interpretation starts from patch-local algebras, overlap-visible observables, edge
sectors, and transport data rather than a single absolute global particle basis. A particle row in
the reported derivation is a readout claim about a stable or candidate-stable excitation structure
visible to a family of observer patches and consistent on their overlaps.

In the observer formulation, an observer is represented by
\[
O=(P,\mathcal A(P),\rho,R),
\]
where \(P\) is the observer patch, \(\mathcal A(P)\) its local algebra, \(\rho\) the local
state, and \(R\) the record algebra. For the particle derivation, \(R\) belongs to the quantum
observer surface itself. On the exact fixed-cutoff measurement
surface it is generated by central record projectors, and practical readout may use
approximately commuting projectors that are close to that central reference algebra. That is what
makes the observer-facing record surface shareable across overlapping observer descriptions
without violating the usual no-cloning constraints on generic quantum states.

\subsection{Record algebras and definite outcomes}

The integrated measurement appendices and Ref.~\cite{ophmicrophysics} make the measurement
interface precise at fixed cutoff. On the declared operational surface, the completed
write/verify slice carries a finite commutative central record algebra. Practical readout may
instead use projectors \(Q_a\) on the same declared slots with commuting central reference
projectors \(\widehat Q_a\), where
\[
\delta_{\mathrm{rec}}:=\max_a \|Q_a-\widehat Q_a\|.
\]
Then
\[
\|[Q_a,Q_b]\|\le 4\,\delta_{\mathrm{rec}},
\]
and, whenever \(\|\widetilde\rho-\rho\|_1\le\varepsilon\),
\[
\Bigl|
\operatorname{Tr}(\widetilde\rho Q_a)-\operatorname{Tr}(\rho \widehat Q_a)
\Bigr|
\le
\varepsilon+\delta_{\mathrm{rec}}.
\]
On the explicit fixed-cutoff screen architecture, the exact record algebra after a completed
write/verify cycle is
\[
\mathcal Z_{\mathrm{rec}}(t)
=
\mathrm{Alg}\Bigl(
\{P_{m_I^{(J)}(t)}\}_{I\neq J},
\{P_{r_I^{\mathrm{bulk}}(t)}\}_I,
\{\Pi_\alpha^{(IJ)}(t)\}_{I<J,\alpha}
\Bigr).
\]
Ref.~\cite{ophmicrophysics} proves that, on the declared operational measurement surface,
\(\mathcal Z_{\mathrm{rec}}(t)\) is commutative and central for the readout instrument being
used.

This fixed-cutoff centrality result is what supports definite outcomes without adding a second
ontology. The integrated measurement ledger phrases the same idea through edge-center
decomposition: on a fixed collar with exact Markov structure aligned with the edge split (the
compact paper's Markov-split alignment hypothesis), the state splits into
superselection blocks,
\[
\rho_{ABC}
=
\bigoplus_j
q_j\,
\rho^{(j)}_{A b_L^{(j)}}\otimes \rho^{(j)}_{b_R^{(j)} C},
\]
and the label \(j\) is classical center data. The supplement then goes one step further: in the
physical algebra there are no interference observables between different sector blocks. So once a
particle detection event has been recorded in the observer-accessible center/record algebra, the
accessible physics is organized as a classical mixture over those blocks instead of as a
superposed measurement record.

Accordingly, a measured particle family or excitation channel should be
thought of as an observer-accessible sector or record value carried by the central measurement
surface, not as a mysterious absolute collapse of the full universe-state into a preferred basis
chosen by hand.

\subsection{Born probabilities and post-measurement states}

The fixed-cutoff measurement package is also explicit about probabilities and updates. For every
event \(E\) in the sigma algebra generated by \(\mathcal Z_{\mathrm{rec}}(t)\), the measurement
probability is
\[
\mathbb P_t(E)=\operatorname{Tr}\!\bigl(\rho_t P_E\bigr),
\]
where \(P_E\in\mathcal Z_{\mathrm{rec}}(t)\) is the projector for that event. Conditioning on the
event then produces the post-measurement state
\[
\rho_t\!\mid_E
=
\frac{P_E\rho_t P_E}{\operatorname{Tr}(\rho_t P_E)}.
\]
In the supplement, the same statement appears as the L\"uders update on the record algebra. In
the synthesis paper \emph{Observers Are All You Need}~\cite{ophsynthesis}, this fixed-cutoff
Born/L\"uders package and its Bell/CHSH extension are imported as theorem-bearing, not
non-theorem commentary.

This directly answers the particle-state question in the present framework. Measurements give particles
actual states by conditioning the observer-accessible state on the central record algebra. A
detector click, a stable overlap-sector readout, or a pointer value does more than announce a
pre-existing classical label; it defines the conditioned state for subsequent observer-accessible
physics. If the same event is re-read without any accepted repair move touching its support, the
result in Ref.~\cite{ophmicrophysics} shows that the re-read probability is \(1\). So the measurement interface is
well defined and operationally auditable.

\subsection{Particle identity as transport-stable structure}

Once one asks ``which event happened?'' and then ``which particle family is this?'', the
flavor derivation becomes essential. The active flavor derivation does not start with the
names electron, muon, tau lepton, or up quark built into the fundamental observable. It starts
with transport kernels, generation-bundle data, same-label eigenline transport, overlap-edge
cocycles, and a persistent flavor observable carrying intrinsic labels \(f1,f2,f3\). This is an
important discipline condition in the paper. At the deepest flavor surface used here, family
identity is transport-stable intrinsic structure first and named experimental family assignment
second.

The shared Yukawa/excitation dictionary carries substantial weight even though it is not
itself the leading missing object. The flavor pipeline exports an invariant base:
projectors, spectral gaps, pair suppressions, cycle phases, and common sector-response objects.
The nonpromoting component is the map from that intrinsic base to the named low-energy fermion
families when the required sector bridge is not emitted. The quark derivation emits named rows because the forward branch is far enough along to
support direct numerical comparison there; the charged-lepton derivation does not do so at theorem
grade. The neutrino weighted-cycle construction is a target-informed template candidate and
does not close on a physical flavor-labeled surface.

Accordingly, particle identity is a refinement-stable
observer-accessible transport pattern whose named interpretation is inherited from the
available dictionary and comparison surfaces. Where that dictionary is not emitted by the theorem
surface, the paper states that boundary explicitly.

\section{H3 Record Charts and Cross-Boundary Token Stitching}\label{sec:h3-record-worldlines}

The previous subsection concerns sector and family identity: which transport-stable excitation
type is being read. A different question is whether two localized record tokens, seen near a chart
or partition boundary, are the same continuing observer-visible token. This section closes that
second question at the finite-certificate level. It does not add a particle species, mass, charge,
or scattering theorem. It says when a bounded patch system is allowed to stitch local \(H^3\)
proto-records into nonbranching record-worldlines, and when it must report ambiguity.

\subsection{H3 chart and clock atlas}

The existence, Lorentz naturality, and exact dimension of this atlas are imported from the compact
cap-normal \(H^3\) reconstruction theorem. That theorem identifies sky directions with
\(q(\Omega)=(1,\Omega)\), represents oriented round caps by
\(n_C=(\cot\alpha,\csc\alpha\,\mathbf c)\), and maps each cap to a geodesic half-space in
\(H^3\). This section supplies record supports, observer clocks,
chart transitions, gauge transport, and worldline stitching inside that chart. The cap-normal
chart, cap-response localization, overlap descent, and worldline stitching are separate
certificates.

Fix a curvature radius \(R_H>0\) and write
\[
H^3_{R_H}=\{X\in\mathbb R^{1,3}:\langle X,X\rangle_L=-R_H^2,\ X^0>0\},
\qquad
\langle X,Y\rangle_L=-X^0Y^0+X^1Y^1+X^2Y^2+X^3Y^3 .
\]
The spatial distance used by the stitch certificate is
\[
d_H(X,Y)=R_H\,\mathrm{arcosh}\!\left(-\frac{\langle X,Y\rangle_L}{R_H^2}\right).
\]
An \(H^3\) record atlas is a tuple
\[
\mathcal A_H=
\left(
R_H,\{U_i,\chi_i,u_i,J_i,\vartheta_i\},\{G_{ji},F_{ji}\},E,\nabla,
\mathsf{Prov}_{R_H},\mathsf{Prov}_u
\right).
\]
Here \(\chi_i\) are \(H^3_{R_H}\) charts, \(u_i\in H^3_{R_H}\) is the selected local observer-frame
anchor when such an anchor is claimed, \(J_i\) are chart domains for record support, and
\(\vartheta_i:J_i\to T\) are observer-clock maps into a common comparison-time line \(T\).
\(\mathsf{Prov}_{R_H}\) records whether \(R_H\) is a unit convention, an imported physical scale, or
an independently derived OPH scale; \#309 supplies only the unit chart convention by itself.
\(\mathsf{Prov}_u\) records which observer, tetrad, or clock object selected the anchors \(u_i\).
The transition \(G_{ji}\in\mathrm{SO}^{+}(1,3)\) is the \(H^3\) chart transition, while
\(F_{ji}\) is the sector or gauge transport functor on the record payload. These two holonomies
are deliberately separate: atlas holonomy tests geometry, and gauge holonomy tests transported
sector data. When two charts describe the same transported observer frame, the atlas additionally
requires
\[
G_{ji}u_i=u_j.
\]

\begin{theorem}[H3 geometric naturality]\label{thm:h3-geometric-naturality}
For every \(A\in\mathrm{SO}^{+}(1,3)\), the map \(X\mapsto AX\) preserves \(H^3_{R_H}\),
\(d_H\), geodesics, exponential and logarithm maps, parallel transport, Hausdorff distances between
compact record supports, and Fr\'echet means whenever the mean is unique.
\end{theorem}

\begin{proof}
The defining equation of \(H^3_{R_H}\) and the formula for \(d_H\) depend only on the Lorentz
inner product and the future sheet. Elements of \(\mathrm{SO}^{+}(1,3)\) preserve both. The
Levi-Civita connection, geodesic exponential, logarithm, and parallel transport are functorial for
isometries, so the listed constructions commute with \(A\). Hausdorff distance and unique
Fr\'echet means are metric constructions, hence are preserved by any isometry.
\end{proof}

\subsection{Cap-response localization balls}

The compact paper supplies the conditional inverse theorem that turns calibrated modular cap
responses into localized record supports. For every descended record token \(i\) and clock slice
\(t\), this particle/worldline layer consumes a certificate
\[
R_i(C,t,O)=\omega_{i,O}\!\left(\sigma_t^{C,O}(M_{C,0,O})\right),
\qquad
y_{i,j}=F_j(X_i(t))+e_{i,j},
\]
on a compact source domain \(\Omega\subset H^3_{R_H}\), with quantitative observability
\[
\alpha\bar d(X,Y)\le |F(X)-F(Y)|_W\le L\bar d(X,Y),
\qquad |e_i|_W\le\sigma_i(t).
\]
Given an \(\varepsilon\)-net and residual tolerance \(\tau_i(t)\), the imported theorem emits
\[
S_i(t)=B_H(\widehat X_i(t),r_i(t)),
\qquad
r_i(t)=R_H\left[
\frac L\alpha\varepsilon+\frac2\alpha\sigma_i(t)+\frac1\alpha\tau_i(t)
\right].
\]
A point estimate \(\widehat X_i(t)\) is a unique finite localization output only when
\(\Delta_{\mathrm{loc},i}(t)>0\). If this gap fails, the input to the stitch theorem is an
ambiguity set or support ball, not an arbitrary token-ID tie-break. Distinct record sources at a
shared clock slice require the certified separation
\[
d_H(\widehat X_i,\widehat X_k)>r_i+r_k+m_{\mathrm{sep}}
\]
for a declared positive margin \(m_{\mathrm{sep}}\).

\subsection{Proto-record extraction and descent}

A raw local proto-record is a connected component of an invariant detector scalar above a frozen
threshold, together with its certified localization ball or support set, local clock interval,
sector payload, and uncertainty collar.
The detector scalar is built from record-algebra data, not from implementation object IDs. Extraction is
chart-natural when transporting the raw field by \(G_{ji}\) and \(F_{ji}\) carries extracted
components in chart \(i\) to the extracted components in chart \(j\), up to the declared support
collar.

Before temporal stitching, overlapping charts are descended. The overlap-descent relation joins
two proto-records only when their transported supports lie in the same connected component on a
real chart overlap and the same-component margin dominates the support collar. If distinct
component separation does not dominate the collar, the descent emits an ambiguity certificate
instead of choosing a label. On component-connected covers, this descent produces canonical
global descended records by the usual sheaf gluing argument: cocycle-compatible local components
with positive separation margins have a unique global component, and the construction is natural
under atlas restriction.

\subsection{Boundary-crossing germs and transported residuals}

A cross-boundary candidate edge is admissible only when the two descended records meet a real
patch interface in adjacent observer-clock intervals and share a common comparison chart. The
interface contact must be oriented and transverse: the signed-distance interval brackets the
interface, and the normal velocity has a positive lower bound. Tangencies, grazing contacts, and
file or shard-name boundaries without a physical interface are rejected. If the local slab contains
an interaction marker, the record is routed to the interaction solver instead of stitched by the
free-propagation rule.

For an admissible pair, the certificate transports the left record into the right comparison chart
and evaluates frozen residuals
\[
r_t,\quad r_{\partial},\quad r_x,\quad r_v,\quad r_{\mathrm{int}},
\]
for clock adjacency, boundary contact, \(H^3\) position, transported velocity or finite
difference, and interaction absence. The cost function is declared before seeing the proposed
matching. Record IDs, global IDs, shard-local IDs, and stitch keys are forbidden in admissibility,
costs, and tie-breaking.
All \(H^3\) position and velocity residuals are evaluated as certified interval or
uncertainty-ball distances. A point-estimate residual may be printed as a diagnostic, but it cannot
replace the localization radius in an admissibility or matching certificate.

\begin{lemma}[Candidate-edge naturality]\label{lem:h3-candidate-edge-naturality}
The admissible candidate-edge set and its residual intervals are invariant under common
\(\mathrm{SO}^{+}(1,3)\) chart changes and under gauge-bundle trivialization changes.
\end{lemma}

\begin{proof}
Every geometric term is computed from \(d_H\), oriented interface data, and the common clock line,
all of which are preserved by Theorem~\ref{thm:h3-geometric-naturality} and clock-map
compatibility. Sector and gauge residuals compare transported payloads after applying the declared
connector, so changing a local trivialization conjugates both sides of the comparison. The
residual intervals and the resulting admissibility predicate are therefore unchanged.
\end{proof}

\subsection{Assignment gap, nonbranching paths, and ambiguity}

Let \(L\) and \(R\) be the descended record sets on two adjacent time slabs. A stitch matching is a
partial one-to-one assignment between \(L\) and \(R\), augmented by declared appearance and
disappearance interval costs. Its certified gap is
\[
\Delta_{\mathrm{cert}}
=\min_{M\ne M_*}\underline C(M)-\overline C(M_*),
\]
where \(M_*\) is the proposed matching, \(\overline C(M_*)\) is its upper cost bound, and
\(\underline C(M)\) is the lower cost bound for every competing matching.
The localization gap \(\Delta_{\mathrm{loc}}\) and stitch-assignment gap
\(\Delta_{\mathrm{cert}}\) are distinct receipts: the first certifies a unique spatial support
inside one clock slice, while the second certifies a unique temporal assignment across slices.

\begin{theorem}[Robust ID-independent stitch assignment]\label{thm:h3-robust-stitch-assignment}
If the candidate-edge set is natural, the costs are frozen and ID-independent, the proposed
matching is one-to-one, and \(\Delta_{\mathrm{cert}}>0\), then \(M_*\) is the unique certified
minimum-cost stitch. The certified matching is natural under chart changes, gauge
trivializations, and relabeling of implementation record IDs.
\end{theorem}

\begin{proof}
The strict gap says every competing matching has lower cost bound above the proposed matching's
upper cost bound, so no competitor can tie or beat it under any realization inside the certified
intervals. Lemma~\ref{lem:h3-candidate-edge-naturality} makes the candidate graph and cost
intervals invariant under chart and gauge presentation changes. Because IDs are excluded from
admissibility and costs, relabeling them changes only bookkeeping fields, not the optimum.
\end{proof}

\begin{theorem}[Unavoidable ambiguity]\label{thm:h3-unavoidable-ambiguity}
If a nontrivial permutation of the visible descended records preserves the atlas, clocks,
interfaces, sector/gauge payloads, candidate graph, and cost intervals, then no natural
ID-independent finite procedure can certify a unique stitch across that slab.
\end{theorem}

\begin{proof}
A natural procedure must commute with the permutation. A unique certified output would therefore
be fixed by the permutation. But the permutation sends one equally admissible stitch to a distinct
equally admissible stitch. Choosing between them would introduce a label or presentation
dependence. The only natural finite output is an ambiguity label.
\end{proof}

The certified stitch graph is then the graph obtained by chaining the unique matchings across
adjacent slabs. Since each matching is partial one-to-one, every connected component is a
nonbranching path, ray, or isolated record. These are the \(H^3\) record-worldlines of this
section.

\subsection{Refinement naturality and optional Lorentzian lift}

For a refinement \(s\to r\), let \(Q_{sr}\) contract fine descended records to coarse records and
let \(\eta_{sr}\) bound the cost distortion between the fine and coarse stitch problems. If the
coarse certified gap obeys \(\Delta_r>2\eta_{sr}\), then the contracted fine optimum is
isomorphic to the coarse optimum. This is the refinement-naturality gate: a cross-boundary stitch
is not paper-grade until the coarse/fine contraction agrees within the declared margin.

The \(H^3\) certificate is spatial and observer-record-level. A Lorentzian worldline lift is a
separate optional criterion requiring lapse/shift data, a timelike speed margin, and compatibility
with the emergent four-dimensional chart. Without that lift, the certificate proves
record-token continuation in the \(H^3\) observer chart, not a full spacetime geodesic theorem.

\begin{theorem}[Certified H3 record-worldline stitching]\label{thm:h3-record-worldline-stitching}
Given a valid observer-facing \(H^3\) chart atlas, a common operational clock atlas, and a passing
cap-response localization certificate for every descended record token and clock slice,
chart-natural proto-record extraction, overlap descent with positive localization-ball separation
margins, real transverse interface-crossing germs, sector and gauge transport continuity, a frozen
ID-independent one-to-one assignment with \(\Delta_{\mathrm{cert}}>0\), no interaction marker in
the free-propagation slab, and a passing coarse/fine contraction gate, the emitted stitch graph is
a unique natural nonbranching \(H^3\) record-worldline graph. If any positive localization gap,
separation margin, or assignment gap is absent, the natural output is \(\mathrm{AMBIGUOUS}\); if an
interaction marker is present, the output is \(\mathrm{INTERACTION\ REQUIRED}\); if atlas,
interface, localization, or transport gates fail, the stitch is rejected.
\end{theorem}

\begin{proof}
The atlas, localization, and extraction gates make local proto-records presentation-independent.
The positive \(\Delta_{\mathrm{loc}}\) gate licenses point outputs; otherwise the procedure carries
support balls or ambiguity sets forward. Descent turns overlap-compatible local components into
descended records before any temporal decision is made.
The interface and transport gates define a natural candidate-edge graph by
Lemma~\ref{lem:h3-candidate-edge-naturality}. The strict assignment gap gives a unique
ID-independent matching by Theorem~\ref{thm:h3-robust-stitch-assignment}. Chaining partial
one-to-one matchings gives nonbranching components, and the refinement gate makes the result
stable under the declared coarse/fine contraction. The alternative statuses are exactly the
negations of these gates, with Theorem~\ref{thm:h3-unavoidable-ambiguity} forcing ambiguity when
the visible data do not separate a unique continuation.
\end{proof}

\subsection{Scope boundary of the measurement chapter}

This whole chapter has to preserve one final distinction. The fixed-cutoff measurement interface
is theorem-bearing on the microphysics surface: central record algebra, Born
probabilities, L\"uders conditioning, repeated-read stability, and the stated two-wing
Bell/CHSH theorem stack are written there as fixed-cutoff statements. What is \emph{not} proved
at the same level is the stronger global observer-continuation or strange-loop proposal sometimes
discussed in broader OPH notes. This paper should therefore use the fixed-cutoff measurement
package directly, while leaving the larger metaphysical closure proposal out of the main technical
chain.

% \section{Results at a Glance}
\section{Discussion: Matter, Antimatter, Supersymmetry, and Continuation Questions}

This section is intentionally not written in the same voice as the structural and quantitative
sections above. The OPH theorem surface distinguishes sharply between recovered-core
results, quantitative-closure branches, secondary quantitative branches, and continuations. Matter versus
antimatter, supersymmetry, and several adjacent questions lie on that boundary. They are too
important to omit from a particle-spectrum paper, but they are not part of the closed
prediction surface. The quark section above is likewise an obstruction and target-audit surface,
not a set of public numeric rows. Its \(S_3/D12\) numerical table is a target-anchored
diagnostic, not an exception to that rule.

\subsection{Matter and antimatter}
\label{subsec:baryogenesis-source}

The Standard Model structural branch fixes the chiral gauge architecture within which
matter and antimatter are defined. Once a chiral fermion representation is present, its conjugate
representation gives the corresponding antiparticle content. In that limited sense, OPH explains
why matter and antimatter both exist: they are paired by the realized gauge and chiral structure
of the low-energy branch.

The cosmological asymmetry requires a separate finite source object. At regulator \(r\), write
\[
\mathfrak B_r=(Q_r,C_r,E_r,\omega_{R,r},L_{r,T},p_{r,i},
\tau_r,T_r,\rho_{R,r}).
\]
Here \(Q_r\) is the physical quotient, \(C_r\) is its CP involution,
\(\omega_{R,r}\) is a CP-odd integral winding cocycle, \(L_{r,T}\) is the
physical quotient repair generator, \(p_{r,i}\) is the initial law,
\(\tau_r\) and \(T_r(\tau)\) are the physical clock and temperature map, and
\(\rho_{R,r}\) assigns record charges \(r_\psi\) to chiral matter.

For a record phase acting as
\(\psi\mapsto e^{ir_\psi\Theta_R}\psi\), the finite-quotient
anomaly-and-current theorem gives
\[
\boxed{
k_R=\sum_{\psi\,{\rm LH}}r_\psi\,2T_2(R_\psi)
}
\]
and
\[
\boxed{
\dot\Theta_R(T)=\theta_0
\sum_{q,q'\in Q_r}p_r(q,T)L_{r,T}(q,q')
\omega_{R,r}(q,q')
},
\qquad \theta_0=\frac{2\pi}{m_R}.
\]
The first equation is the mixed electroweak anomaly index. The second is an
oriented probability current in proper time. Quotient settlement supplies no
rate matrix, and a repair iteration count supplies no physical clock.

The theorem also fixes the sign boundary. If the initial law and generator are
CP symmetric,
\[
p_{r,i}(C_rq)=p_{r,i}(q),
\qquad
L_{r,T}(C_rq,C_rq')=L_{r,T}(q,q'),
\]
then every transition cancels its CP image and
\[
\boxed{\dot\Theta_R(T)=Y_B=0.}
\]
An oriented register and a normal-form map therefore select no cosmological
matter sign. A nonzero sign requires a quotient-intrinsic CP-odd boundary
condition, action term, or transition affinity.

The realized \(\mathbb Z_6\) determinant/deck winding supplies an integral
phase coordinate with inversion under orientation reversal. Its natural
Standard Model attachment is the central hypercharge direction. The associated
coefficient vanishes generation by generation:
\[
\boxed{
k_R^{YWW}=N_g(3Y_Q+Y_L)
=3\left(3\cdot\frac16-\frac12\right)=0.
}
\]
This is the mixed \(SU(2)_L^2U(1)_Y\) anomaly-cancellation condition. The
gauge/deck phase cannot directly supply the required
\(\Theta_R W\widetilde W\) coupling. Quotient periodicity and the twelve
electroweak zero modes do not select \(k_R=1\) or any other nonzero record
attachment.

A distinct global record attachment can carry a nonzero index. On the
three-generation branch,
\[
k_R(Y)=k_R(B-L)=0,
\quad k_R(B)=k_R(L)=3,
\quad k_R(B+L)=6.
\]
A quotient-visible gauge-singlet \(B+L\) record phase is therefore a viable
conditional attachment. Its derivation from OPH repair data is open. On
that branch, the freeze-out transport functional requires
\[
\left\langle\frac{\dot\Theta_R}{T}\right\rangle_{\rm fo}
=(4.463\pm0.028)\times10^{-9}.
\]
This is a source-generator target. It may not be consumed to select
\(L_{r,T}\), the attachment, or the CP-odd affinity.

The exact source theorem and gauge/deck no-go are closed at finite quotient.
The physical baryon abundance is work in progress. Its open source
packet consists of a distinct anomalous record phase, source-only transition
rates, a proper-time and temperature calibration, a CP-odd initial or boundary
law, sign-domain coherence, and the resulting full source history. Transport,
washout, and freeze-out must then be evaluated on that emitted history.

\subsection{Why the declared branch does not require supersymmetry}

The D10 appendix below contains the dedicated gauge-coupling quantitative-closure discussion
unification. Its main point is not that supersymmetry has been mathematically ruled out in every
possible extension of the program. Its point is narrower and more relevant to this paper: the declared
branch does not \emph{need} a supersymmetric partner spectrum in order to explain why
unification-like running might appear.

At one loop, the integrated D10 appendix records the familiar fact that Standard Model beta-function
coefficients do not produce successful naive unification, whereas MSSM-like coefficients do. The
The conditional calibration hypothesis is that edge-sector heat-kernel weights shift the effective
beta-function coefficients in an MSSM-like direction through geometric or entropic sector
multiplicity instead of through an actual low-energy superpartner spectrum. In that reading,
``unification-like behavior'' and ``a supersymmetric particle zoo'' come apart.

This calibration is not a derived unification theorem. The appendix records one
calibration branch in which Peter--Weyl multiplicities, the printed one-loop running frame, the
threshold conventions, and an additional fermionic-grading restriction together produce an
MSSM-like benchmark shift. It does not prove that those effective loop multiplicities,
statistics restrictions, and decoupling conventions are uniquely forced by OPH edge sectors alone.

For this paper, the correct conclusion is therefore scope-limited and modest:
\begin{enumerate}[leftmargin=2em]
\item the derivation reconstructs the realized Standard Model branch without having to
introduce a supersymmetric partner for every known field;
\item the integrated D10 appendix contains a conditional calibration in which unification-style running
features could arise from edge-sector structure instead of MSSM particle content;
\item the beta-shift package itself sits on explicit D10 calibration assumptions instead of a closed edge-sector theorem;
\item none of this is a universal theorem that supersymmetry is impossible.
\end{enumerate}

So if the paper asks ``why is there no supersymmetry in the derived spectrum?'', the best
technical answer is: because the realized branch stops at the Standard
Model content, and the existing unification discussion is trying to reproduce the relevant
running behavior through a separate D10 calibration package, without extending the particle content to a full
supersymmetric multiplet structure.

\subsection{Three Generations And Flavor Closure}

The structural Standard Model branch fixes \(N_g=3\) on the realized MAR-admissible
branch. This result is separate from the harder flavor
closure problem. Three generations are structurally recovered; the full flavor dictionary,
excitation map, and CKM closure are not emitted on this theorem surface, while the charged source landing from \(P\) to physical charged data is a corpus-limited no-go and the neutrino
lane has no source-closed flavor prediction. Its weighted-cycle candidate fails the NuFIT 6.1
correlated profile~\cite{nufit61}.
The flavor and family chapters therefore spend substantial space on
constructive object boundaries as well as final numbers.

\subsection{Ordinary matter and the hadron backend boundary}

A second distinction belongs in discussion form instead of in the results sections: the
difference between ``elementary rows exist'' and ``ordinary matter is fully derived.'' Ordinary
visible matter is dominated by hadrons, especially protons and neutrons. The OPH
derivation has a hadron pipeline, with source-only hadron masses gated by a production
backend export. The paper emits the eligible elementary rows and retains the weighted-cycle
coordinates only as a rejected comparison record, while the full
observed matter spectrum is not emitted from a source-only OPH backend.

\subsection{Claim boundary for conditional branches}

The reader should therefore interpret this whole discussion section with the same rule used by
\emph{Recovering Relativity and the Standard Model from Observer Overlap Consistency}~\cite{ophcompact}. Quark source-spread selection, charged-lepton source landing, fuller flavor-labeled neutrino, and source-only hadron mass theorems sit beyond the closed structural core; failure in those continuation branches would not undo the recovered-core gauge and gravity chain or erase the independently supported bosonic rows.

\section{Conclusion}

The particle-spectrum derivation is nontrivial and much more concrete than a qualitative aspiration.
The structural theorem surface fixes the realized Standard Model branch and its carrier roles.
The conditional quadratic-action theorem supplies classical Maxwell, perturbative pure-Yang--Mills,
and pure-Einstein modes, while their quantum-particle gates are open. The claim ledger
records those fail-closed carrier gates, the hierarchy/naturality bridge on the selected \(P/N\) branch, the Higgs/top row
on the declared D10/D11 surface, and the two-modulus quark source-spread obstruction on \(f_P\).
The quark audit supplies exact \(S_3\) heat-spectrum and \(L,Q\)-basis lemmas plus a
target-informed template ansatz whose mixed-chart residual is \(0.2946\%\). RSCC has an unproved
target-informed module/effect/sign/cumulant ledger; its lower-order ablation fits the same table
better, and F1--F6 are open. Neither diagnostic adds a dimensionful source mass, a common-scale
Yukawa packet, or a prediction row. The neutrino lane records an
isotropic no-go, a rejected target-informed weighted-cycle candidate, and open source-basis,
mass-label, and operator gates. The ledger also marks four separate row classes: the incomplete
\(P\)-closure declared-map audit, the electroweak \(W/Z\) chart audit, the charged-lepton
absolute-mass landing, and the hadron backend plus empirical closure split.
The \(W/Z\) row contains running/chart coordinates and a separate reference-fitted compare-only adapter.
Neither is a physical near-hit record. The
Ward-projected Thomson lane records the source/root audit, the mixed-provenance
root-plus-comparison-pixel diagnostic, and the measured endpoint as separate rows on the declared
\(P\)-closure surface. In the charged lane the theorem surface carries an exact
same-family witness, a determinant-line lift, and an algebraic mass readout from
theorem-grade \(A_{\mathrm{ch}}(P)\), while the source landing from the D10 quantitative
descendants of \(P\) to physical charged data is closed as a corpus-limited no-go. Source-only
hadron masses require a working OPH hadron backend. The empirical route requires a separate
\(e^+e^-\to\mathrm{hadrons}\) payload class; no integrated endpoint payload is emitted here.

That accounting belongs on the page because it is part of the derivation. The paper puts the
structural results, the incomplete declared-map fine-structure output, the emitted quantitative rows, and the
claim-boundary ledger onto one common surface. The bosonic \(P\)-driven trunk is explicit. The local
extension surface identifies the shared edge entropy with the D10 nonabelian sum on the lifted
product presentation of the realized quotient branch, and the same surface fixes that scalar to
\(P/4\). The local/global hierarchy bridge adds the mathematical coordinate
\(N_{\mathrm{CRC}}^{\mathrm{EW}}\), the zero bridge residual, the independent product-adjoint
count \(m_{\rm rep}=24\), and \(\epsilon_H=0\). The conditional screen has 12 ports and an
oriented 24-slot register, with no physical identification receipt between the equal counts. The particle spectrum is
organized by one local screen scale instead of by a list of unrelated particle inputs. The claim
ledger records scope boundaries for the weak validation pair, charged-lepton witnesses,
direct-top auxiliary codomain, strong-CP phase boundary,
neutrino comparison tension, source-only hadron backend rows, and empirical hadron closure rows.
Equality of the bridge coordinate with cosmic capacity is conditional on CP-1 to CP-3 and an
unconstructed physical map \(F\). The bridge value near \(3.532\times10^{122}\) and the
Planck-\(\Lambda\) central capacity near \(3.313\times10^{122}\) differ by about \(6.6\) percent.

The closure ledger groups its open quantitative rows under four generating objects and one
selector. This bookkeeping does not establish that the underlying prescription families are
exhaustive. G1 is the Ward-projected hadronic transport for CL-1 and CL-2. Its broad
internal-mass bracket \(S_{\mathrm{eff}}\in[0.5578,1.0543]\) contains the comparison value 0.8954,
but it spans about \(1.16\times10^8\) declared tolerance widths. The target is present only in the
compare-only payload, and the evaluation pixel differs from the comparison pixel. This bracket
carries no promotion. G2 requires a physical capacity readback map \(F\); the interval certificate
near \(3.53\times10^{122}\) applies only to the conditional electroweak bridge under CP-1 to CP-3.
G3 carries the certified forward electroweak chain. The 96-entry one-loop sweep excludes entries
within its frozen menu only. Higher-order, matching, threshold, tadpole, scale, and field-content
prescriptions are open. G4 supplies solver hygiene for the declared map. Twelve
preregistered flavor candidates were evaluated and excluded, which excludes those candidates
rather than the broader selector family.

The evidence ledger contains zero discriminating precommitted target-blind hits. Electroweak checks that
freeze analysis choices against known targets are audits, not target-blind predictions.
Other positive rows are theorem-level, conditional, target-exposed, or non-discriminating.

\appendix
\section{D10 Heat-Kernel Calibration Supplement}
\label{app:particle-d10-heat-kernel}

This appendix carries the D10 compact-group heat-kernel items needed by the particle surface.
The theorem-bearing leaf for the same-overlap thermal/Casimir law sits on the
screen-microphysics paper surface at fixed cutoff. What is recorded here is the
compact-group / Peter--Weyl lift and beta-shift bookkeeping used by the particle-side
calibration branch. This claim tier does not derive the physical one-loop beta coefficients
from edge sectors alone.

\subsection{\texorpdfstring{\(\mathcal R_U\)}{R\_U} Unified-Coupling Certificate}

The \(\mathcal R_U\) certificate is the D10 gate needed by the hierarchy discussion in the
synthesis paper. It certifies the unified diffusion coupling used in the electroweak
transmutation factor
\[
\frac{v}{E_\star}
=
P^{-1/2}\exp\!\left[-\frac{2\pi}{4\alpha_U(P)}\right].
\]
It does not certify the full clock branch
\[
\mathcal R_\gamma
=
\mathcal R_U+\mathcal R_\alpha+\mathcal R_e^{\mathrm{abs}}
+\mathcal R_{\mathrm{QCD/nuc}}^{133\mathrm{Cs}}
+\mathcal R_{\mathrm{atom}}^{133\mathrm{Cs}}.
\]
The full \(\mathcal R_\gamma\) stack requires source-only electromagnetic endpoint,
charged-lepton absolute-scale, cesium nuclear, and atomic spectral enclosures. The numerical package is
a numerical witness for \(\mathcal R_U\), with formal promotion gated by outward-rounded interval
arithmetic.

\paragraph{Frozen source packet.}
The public-pixel branch uses
\[
P_C
=
\decnum{1.6309682094039593248792798477826489413359828516279250606661507533907793398933432}.
\]
The source-audit branch uses the certified coordinate
\[
P_{\mathrm{fwd}}
=
1.630972095858897\ldots.
\]
It is interval-certified unique on its interval and, by the domain-global certificate, the only fixed point of its
readout map on the declared physical domain \(\alpha^{-1}\in[100,200]\)
(the interval contraction certificate of 2026-07-14,
\texttt{p\_global\_uniqueness\_certificate\_2026-07-14.json};
closure row CL-6, closed).
The finite heat-kernel cutoffs are
\[
N_2=128,\qquad N_3=64.
\]
The one-loop running packet is
\[
(b_1,b_2,b_3)=\left(\frac{33}{5},1,-3\right),
\qquad
N_c=3,
\qquad
\beta_{\mathrm{EW}}=N_c+1=4.
\]
This packet is frozen in the certificate. It is not silently derived in this appendix from the
recovered OPH core.

\begin{definition}[No-hidden-free-variable rule for \(\mathcal R_U\)]
The following objects must be fixed before comparison with electroweak or gravity data:
\[
(b_1,b_2,b_3),\quad
N_2,\quad
N_3,\quad
M_U(P),\quad
E_{\mathrm{cell}}(P),\quad
\beta_{\mathrm{EW}},\quad
\hbox{Casimir normalization},\quad
\hbox{matching scheme}.
\]
Changing any of them after comparison with \(M_W\), \(M_Z\), \(\alpha_s(M_Z)\), low-energy
electroweak couplings, or \(G_{\mathrm{exp}}\) demotes the row to calibration.
\end{definition}

\paragraph{Dimensionless source map.}
Set \(E_\star=1\) for the dimensionless calculation. For a trial coupling \(a\), define
\[
M_U(P)=e^{-2\pi}P^{1/6},
\qquad
E_{\mathrm{cell}}(P)=P^{-1/2},
\]
\[
v(P,a)=E_{\mathrm{cell}}(P)\exp\!\left[-\frac{2\pi}{4a}\right].
\]
For \(i=1,2,3\),
\[
\alpha_i^{-1}(\mu;P,a)
=
a^{-1}
+
\frac{b_i}{2\pi}\log\!\left(\frac{M_U(P)}{\mu}\right),
\qquad
\alpha_Y(\mu;P,a)=\frac35\alpha_1(\mu;P,a).
\]
The source \(Z\)-scale is the positive solution of
\[
\mu
=
\frac{v(P,a)}{2}
\sqrt{4\pi\alpha_2(\mu;P,a)+4\pi\alpha_Y(\mu;P,a)}.
\]
For \(\mathrm{SU}(2)\), with \(j=n/2\) and \(0\le n\le N_2\),
\[
d_j=2j+1,\qquad C_2(j)=j(j+1).
\]
For \(\mathrm{SU}(3)\), with \(0\le p,q\le N_3\),
\[
d_{p,q}=\frac{(p+1)(q+1)(p+q+2)}{2},
\qquad
C_2(p,q)=\frac{p^2+q^2+pq+3p+3q}{3}.
\]
The group sums are
\[
Z_G(t)=\sum_R d_R e^{-tC_2(R)},
\qquad
\bar\ell_G(t)=
\frac{1}{Z_G(t)}
\sum_R d_R e^{-tC_2(R)}\log d_R.
\]
Define
\[
t_2(P,a)=4\pi^2\alpha_2(\mu_Z(P,a);P,a),
\qquad
t_3(P,a)=4\pi^2\alpha_3(\mu_Z(P,a);P,a),
\]
and
\[
\Phi_U(P,a)
=
\bar\ell_{\mathrm{SU}(2)}(t_2(P,a))
+
\bar\ell_{\mathrm{SU}(3)}(t_3(P,a))
-
\frac{P}{4}.
\]
The certificate value is the zero \(\Phi_U(P,\alpha_U(P))=0\).

\paragraph{Numerical witness.}
For \(P=P_C\), the checker emits
\[
\begin{aligned}
\alpha_U(P_C)&=0.041124336195630495,\\
\alpha_U(P_C)^{-1}&=24.316501918546496.
\end{aligned}
\]
The corresponding dimensionless hierarchy data are
\[
\begin{aligned}
\frac{M_U(P_C)}{E_\star}&=0.0020260720012100037,\\
\frac{v(P_C)}{E_\star}&=2.0199803239725553\times10^{-17},\\
\frac{\mu_Z(P_C)}{E_\star}&=7.502403238986022\times10^{-18}.
\end{aligned}
\]
The emitted couplings at \(\mu_Z\) are
\[
\begin{aligned}
\alpha_1&=0.016885708038988166,\\
\alpha_Y&=0.010131424823392899,\\
\alpha_2&=0.033777889448908055,\\
\alpha_3&=0.11833602747445048.
\end{aligned}
\]
The heat-kernel parameters and entropy readouts are
\[
\begin{aligned}
t_2&=1.3334976254578108,\\
t_3&=4.6717191102770705,
\end{aligned}
\]
\[
\begin{aligned}
\bar\ell_{\mathrm{SU}(2)}&=0.39488144681077636,\\
\bar\ell_{\mathrm{SU}(3)}&=0.012860605540213493.
\end{aligned}
\]
Thus
\[
\begin{aligned}
\bar\ell_{\mathrm{SU}(2)}+\bar\ell_{\mathrm{SU}(3)}
&=0.40774205235098987,\\
P_C/4&=0.4077420523509898,
\end{aligned}
\]
and
\[
\begin{aligned}
\Phi_U(P_C,0.041124336195630495)
&=5.55\times10^{-17}.
\end{aligned}
\]
A centered derivative check gives
\[
\partial_a\Phi_U(P_C,a)\simeq -10.990524683118785.
\]
With \(\delta=10^{-6}\), the bracket signs are
\[
\Phi_U(P_C,\alpha_U-\delta)=1.0990748697592423\times10^{-5},
\]
\[
\Phi_U(P_C,\alpha_U+\delta)=-1.0990300680135956\times10^{-5}.
\]

For the certified source-audit pixel,
\[
\begin{aligned}
P_{\mathrm{fwd}}&=1.630972095858897\ldots,\\
\alpha_U(P_{\mathrm{fwd}})&=0.041124247441816685\ldots,\\
\frac{v(P_{\mathrm{fwd}})}{E_\star}&=2.0198114078576331\times10^{-17}.
\end{aligned}
\]
This branch is the one to cite when avoiding the measured Thomson endpoint in the upstream pixel
solve.

\begin{theorem}[Electroweak unified-coupling numerical witness]\label{thm:particle-ru-witness}
Assume the frozen source map, running packet, heat-kernel cutoff policy, and
no-hidden-free-variable rule above. Let
\[
I_U=[0.041123336195630494,\;0.041125336195630496].
\]
If outward-rounded interval arithmetic verifies
\[
0\in \Phi_U(P_C,I_U),
\qquad
0\notin \partial_a\Phi_U(P_C,I_U),
\]
then there is a unique \(\alpha_U(P_C)\in I_U\) satisfying
\[
\Phi_U(P_C,\alpha_U(P_C))=0.
\]
Moreover, the dependency graph of this zero contains no measured \(G\), no Planck unit formed
using measured \(G\), no measured \(\Lambda\), no measured \(M_Z\), no measured \(M_W\), and no
measured low-energy gauge coupling.
\end{theorem}

\begin{proof}
Continuity follows because the running couplings, self-consistent \(Z\)-scale solution, and
finite heat-kernel sums are continuous on the declared positive interval. The endpoint sign
change gives existence by the intermediate value theorem. The derivative exclusion gives
uniqueness. The dependency statement follows by inspection of the frozen source map: the residual
is built only from \(P_C\), the frozen running packet, the self-consistent \(Z\)-scale equation,
and finite compact-group representation sums.
\end{proof}

\paragraph{Precision policy.}
The exponential map makes downstream gravity displays sensitive to \(\alpha_U\):
\[
\frac{\partial\log G}{\partial \alpha_U}
\simeq
\frac{\pi}{\alpha_U^2}
\approx
1.86\times10^3.
\]
Therefore the number of digits printed for any downstream \(G\) or \(\varepsilon_{\mathrm{Cs}}\)
row may not exceed the interval precision certified by
\[
\mathcal R_U+\mathcal R_\alpha+\mathcal R_e^{\mathrm{abs}}
+\mathcal R_{\mathrm{QCD/nuc}}^{133\mathrm{Cs}}
+\mathcal R_{\mathrm{atom}}^{133\mathrm{Cs}}.
\]
The witness supports an electroweak hierarchy certificate. It does not support a 52-digit
source-only gravity prediction.

\subsection{One-Loop Calibration Frame}

At one loop, if couplings unify at \((M_U,\alpha_U)\),
\[
\alpha_i^{-1}(M_Z)=\alpha_U^{-1}+\frac{b_i}{2\pi}\ln\frac{M_U}{M_Z}.
\]
On the D10 branch these are the branch values \((M_U(P),\alpha_U(P))\) emitted by the forward
D10 solve; they are not inferred by inverse readback from measured low-energy couplings.
Writing \(A_i:=\alpha_i^{-1}(M_Z)\) and \(L:=\ln(M_U/M_Z)\), one finds
\[
L=\frac{2\pi}{b_1-b_2}(A_1-A_2),
\]
and the corresponding consistency relation for the D10 lane is
\[
A_3^{\mathrm{pred}}
=
\frac{b_3-b_2}{b_1-b_2}A_1
+
\frac{b_1-b_3}{b_1-b_2}A_2.
\]
This appendix keeps that algebra on the page as a calibration relation, not as a theorem that OPH
derives a supersymmetric UV spectrum.

\subsection{Benchmark coefficient comparison}

For the standard one-loop benchmark coefficients,
\[
\begin{array}{c|c|c}
\text{Model} & (b_1,b_2,b_3) & \alpha_s(M_Z)^{\mathrm{pred}} \\
\hline
\text{SM} & (41/10,-19/6,-7) & \approx 0.071 \\
\text{MSSM} & (33/5,1,-3) & \approx 0.116 \\
\text{Observed} & & 0.1179\pm 0.0010
\end{array}
\]
At one loop, the printed MSSM-style coefficients reproduce the observed closure far better than
the plain SM coefficients. On this D10 branch this is a benchmark matched by the D10
calibration branch, not a theorem that OPH derives a supersymmetric UV spectrum.

\subsection{Heat-Kernel Sector Law}

Read this subsection as the compact-group merge boundary above the fixed-cutoff microphysics
theorem. MaxEnt plus the bi-invariant compact-group package yield the familiar
\(d_R e^{-tC_2(R)}\) law once the Peter--Weyl lift is supplied; they are not a second,
independent proof of the finite-cutoff same-overlap Casimir branch.

\begin{theorem}[Heat-kernel edge-sector weights]
Under MaxEnt with bi-invariant constraints on a compact Lie group \(G\), the sector probabilities
take heat-kernel form:
\[
p_R(t)\propto d_R\,e^{-tC_2(R)},
\]
where \(d_R=\dim R\) and \(C_2(R)\) is the quadratic Casimir.
\end{theorem}

\begin{lemma}[Bi-invariant operators]
If \(G=\prod_i G_i\) is a compact semisimple Lie group written as a product of compact simple
factors, then any bi-invariant second-order differential operator on \(G\) has the form
\[
D=c_0\mathbf 1-\sum_i c_i\Delta_{G_i},
\]
where \(\Delta_{G_i}\) is the Laplace--Beltrami operator on the factor \(G_i\).
\end{lemma}

\begin{remark}
Bi-invariance is equivalent to \(D\in Z(U(\mathfrak g))\). Linear terms vanish because there are
no invariant vectors in the adjoint representation, and the quadratic part is proportional to the
Killing form on each simple factor, giving one independent coefficient per factor. The Casimir
element then acts as \(-\Delta_G\) in the regular representation.
\end{remark}

\subsection{Peter--Weyl Multiplicity and Beta Shifts}

By Peter--Weyl decomposition, the effective refinement-limit edge representation space is
\[
L^2(G)\cong \bigoplus_R V_R\otimes V_R^*.
\]
Entanglement traces over one factor give multiplicity \(d_R\) in \(p_R\), but loops see both
factors, so the effective multiplicity entering the calibration branch is
\[
N_{\mathrm{eff}}(R)=d_R\,p_R.
\]
The one-loop beta shift from edge sectors is then
\[
\Delta b_a=\sum_R p_R\,d_R\,T_a(R),
\]
where \(T_a(R)\) is the Dynkin index for gauge factor \(a\). This is the exact mathematical point
on the D10 lane where edge-sector heat-kernel weights feed the calibration branch. Matching to
MSSM-like one-loop running is a comparison benchmark, not a theorem-level
MSSM spectrum claim.

The branch boundary for theorem-level unification closure is explicit. One requires a derived
refinement-limit identification of which transportable sectors actually contribute to the running
carrier, a derived statistics/grading rule rather than an imposed fermionic-grading restriction,
controlled threshold and decoupling conventions on that same carrier, a derived representation
content/truncation rule for the sectors retained in the comparison, and a proof that the effective
loop multiplicities entering \(\Delta b_a\) are exactly the ones used in this calibration package
rather than nearby alternatives. The declared runtime surface does not emit the
full RG/matching/threshold/scheme packet: scheme lock, threshold map,
beta provenance, and interval composition are separate certificate requirements.

At the branch unification diffusion parameter \(t_U(P)\approx 1.64\) selected by the forward D10
solve, the benchmark shift is
\[
\Delta b\approx (2.49,4.38,3.97)
\qquad
\text{vs MSSM}
\qquad
(2.50,4.17,4.00).
\]
With the additional D10 calibration assumption of a fermionic-grading restriction to
half-integer \(\mathrm{SU}(2)\) sectors, this sharpens to
\[
\Delta b\approx (2.50,4.17,3.97),
\]
which matches the MSSM-style one-loop benchmark at the percent level under that declared
calibration package. Equivalently, the benchmark ratio
\[
\frac{\Delta b_3}{\Delta b_2}\approx 0.91
\]
sits close to the MSSM comparison value \(0.96\).

\section{H3 Worldline-Stitch Certificate}\label{app:h3-worldline-certificate}

This appendix records the finite evidence-facing certificate for
Theorem~\ref{thm:h3-record-worldline-stitching}. It is separate from the D10 heat-kernel
calibration package above. A passing certificate is a record-continuation statement in a declared
\(H^3\) observer chart, not a particle-species, mass, charge, or scattering-amplitude theorem.

The certificate payload has the following required blocks.
\begin{enumerate}
\item \textbf{Atlas.} It declares the hyperboloid model \(H^3_{R_H}\subset\mathbb R^{1,3}\),
the Lorentz signature \((-+++)\), the curvature radius \(R_H\), chart domains, and chart
transitions \(G_{ji}\in\mathrm{SO}^{+}(1,3)\). Transition residuals must certify Lorentz
inner-product preservation, orientation, and future-sheet preservation.
\item \textbf{Clock map.} Each local record interval is mapped to a common comparison-time line.
The observer-time adjacency margin must dominate the declared clock uncertainty.
\item \textbf{Chart-natural extraction.} Raw local records are connected components of a frozen
detector scalar on the \(H^3\) chart. The detector scalar, thresholds, support collars, and
chart-naturality residuals are included. Implementation record IDs are not inputs.
\item \textbf{Overlap descent.} Records are first descended across genuine chart overlaps. The
payload records same-component join margins, distinct-component separation margins, support
errors, and triple-overlap cocycle residuals.
\item \textbf{Interface crossing.} A candidate stitch must cross a real interface in adjacent
clock intervals. It includes oriented signed-distance margins and a positive lower bound on
normal velocity. Tangential, grazing, and file-boundary-only contacts are rejected.
\item \textbf{Transport.} The candidate edge is evaluated in a common chart after sector and
gauge transport. The atlas holonomy comparison is separate from the gauge-holonomy comparison.
\item \textbf{Assignment.} The matching is one-to-one with appearance/disappearance interval
costs. The proposed winner has an upper cost bound, every competitor has a lower cost bound, and
\[
\Delta_{\mathrm{cert}}
=\min_{M\ne M_*}\underline C(M)-\overline C(M_*)>0 .
\]
Record IDs, global IDs, stitch keys, and shard-local IDs are forbidden in admissibility, costs,
and tie-breaking.
\item \textbf{Refinement.} A coarse/fine pair supplies \(Q_{sr}\), a distortion bound
\(\eta_{sr}\), and a contracted-graph isomorphism check. The coarse gap must satisfy
\(\Delta_r>2\eta_{sr}\).
\item \textbf{Interaction firewall.} A free-propagation slab is required. If an interaction marker
is present, the terminal label is \(\mathrm{H3\_INTERACTION\_REQUIRED}\) and the event is handed
to the interaction solver.
\end{enumerate}

The allowed terminal statuses are
\[
\begin{gathered}
\mathrm{H3\_STITCH\_CERTIFIED},\quad
\mathrm{H3\_STITCH\_AMBIGUOUS},\quad
\mathrm{H3\_STITCH\_REJECTED},\\
\mathrm{H3\_INTERACTION\_REQUIRED},\quad
\mathrm{H3\_ATLAS\_INVALID},\quad
\mathrm{H3\_CERTIFICATE\_INCOMPLETE}.
\end{gathered}
\]
\(\mathrm{H3\_STITCH\_CERTIFIED}\) is emitted only when every block above passes and the robust
assignment gap is positive. Equal-cost matchings, preserved nontrivial permutations of the visible
data, or coarse/fine gaps below margin force \(\mathrm{H3\_STITCH\_AMBIGUOUS}\). Missing fields
force \(\mathrm{H3\_CERTIFICATE\_INCOMPLETE}\). Failed atlas, interface, metric, or transport gates
force rejection or atlas invalidity. In particular, Euclidean distances on
\(\texttt{h3SpatialPoint}\) coordinates are not certificate distances; the certificate distance is
the hyperboloid geodesic \(d_H\) from Section~\ref{sec:h3-record-worldlines}.


\begin{thebibliography}{9}

\bibitem{observableNormalForms}
B.\ M\"uller, D.\ Matscheko, and J.\ Hill,
\emph{Observation-Determined Normal Forms: Stability, Obstructions, and Refinement in Constraint
and Rewrite Systems},
2026.\\
Available at \url{https://github.com/FloatingPragma/observer-patch-holography/blob/main/extra/observable_normal_forms.pdf}.

\bibitem{ophsynthesis}
B.\ M\"uller, A.\ Osika, M.\ Poneder, K.\ Xue, B.\ Cassie, P.\ Nguyen, M.\ A.\ Visser, K.\ A.\ Anirudha, D.\ Matscheko, and J.\ Hill,
\emph{Observers Are All You Need}.
GitHub PDF:
\url{https://github.com/FloatingPragma/observer-patch-holography/blob/main/paper/observers_are_all_you_need.pdf}

\bibitem{ophconsensus}
B.\ M\"uller, K.\ Xue, K.\ A.\ Anirudha, D.\ Matscheko, and J.\ Hill,
\emph{Reality as a Consensus Protocol: The Fixed-Point Computation That Implements Physics}.
GitHub PDF:
\url{https://github.com/FloatingPragma/observer-patch-holography/blob/main/paper/reality_as_consensus_protocol.pdf}

\bibitem{ophmicrophysics}
B.\ M\"uller, A.\ Osika, K.\ Xue, B.\ Cassie, M.\ A.\ Visser, and D.\ Matscheko,
\emph{Federated Echosahedral Screen Microphysics: Patch Hardware, Records, and Observer Synchronization in OPH}.
GitHub PDF:
\url{https://github.com/FloatingPragma/observer-patch-holography/blob/main/paper/screen_microphysics_and_observer_synchronization.pdf}

\bibitem{ophcompact}
B.\ M\"uller, A.\ Osika, M.\ Poneder, K.\ Xue, P.\ Nguyen, M.\ A.\ Visser, and D.\ Matscheko,
\emph{Recovering Relativity and the Standard Model from Observer Overlap Consistency}.
GitHub PDF:
\url{https://github.com/FloatingPragma/observer-patch-holography/blob/main/paper/recovering_relativity_and_standard_model_structure_from_observer_overlap_consistency_compact.pdf}

\bibitem{codata2022}
National Institute of Standards and Technology,
\emph{2022 CODATA Recommended Values of the Fundamental Constants of Physics and Chemistry},
NIST SP 959, May 2024.
\url{https://physics.nist.gov/cuu/pdf/wallet_2022.pdf}

\bibitem{koide1983}
Y. Koide,
\emph{New view of quark and lepton mass hierarchy},
Physical Review D \textbf{28}, 252 (1983).

\bibitem{brannen2006}
C. A. Brannen,
\emph{The Lepton Masses},
preprint, 2006.
\url{http://brannenworks.com/MASSES2.pdf}

\bibitem{antusch2025running}
S. Antusch, K. Hinze, and S. Saad,
\emph{Updated Running Quark and Lepton Parameters at Various Scales},
arXiv:2510.01312, 2025.
\url{https://arxiv.org/abs/2510.01312}

\bibitem{planck18}
Planck Collaboration,
\emph{Planck 2018 results. VI. Cosmological parameters},
Astronomy \& Astrophysics 641, A6, 2020.
\url{https://arxiv.org/abs/1807.06209}

\bibitem{desidr2nu}
DESI Collaboration,
\emph{Constraints on Neutrino Physics from DESI DR2 BAO and DR1 Full Shape},
arXiv:2503.14744, 2025.
\url{https://arxiv.org/abs/2503.14744}

\bibitem{nufit61}
I. Esteban, M. C. Gonzalez-Garcia, M. Maltoni, I. Martinez-Soler, J. P. Pinheiro, and T. Schwetz,
\emph{NuFit-6.0: Updated global analysis of three-flavor neutrino oscillations},
JHEP \textbf{12} (2024) 216, arXiv:2410.05380,
with the NuFIT 6.1 (2025) profile-table release at
\url{https://www.nu-fit.org/?q=node/309}.

\bibitem{kniehl2013}
B. A. Kniehl,
\emph{All-order renormalization of the propagator matrix for fermionic systems
with flavor mixing},
Physical Review D \textbf{89}, 096005 (2014), arXiv:1308.3140.

\bibitem{gambino1999}
P. Gambino and P. A. Grassi,
\emph{The Nielsen identities of the Standard Model and the definition of mass},
Physical Review D \textbf{62}, 076002 (2000), arXiv:hep-ph/9907254.

\bibitem{buchholz1986}
D. Buchholz,
\emph{Gauss' law and the infraparticle problem},
Physics Letters B \textbf{174}, 331--334 (1986).

\bibitem{duch2023}
P. Duch and W. Dybalski,
\emph{Infrared problem in quantum electrodynamics},
arXiv:2307.06114.

\end{thebibliography}
\end{document}
